%Root main.tex
%---------------------------------------------------------------------------
%付録F　付録
%---------------------------------------------------------------------------

\chapter{VHDLプログラム}
\thispagestyle{fancy}
\label{chap:}
\section{MMCプログラム構成}
今回実装するMMCのベースとなるtopモジュールは，Headspring社から提供されている．そのため，三相MMCにて，新規で作成，変更したプログラムは主に，Control2.topモジュールとなる．また，FPGA内部モジュールの入出力に関して簡易図を示す．

 \begin{figure}[htbp]
  \begin{center}
  \includegraphics[width=150mm]
  {gaze/app_F/module.png}
  \caption{FPGA内部モジュールのブロック図}
  \label{vhdlmodule}
  \end{center}
 \end{figure}

topモジュールの直下にはAD変換を行う「AD7357」モジュール，MSSRDCDB制御及びSVPWM制御を行う「Control\_top2」モジュールがある．Control\_top2モジュールをSimulnkの拡張機能HDL\_Coderによって自動生成している．Control\_top2モジュールの直下にはデッドビート制御を行う「SRDCDB2」モジュール,空間ベクトル変調を行う「SVPWM\_calcurate」モジュール,バランス制御を行う「Cell\_balance」モジュールがある．

\begin{table}[H]
\begin{center}
 \caption{作成プログラムのバージョン履歴}
 \label{table:ver_log}
 \centering
 \scalebox{0.8}[1.0]{
  \begin{tabular}{c|c|c}
   \hline
   内容 & Simulinkファイル & ISEバージョン\\
   \hline \hline
   
   OPEN制御プログラム & SVPWM\_open\_VHDL\_17\_delay15us & ver44\\
   
   \hline
   
   MSSRDCDB制御を追加 & MMC\_9level\_half\_cel\_ver7 & ver68\\
   
   \hline
   
   オーバーサンプリング係数を100に変更 & MMC\_9level\_half\_cel\_ver7 & ver69\\
   
   \hline
   
   デッドビート制御の状態量ADC結果を校正 & MMC\_9level\_half\_cel\_ver7 & ver70\\
   
   \hline
   
   バランス制御の状態量ADC結果を校正 & MMC\_9level\_half\_cel\_ver7 & ver71\\
   
  \end{tabular}
  }
 \end{center}
\end{table}

上記ファイルは横山研究室のサーバー"Canan"内のディレクトリ\\
「¥¥192.168.253.51/Canaan/MMC」フォルダ内に保存されている．
以下に本論文の結果取得に使用したVHDLファイル（HSP0056-DA0-0001C\_20240120\_ver71）のソースコードを示す．

%Root main.tex
%---------------------------------------------------------------------------
%付録A　付録
%---------------------------------------------------------------------------

%\lstnewenvironment{mylisting}[1][]
%    {\lstset{
%        frame=single,
%        numbersep=10pt,
%        tabsize=2,
%        #1
%    }
%}{}
\begin{mylisting}[language=VHDL,caption=トップモジュール,label=appF_Top_module_VHDL]

----------------------------------------------------------------------------------
-- FPGA Logic for MMC Controller Board
--
--		Logic No.:		HSP0056-DA0-0001
--		Target Board:	HSP0056-B00-0001
--		Target FPGA:	XC6SLX75-2FGG676C
--
--		Version:		1.0 Build4
--		Date:			2016/12/15
--		Copyright:		Headspring Inc.
--
----------------------------------------------------------------------------------
library IEEE;
Library UNISIM;
use IEEE.STD_LOGIC_1164.ALL;
use IEEE.STD_LOGIC_ARITH.ALL;
use IEEE.STD_LOGIC_UNSIGNED.ALL;
use UNISIM.vcomponents.all;

-- Uncomment the following library declaration if using
-- arithmetic functions with Signed or Unsigned values
--use IEEE.NUMERIC_STD.ALL;


entity top is
	generic (
		FPGA_VERSION_MAJOR		: natural := 1;			-- 4bit
		FPGA_VERSION_MINOR		: natural := 0;			-- 4bit
		FPGA_VERSION_BUILD		: natural := 4			-- 8bit
	);

	Port ( 
		-- ====================
		-- CLK and Reset
		-- ====================
		CLK20_IN			: in std_logic;
		nRESET_IN			: in std_logic;

		BUS_CLK_IN			: in std_logic;

		-- ====================
		-- TI Bus
		-- ====================
		BUS_A_IN				: in  std_logic_vector (15 downto 0);
		BUS_BA1_IN				: in  std_logic;
		BUS_D_IO				: inout  std_logic_vector (15 downto 0);
		nBUS_CS2_IN				: in  std_logic;
		nBUS_CS3_IN				: in  std_logic;
		nBUS_RD_IN				: in  std_logic;
		nBUS_WR_IN				: in  std_logic;
		nBUS_WAIT_OUT			: out std_logic;
		BUS_RnW_IN				: in  std_logic;


		-- ====================
		-- DIO	(External I/O)
		-- ====================
		DI_IN					: in  std_logic_vector (15 downto 0);
		DO_OUT					: out std_logic_vector (15 downto 0);

		nLED_OUT				: out std_logic_vector (3 downto 0);
		DSW_IN					: in  std_logic_vector (3 downto 0);
		DSW_BOOT_IN				: in  std_logic_vector (1 downto 0);

		-- ====================
		-- DIO	(TI CPU)
		-- ====================
		DI_OUT					: out std_logic_vector (15 downto 0);
		DO_IN					: in  std_logic_vector (15 downto 0);

		-- ====================
		-- Gate Signal
		-- ====================
		INV_PWM_IN				: in  std_logic_vector (11 downto 0);
		nINV_PWM_OUT			: out  std_logic_vector (11 downto 0);
		nMMC_PWM1_OUT			: out  std_logic_vector (47 downto 0);
		nMMC_PWM2_OUT			: out  std_logic_vector (47 downto 0);

		-- ====================
		-- ADC (AD7357)
		-- ====================
		nADC_GR0_CS_OUT			: inout  std_logic;
		ADC_GR0_SCLK_OUT		: inout  std_logic;
		ADC_GR0_DATA_IN			: in  std_logic_vector (3 downto 0);

		nADC_GR1_CS_OUT			: inout  std_logic;
		ADC_GR1_SCLK_OUT		: inout  std_logic;
		ADC_GR1_DATA_IN			: in  std_logic_vector (3 downto 0);

		nADC_GR2_CS_OUT			: inout  std_logic;
		ADC_GR2_SCLK_OUT		: inout  std_logic;
		ADC_GR2_DATA_IN			: in  std_logic_vector (3 downto 0);

		nADC_GR3_CS_OUT			: inout  std_logic;
		ADC_GR3_SCLK_OUT		: inout  std_logic;
		ADC_GR3_DATA_IN			: in  std_logic_vector (3 downto 0);

		nADC_GR4_CS_OUT			: inout  std_logic;
		ADC_GR4_SCLK_OUT		: inout  std_logic;
		ADC_GR4_DATA_IN			: in  std_logic_vector (3 downto 0);

		nADC_GR5_CS_OUT			: inout  std_logic;
		ADC_GR5_SCLK_OUT		: inout  std_logic;
		ADC_GR5_DATA_IN			: in  std_logic_vector (3 downto 0);

		nADC_GR6_CS_OUT			: inout  std_logic;
		ADC_GR6_SCLK_OUT		: inout  std_logic;
		ADC_GR6_DATA_IN			: in  std_logic_vector (3 downto 0);

		nADC_GR7_CS_OUT			: inout  std_logic;
		ADC_GR7_SCLK_OUT		: inout  std_logic;
		ADC_GR7_DATA_IN			: in  std_logic_vector (3 downto 0);

		nADC_GR8_CS_OUT			: inout  std_logic;
		ADC_GR8_SCLK_OUT		: inout  std_logic;
		ADC_GR8_DATA_IN			: in  std_logic_vector (3 downto 0);

		nADC_GR9_CS_OUT			: inout  std_logic;
		ADC_GR9_SCLK_OUT		: inout  std_logic;
		ADC_GR9_DATA_IN			: in  std_logic_vector (3 downto 0);

		CPU_FPGA_RES			: out  std_logic_vector(9 downto 0);
		FPGA_RES				: out  std_logic_vector(9 downto 0)
		
	 
	);
	

end top;

architecture Behavioral of top is

	signal SPR			: std_logic_vector(15 downto 0)		:= "00000000" & "00000000";
	signal SPR_ST1		: std_logic_vector(15 downto 0)		:= "00000000" & "00000000";
	signal SPR_ST2		: std_logic_vector(15 downto 0)		:= "00000000" & "00000000";
	signal DUMMY		: std_logic_vector(4 downto 0)		:= "00000";

	signal RESET			: std_logic							:= '0';
	signal RESET_B			: std_logic							:= '0';
	signal RESET_BB			: std_logic							:= '0';

	signal CLK100			: std_logic							:= '0';
	signal CLK100FX			: std_logic							:= '0';
	signal CLK0				: std_logic							:= '0';
	signal CLKFB			: std_logic							:= '0';
	signal CLK20B			: std_logic							:= '0';
	signal LOCKED_100		: std_logic							:= '0';
	signal LOCK_STATUS_100	: std_logic_vector(7 downto 0)		:= "00000000";
	signal RESET_DCM_100	: std_logic							:= '0';
	signal DCM_OK_100		: std_logic							:= '0';
	signal LOCK_CNT_100		: std_logic_vector(18 downto 0)		:= "000" & "00000000" & "00000001";

	signal CLK150			: std_logic							:= '0';
	signal CLK150FX			: std_logic							:= '0';
	signal CLK0_2			: std_logic							:= '0';
	signal CLKFB_2			: std_logic							:= '0';
	signal LOCKED_150		: std_logic							:= '0';
	signal LOCK_STATUS_150	: std_logic_vector(7 downto 0)		:= "00000000";
	signal RESET_DCM_150	: std_logic							:= '0';
	signal DCM_OK_150		: std_logic							:= '0';
	signal LOCK_CNT_150		: std_logic_vector(18 downto 0)		:= "000" & "00000000" & "00000001";


	signal BUS_CS2_B		: std_logic							:= '0';
	signal BUS_CS2_BB		: std_logic							:= '0';
	signal BUS_CS3_B		: std_logic							:= '0';
	signal BUS_CS3_BB		: std_logic							:= '0';
	signal BUS_RD_B			: std_logic							:= '0';
	signal BUS_RD_BB		: std_logic							:= '0';
	signal BUS_WR_B			: std_logic							:= '0';
	signal BUS_WR_BB		: std_logic							:= '0';
	signal BUS_RnW_B		: std_logic							:= '0';
	signal BUS_RnW_BB		: std_logic							:= '0';
	signal BUS_A_B			: std_logic_vector(16 downto 0)		:= "000000000" & "00000000";
	signal BUS_A_BB			: std_logic_vector(16 downto 0)		:= "000000000" & "00000000";
	signal BUS_DW_B			: std_logic_vector(15 downto 0)		:= "00000000" & "00000000";
	signal BUS_DW_BB		: std_logic_vector(15 downto 0)		:= "00000000" & "00000000";
	signal BUS_CNT			: integer range 0 to 20				:= 0;

	signal BUS_RD_0			: std_logic							:= '0';
--	signal BUS_WR_0			: std_logic							:= '0';
	signal BUS_RD_1			: std_logic							:= '0';
	signal BUS_WR_1			: std_logic							:= '0';
	signal BUS_RD_2			: std_logic							:= '0';
	signal BUS_WR_2			: std_logic							:= '0';
	signal BUS_RD_3			: std_logic							:= '0';
	signal BUS_WR_3			: std_logic							:= '0';
	signal BUS_RD_FF		: std_logic							:= '0';
	signal BUS_WR_FF		: std_logic							:= '0';
	signal BUS_DR_0			: std_logic_vector(15 downto 0)		:= "00000000" & "00000000";
	signal BUS_DR_1			: std_logic_vector(15 downto 0)		:= "00000000" & "00000000";
	signal BUS_DR_1_2		: std_logic_vector(15 downto 0)		:= "00000000" & "00000000";
	signal BUS_DR_1_3		: std_logic_vector(15 downto 0)		:= "00000000" & "00000000";
	signal BUS_DR_2			: std_logic_vector(15 downto 0)		:= "00000000" & "00000000";
	signal BUS_DR_2_2		: std_logic_vector(15 downto 0)		:= "00000000" & "00000000";
	signal BUS_DR_3			: std_logic_vector(15 downto 0)		:= "00000000" & "00000000";
	signal BUS_DR_FF		: std_logic_vector(15 downto 0)		:= "00000000" & "00000000";

	signal DI_LPF			: std_logic_vector(15 downto 0)		:= "00000000" & "00000000";
	signal DSW_LPF			: std_logic_vector(3 downto 0)		:= "0000";
	signal DSW_BOOT_LPF		: std_logic_vector(1 downto 0)		:= "00";
	signal DO_LPF			: std_logic_vector(15 downto 0)		:= "00000000" & "00000000";
--	signal LED				: std_logic_vector(3 downto 0)		:= "0000";

	signal ADC_ERR				: std_logic							:= '0';
	signal WDT_ERR				: std_logic							:= '0';
	signal GB					: std_logic							:= '0';

	signal AD_CS				: std_logic							:= '0';
	signal AD_CLK				: std_logic							:= '0';

	signal ADC_GR0_CS_B				: std_logic							:= '0';
	signal ADC_GR1_CS_B				: std_logic							:= '0';
	signal ADC_GR2_CS_B				: std_logic							:= '0';
	signal ADC_GR3_CS_B				: std_logic							:= '0';
	signal ADC_GR4_CS_B				: std_logic							:= '0';
	signal ADC_GR5_CS_B				: std_logic							:= '0';
	signal ADC_GR6_CS_B				: std_logic							:= '0';
	signal ADC_GR7_CS_B				: std_logic							:= '0';
	signal ADC_GR8_CS_B				: std_logic							:= '0';
	signal ADC_GR9_CS_B				: std_logic							:= '0';

	signal ADC_GR0_SCLK_B			: std_logic							:= '0';
	signal ADC_GR1_SCLK_B			: std_logic							:= '0';
	signal ADC_GR2_SCLK_B			: std_logic							:= '0';
	signal ADC_GR3_SCLK_B			: std_logic							:= '0';
	signal ADC_GR4_SCLK_B			: std_logic							:= '0';
	signal ADC_GR5_SCLK_B			: std_logic							:= '0';
	signal ADC_GR6_SCLK_B			: std_logic							:= '0';
	signal ADC_GR7_SCLK_B			: std_logic							:= '0';
	signal ADC_GR8_SCLK_B			: std_logic							:= '0';
	signal ADC_GR9_SCLK_B			: std_logic							:= '0';

	signal ADC_GR0_DATA_B			: std_logic_vector(3 downto 0)		:= "0000";
	signal ADC_GR1_DATA_B			: std_logic_vector(3 downto 0)		:= "0000";
	signal ADC_GR2_DATA_B			: std_logic_vector(3 downto 0)		:= "0000";
	signal ADC_GR3_DATA_B			: std_logic_vector(3 downto 0)		:= "0000";
	signal ADC_GR4_DATA_B			: std_logic_vector(3 downto 0)		:= "0000";
	signal ADC_GR5_DATA_B			: std_logic_vector(3 downto 0)		:= "0000";
	signal ADC_GR6_DATA_B			: std_logic_vector(3 downto 0)		:= "0000";
	signal ADC_GR7_DATA_B			: std_logic_vector(3 downto 0)		:= "0000";
	signal ADC_GR8_DATA_B			: std_logic_vector(3 downto 0)		:= "0000";
	signal ADC_GR9_DATA_B			: std_logic_vector(3 downto 0)		:= "0000";

	signal AD_CH_EN					: std_logic_vector(9 downto 0)		:= "1111111111";

	signal nINT_SIG					: std_logic							:= '1';
	signal PWM_ON					: std_logic							:= '0';
	signal WDT_ON					: std_logic							:= '0';
	
	---------------
	signal nMMC_PWM1_buf: std_logic_vector(47 downto 0)		:= X"000000000000";
	signal nMMC_PWM2_buf: std_logic_vector(47 downto 0)		:= X"000000000000";
	------------------
	
	------20220303W
	signal nMMC_PWM1_FPGA_buf_buf_buf		: std_logic_vector(47 downto 0)	:= "10011111" & "11111111" & "11111111" & "11111111" & "11111111" & "00001111";
	signal nMMC_PWM2_FPGA_buf_buf_buf		: std_logic_vector(47 downto 0)	:= "10011111" & "11111111" & "11111111" & "11111111" & "11111111" & "00001111";
	
	-----------------------------

	attribute IOB: string;
	attribute IOB of BUS_A_B : signal is "FORCE";
	attribute IOB of BUS_CS2_B : signal is "FORCE";
	attribute IOB of BUS_RnW_B : signal is "FORCE";
	attribute IOB of BUS_DW_B : signal is "FORCE";
	attribute IOB of ADC_GR0_DATA_B : signal is "FORCE";
	attribute IOB of ADC_GR1_DATA_B : signal is "FORCE";
	attribute IOB of ADC_GR2_DATA_B : signal is "FORCE";
	attribute IOB of ADC_GR3_DATA_B : signal is "FORCE";
	attribute IOB of ADC_GR4_DATA_B : signal is "FORCE";
	attribute IOB of ADC_GR5_DATA_B : signal is "FORCE";
	attribute IOB of ADC_GR6_DATA_B : signal is "FORCE";
	attribute IOB of ADC_GR7_DATA_B : signal is "FORCE";
	attribute IOB of ADC_GR8_DATA_B : signal is "FORCE";
	attribute IOB of ADC_GR9_DATA_B : signal is "FORCE";
	attribute IOB of ADC_GR0_CS_B : signal is "FORCE";
	attribute IOB of ADC_GR1_CS_B : signal is "FORCE";
	attribute IOB of ADC_GR2_CS_B : signal is "FORCE";
	attribute IOB of ADC_GR3_CS_B : signal is "FORCE";
	attribute IOB of ADC_GR4_CS_B : signal is "FORCE";
	attribute IOB of ADC_GR5_CS_B : signal is "FORCE";
	attribute IOB of ADC_GR6_CS_B : signal is "FORCE";
	attribute IOB of ADC_GR7_CS_B : signal is "FORCE";
	attribute IOB of ADC_GR8_CS_B : signal is "FORCE";
	attribute IOB of ADC_GR9_CS_B : signal is "FORCE";
	attribute IOB of ADC_GR0_SCLK_B : signal is "FORCE";
	attribute IOB of ADC_GR1_SCLK_B : signal is "FORCE";
	attribute IOB of ADC_GR2_SCLK_B : signal is "FORCE";
	attribute IOB of ADC_GR3_SCLK_B : signal is "FORCE";
	attribute IOB of ADC_GR4_SCLK_B : signal is "FORCE";
	attribute IOB of ADC_GR5_SCLK_B : signal is "FORCE";
	attribute IOB of ADC_GR6_SCLK_B : signal is "FORCE";
	attribute IOB of ADC_GR7_SCLK_B : signal is "FORCE";
	attribute IOB of ADC_GR8_SCLK_B : signal is "FORCE";
	attribute IOB of ADC_GR9_SCLK_B : signal is "FORCE";

		--ver2--
	signal Vc_UP1_B		: std_logic_vector(13 downto 0);		--RfTdVc
	signal Vc_UP2_B		: std_logic_vector(13 downto 0);
	signal Vc_UP3_B		: std_logic_vector(13 downto 0);
	signal Vc_UP4_B		: std_logic_vector(13 downto 0);
	signal Vc_UN1_B		: std_logic_vector(13 downto 0);
	signal Vc_UN2_B		: std_logic_vector(13 downto 0);
	signal Vc_UN3_B		: std_logic_vector(13 downto 0);
	signal Vc_UN4_B		: std_logic_vector(13 downto 0);
	signal Vc_VP1_B		: std_logic_vector(13 downto 0);
	signal Vc_VP2_B		: std_logic_vector(13 downto 0);
	signal Vc_VP3_B		: std_logic_vector(13 downto 0);
	signal Vc_VP4_B		: std_logic_vector(13 downto 0);
	signal Vc_VN1_B		: std_logic_vector(13 downto 0);
	signal Vc_VN2_B		: std_logic_vector(13 downto 0);
	signal Vc_VN3_B		: std_logic_vector(13 downto 0);
	signal Vc_VN4_B		: std_logic_vector(13 downto 0);
	signal Vc_WP1_B		: std_logic_vector(13 downto 0);
	signal Vc_WP2_B		: std_logic_vector(13 downto 0);
	signal Vc_WP3_B		: std_logic_vector(13 downto 0);
	signal Vc_WP4_B		: std_logic_vector(13 downto 0);
	signal Vc_WN1_B		: std_logic_vector(13 downto 0);
	signal Vc_WN2_B		: std_logic_vector(13 downto 0);
	signal Vc_WN3_B		: std_logic_vector(13 downto 0);
	signal Vc_WN4_B		: std_logic_vector(13 downto 0);
	signal Iz_UP_B			: std_logic_vector(13 downto 0);		--A[dIz
	signal Iz_UN_B			: std_logic_vector(13 downto 0);
	signal Iz_VP_B			: std_logic_vector(13 downto 0);
	signal Iz_VN_B			: std_logic_vector(13 downto 0);
	signal Iz_WP_B			: std_logic_vector(13 downto 0);
	signal Iz_WN_B			: std_logic_vector(13 downto 0);
	signal Vdc_B			: std_logic_vector(13 downto 0);			--ddVdc
	signal Vuv_mmc_B		: std_logic_vector(13 downto 0);		--MMCodVo_mmc
	signal Vwu_mmc_B		: std_logic_vector(13 downto 0);
	signal Vvw_mmc_B		: std_logic_vector(13 downto 0);
	signal Vwu_inv_B		: std_logic_vector(13 downto 0);		--Co[^odVo_inv
	signal Vvw_inv_B		: std_logic_vector(13 downto 0);
	signal Vuv_inv_B		: std_logic_vector(13 downto 0);
	signal IU_inv_B		: std_logic_vector(13 downto 0);		--Co[^odIo_inv
	signal IV_inv_B		: std_logic_vector(13 downto 0);
	signal IW_inv_B		: std_logic_vector(13 downto 0);
	signal MRate_UP1_B	: std_logic_vector(15 downto 0);		--UT
	signal MRate_UP2_B	: std_logic_vector(15 downto 0);
	signal MRate_UP3_B	: std_logic_vector(15 downto 0);
	signal MRate_UP4_B	: std_logic_vector(15 downto 0);
	signal MRate_UN1_B	: std_logic_vector(15 downto 0);
	signal MRate_UN2_B	: std_logic_vector(15 downto 0);
	signal MRate_UN3_B	: std_logic_vector(15 downto 0);
	signal MRate_UN4_B	: std_logic_vector(15 downto 0);
	signal MRate_VP1_B	: std_logic_vector(15 downto 0);		--VT
	signal MRate_VP2_B	: std_logic_vector(15 downto 0);
	signal MRate_VP3_B	: std_logic_vector(15 downto 0);
	signal MRate_VP4_B	: std_logic_vector(15 downto 0);
	signal MRate_VN1_B	: std_logic_vector(15 downto 0);
	signal MRate_VN2_B	: std_logic_vector(15 downto 0);
	signal MRate_VN3_B	: std_logic_vector(15 downto 0);
	signal MRate_VN4_B	: std_logic_vector(15 downto 0);
	signal MRate_WP1_B	: std_logic_vector(15 downto 0);		--WT
	signal MRate_WP2_B	: std_logic_vector(15 downto 0);
	signal MRate_WP3_B	: std_logic_vector(15 downto 0);
	signal MRate_WP4_B	: std_logic_vector(15 downto 0);
	signal MRate_WN1_B	: std_logic_vector(15 downto 0);
	signal MRate_WN2_B	: std_logic_vector(15 downto 0);
	signal MRate_WN3_B	: std_logic_vector(15 downto 0);
	signal MRate_WN4_B	: std_logic_vector(15 downto 0);
		-------
--ver2.3---		
	signal nMMC_PWM1_FPGA_B	: std_logic_vector(47 downto 0)	:= "11111111" & "11111111" & "11111111" & "11111111" & "11111111" & "11111111";
	signal nMMC_PWM2_FPGA_B	: std_logic_vector(47 downto 0)	:= "11111111" & "11111111" & "11111111" & "11111111" & "11111111" & "11111111";
	signal nINT_SIG_FPGA_B		:	std_logic := '1';
	
	signal PWM_ON_FPGA_B		:	std_logic := '0';
	signal GB_FPGA_B			: std_logic;
	-----
	
	--ver2.6.4--
	signal REF_UP1_vector_B	: std_logic_vector(15 downto 0);
	signal REF_UP2_vector_B	: std_logic_vector(15 downto 0);
	signal REF_UP3_vector_B	: std_logic_vector(15 downto 0);
	signal REF_UP4_vector_B	: std_logic_vector(15 downto 0);
	signal REF_UN1_vector_B	: std_logic_vector(15 downto 0);
	signal REF_UN2_vector_B	: std_logic_vector(15 downto 0);
	signal REF_UN3_vector_B	: std_logic_vector(15 downto 0);
	signal REF_UN4_vector_B	: std_logic_vector(15 downto 0);
	signal REF_VP1_vector_B	: std_logic_vector(15 downto 0);
	signal REF_VP2_vector_B	: std_logic_vector(15 downto 0);
	signal REF_VP3_vector_B	: std_logic_vector(15 downto 0);
	signal REF_VP4_vector_B	: std_logic_vector(15 downto 0);
	signal REF_VN1_vector_B	: std_logic_vector(15 downto 0);
	signal REF_VN2_vector_B	: std_logic_vector(15 downto 0);
	signal REF_VN3_vector_B	: std_logic_vector(15 downto 0);
	signal REF_VN4_vector_B	: std_logic_vector(15 downto 0);
	---------------
	
	--ver3.4.2--chipscopefobOp---------------------
	signal vau_B        : std_logic_vector(15 downto 0);
	signal vav_B        : std_logic_vector(15 downto 0);
	signal vb_up1_B     : std_logic_vector(15 downto 0);
	signal vb_up2_B     : std_logic_vector(15 downto 0);
	signal vb_up3_B     : std_logic_vector(15 downto 0);
	signal vb_up4_B     : std_logic_vector(15 downto 0);
	signal vb_un1_B     : std_logic_vector(15 downto 0);
	signal vb_un2_B     : std_logic_vector(15 downto 0);
	signal vb_un3_B     : std_logic_vector(15 downto 0);
	signal vb_un4_B     : std_logic_vector(15 downto 0);
	signal vb_vp1_B     : std_logic_vector(15 downto 0);
	signal vb_vp2_B     : std_logic_vector(15 downto 0);
	signal vb_vp3_B     : std_logic_vector(15 downto 0);
	signal vb_vp4_B     : std_logic_vector(15 downto 0);
	signal vb_vn1_B     : std_logic_vector(15 downto 0);
	signal vb_vn2_B     : std_logic_vector(15 downto 0);
	signal vb_vn3_B     : std_logic_vector(15 downto 0);
	signal vb_vn4_B     : std_logic_vector(15 downto 0);
	-------------------------------------------------
	
	--ver3.7--lp
	signal GATE_EN_B    : std_logic;
	----------
	--ver3.7.1--
	signal VA_no_limit_U_B : std_logic_vector(15 downto 0);
	signal VA_no_limit_V_B : std_logic_vector(15 downto 0);
	---------------------------------
	
	--ver3.7.5--
	signal Voref_B	: std_logic_vector(15 downto 0);
	------------
	
	--ver3.8.5(debug)--
	signal Iz_UP_N_debug_out_B : std_logic_vector(15 downto 0);
	signal Iz_UP_N_out_B 	   : std_logic_vector(15 downto 0);
	-------------------
	
--QMDB_ver1.1--------------------------------------------------------------------------------------
	signal	dVoref_out_B        : std_logic_vector(15 downto 0);
	signal	IC_N_out_B	        : std_logic_vector(15 DOWNTO 0);
	signal	Vo_N_out_B	        : std_logic_vector(15 DOWNTO 0);
	signal	QMDB_OUTPUT_out_B   : std_logic_vector(15 DOWNTO 0);
	signal	Vdc_N_out_B		     : std_logic_vector(15 DOWNTO 0);
	signal   soft_cnt_out_B		  : std_logic_vector(15 DOWNTO 0);
	signal   clk10k_BB			  : std_logic;
	signal   cnt80k_debug_B		  : std_logic_vector(15 downto 0);
---------------------------------------------------------------------------------------------------

--ILA_signals----------------------------------------------------------------------------------------------
	signal	Vuv_ILA					: std_logic_vector(19 downto 0);
	signal	Vvw_ILA					: std_logic_vector(19 DOWNTO 0);
	signal	Ioutu_ILA				: std_logic_vector(19 DOWNTO 0);
	signal	Ioutv_ILA				: std_logic_vector(19 DOWNTO 0);
	signal	ILu_ILA					: std_logic_vector(19 DOWNTO 0);
	signal   ILv_ILA					: std_logic_vector(19 DOWNTO 0);
	signal   Vrefa_ILA				: std_logic_vector(19 DOWNTO 0);
	signal   Vrefb_ILA				: std_logic_vector(19 downto 0);
	signal	deltaTa_ILA				: std_logic_vector(15 downto 0);
	signal	deltaTb_ILA				: std_logic_vector(15 DOWNTO 0);
	signal	deltaTa_hold_ILA		: std_logic_vector(15 DOWNTO 0);
	signal	deltaTb_hold_ILA		: std_logic_vector(15 DOWNTO 0);
	signal	switche_contol_ILA	: std_logic_vector(13 DOWNTO 0);
---------------------------------------------------------------------------------------------------

	signal	CONTROL0	: std_logic_vector(35 DOWNTO 0);

--count_1MHz------------------------------------------------------
	signal	cnt_1MHz					: std_logic;
------------------------------------------------------------------
--count_50kHz------------------------------------------------------
	signal	cnt_50kHz					: std_logic;
------------------------------------------------------------------

---count_1MHz_module----------------------
	component count_1M is
		Port (
			clk	: in std_logic; --100MHz base clock
			cnt_1M	: out std_logic --1MHz count
		);
	end component;
------------------------------------------
---count_50kHz_module----------------------
	component count_50k is
		Port (
			clk	: in std_logic; --1MHz base clock
			cnt_50k	: out std_logic --50kHz count
		);
	end component;
------------------------------------------

---icon_module-----------------------------
	component icon_srdcdb is
		Port (
			CONTROL0	: inout std_logic_vector(35 DOWNTO 0)
		);
	end component;
-------------------------------------------

---ila_module------------------------------
	component ila_srdcdb is
		Port (
			CONTROL	: inout std_logic_vector(35 DOWNTO 0);
			CLK		: in std_logic;
			TRIG0		: in std_logic_vector(19 DOWNTO 0);
			TRIG1		: in std_logic_vector(19 DOWNTO 0);
			TRIG2		: in std_logic_vector(19 DOWNTO 0);
			TRIG3		: in std_logic_vector(19 DOWNTO 0);
			TRIG4		: in std_logic_vector(19 DOWNTO 0);
			TRIG5		: in std_logic_vector(19 DOWNTO 0);
			TRIG6		: in std_logic_vector(19 DOWNTO 0);
			TRIG7		: in std_logic_vector(19 DOWNTO 0);
			TRIG8		: in std_logic_vector(15 DOWNTO 0);
			TRIG9		: in std_logic_vector(15 DOWNTO 0);
			TRIG10	: in std_logic_vector(15 DOWNTO 0);
			TRIG11	: in std_logic_vector(15 DOWNTO 0);
			TRIG12	: in std_logic_vector(13 DOWNTO 0)
			);
	end component;
--------------------------------------------



	component BBB_LPF is
		generic (
			LPF_TIME	: natural;
			INVERT		: natural
		);
		Port (
			LPF_IN	: in std_logic;
			CLK		: in std_logic;
			LPF_OUT	: out std_logic
		);
	end component;


	component AD7357 is
		Port (
			CLK150		: in	std_logic;
			CLK100		: in	std_logic;
			RESET		: in	std_logic;
			A			: in	std_logic_vector(7 downto 0);
			DW			: in	std_logic_vector(15 downto 0);
			DR			: out	std_logic_vector(15 downto 0);
			DR2			: out	std_logic_vector(15 downto 0);
			RD			: in	std_logic;
			WR			: in	std_logic;
			
			AD_CS		: out	std_logic;
			AD_CLK		: out	std_logic;
			AD_CH_EN	: out	std_logic_vector(9 downto 0);
			
			ADC_GR0_CS		: in	std_logic;
			ADC_GR0_SCLK	: in	std_logic;
			ADC_GR0_DATA	: in	std_logic_vector(3 downto 0);

			ADC_GR1_CS		: in	std_logic;
			ADC_GR1_SCLK	: in	std_logic;
			ADC_GR1_DATA	: in	std_logic_vector(3 downto 0);

			ADC_GR2_CS		: in	std_logic;
			ADC_GR2_SCLK	: in	std_logic;
			ADC_GR2_DATA	: in	std_logic_vector(3 downto 0);

			ADC_GR3_CS		: in	std_logic;
			ADC_GR3_SCLK	: in	std_logic;
			ADC_GR3_DATA	: in	std_logic_vector(3 downto 0);

			ADC_GR4_CS		: in	std_logic;
			ADC_GR4_SCLK	: in	std_logic;
			ADC_GR4_DATA	: in	std_logic_vector(3 downto 0);

			ADC_GR5_CS		: in	std_logic;
			ADC_GR5_SCLK	: in	std_logic;
			ADC_GR5_DATA	: in	std_logic_vector(3 downto 0);

			ADC_GR6_CS		: in	std_logic;
			ADC_GR6_SCLK	: in	std_logic;
			ADC_GR6_DATA	: in	std_logic_vector(3 downto 0);

			ADC_GR7_CS		: in	std_logic;
			ADC_GR7_SCLK	: in	std_logic;
			ADC_GR7_DATA	: in	std_logic_vector(3 downto 0);

			ADC_GR8_CS		: in	std_logic;
			ADC_GR8_SCLK	: in	std_logic;
			ADC_GR8_DATA	: in	std_logic_vector(3 downto 0);

			ADC_GR9_CS		: in	std_logic;
			ADC_GR9_SCLK	: in	std_logic;
			ADC_GR9_DATA	: in	std_logic_vector(3 downto 0);
			
			ADC_ERR			: out	std_logic;
			
			--ver2--
			Vc_UP1		: out	std_logic_vector(13 downto 0);		--RfTdVc
			Vc_UP2		: out	std_logic_vector(13 downto 0);
			Vc_UP3		: out	std_logic_vector(13 downto 0);
			Vc_UP4		: out	std_logic_vector(13 downto 0);
			Vc_UN1		: out	std_logic_vector(13 downto 0);
			Vc_UN2		: out	std_logic_vector(13 downto 0);
			Vc_UN3		: out	std_logic_vector(13 downto 0);
			Vc_UN4		: out	std_logic_vector(13 downto 0);
			Vc_VP1		: out	std_logic_vector(13 downto 0);
			Vc_VP2		: out	std_logic_vector(13 downto 0);
			Vc_VP3		: out	std_logic_vector(13 downto 0);
			Vc_VP4		: out	std_logic_vector(13 downto 0);
			Vc_VN1		: out	std_logic_vector(13 downto 0);
			Vc_VN2		: out	std_logic_vector(13 downto 0);
			Vc_VN3		: out	std_logic_vector(13 downto 0);
			Vc_VN4		: out	std_logic_vector(13 downto 0);
			Vc_WP1		: out	std_logic_vector(13 downto 0);
			Vc_WP2		: out	std_logic_vector(13 downto 0);
			Vc_WP3		: out	std_logic_vector(13 downto 0);
			Vc_WP4		: out	std_logic_vector(13 downto 0);
			Vc_WN1		: out	std_logic_vector(13 downto 0);
			Vc_WN2		: out	std_logic_vector(13 downto 0);
			Vc_WN3		: out	std_logic_vector(13 downto 0);
			Vc_WN4		: out	std_logic_vector(13 downto 0);
			Iz_UP			: out	std_logic_vector(13 downto 0);		--A[dIz
			Iz_UN			: out	std_logic_vector(13 downto 0);
			Iz_VP			: out	std_logic_vector(13 downto 0);
			Iz_VN			: out	std_logic_vector(13 downto 0);
			Iz_WP			: out	std_logic_vector(13 downto 0);
			Iz_WN			: out	std_logic_vector(13 downto 0);
			Vdc			: out	std_logic_vector(13 downto 0);		--ddVdc
			Vuv_mmc		: out	std_logic_vector(13 downto 0);		--MMCodVo_mmc
			Vwu_mmc		: out	std_logic_vector(13 downto 0);
			Vvw_mmc		: out	std_logic_vector(13 downto 0);
			Vwu_inv		: out	std_logic_vector(13 downto 0);		--Co[^odVo_inv
			Vvw_inv		: out	std_logic_vector(13 downto 0);
			Vuv_inv		: out	std_logic_vector(13 downto 0);
			IU_inv		: out	std_logic_vector(13 downto 0);		--Co[^odIo_inv
			IV_inv		: out	std_logic_vector(13 downto 0);
			IW_inv		: out	std_logic_vector(13 downto 0)		
			-------			
			);
	end component;


	component MMC_PWM is
		Port (
			CLK100		: in	std_logic;
			RESET		: in	std_logic;
			A			: in	std_logic_vector(7 downto 0);
			DW			: in	std_logic_vector(15 downto 0);
			DR			: out	std_logic_vector(15 downto 0);
			DR2			: out	std_logic_vector(15 downto 0);
			DR3			: out	std_logic_vector(15 downto 0);
			RD			: in	std_logic;
			WR			: in	std_logic;
			
			nMMC_PWM1	: out	std_logic_vector(47 downto 0);
			nMMC_PWM2	: out	std_logic_vector(47 downto 0);
			nINT_SIG	: out	std_logic;
			
			PWM_ON		: out	std_logic;
			GB			: in	std_logic;
			
			--20190814_PFPGA_PIfobOp--
			MRate_UP1	: in	std_logic_vector(15 downto 0);		--UT
			MRate_UP2	: in	std_logic_vector(15 downto 0);
			MRate_UP3	: in	std_logic_vector(15 downto 0);
			MRate_UP4	: in	std_logic_vector(15 downto 0);
			MRate_UN1	: in	std_logic_vector(15 downto 0);
			MRate_UN2	: in	std_logic_vector(15 downto 0);
			MRate_UN3	: in	std_logic_vector(15 downto 0);
			MRate_UN4	: in	std_logic_vector(15 downto 0);
			MRate_VP1	: in	std_logic_vector(15 downto 0);		--VT
			MRate_VP2	: in	std_logic_vector(15 downto 0);
			MRate_VP3	: in	std_logic_vector(15 downto 0);
			MRate_VP4	: in	std_logic_vector(15 downto 0);
			MRate_VN1	: in	std_logic_vector(15 downto 0);
			MRate_VN2	: in	std_logic_vector(15 downto 0);
			MRate_VN3	: in	std_logic_vector(15 downto 0);
			MRate_VN4	: in	std_logic_vector(15 downto 0);
			MRate_WP1	: in	std_logic_vector(15 downto 0);		--WT
			MRate_WP2	: in	std_logic_vector(15 downto 0);
			MRate_WP3	: in	std_logic_vector(15 downto 0);
			MRate_WP4	: in	std_logic_vector(15 downto 0);
			MRate_WN1	: in	std_logic_vector(15 downto 0);
			MRate_WN2	: in	std_logic_vector(15 downto 0);
			MRate_WN3	: in	std_logic_vector(15 downto 0);
			MRate_WN4	: in	std_logic_vector(15 downto 0);
			-----------------------------------
			
			--ver2.6.4--
			REF_UP1_vector	: out std_logic_vector(15 downto 0);
			REF_UP2_vector	: out std_logic_vector(15 downto 0);
			REF_UP3_vector	: out std_logic_vector(15 downto 0);
			REF_UP4_vector	: out std_logic_vector(15 downto 0);
			REF_UN1_vector	: out std_logic_vector(15 downto 0);
			REF_UN2_vector	: out std_logic_vector(15 downto 0);
			REF_UN3_vector	: out std_logic_vector(15 downto 0);
			REF_UN4_vector	: out std_logic_vector(15 downto 0);
			REF_VP1_vector	: out std_logic_vector(15 downto 0);
			REF_VP2_vector	: out std_logic_vector(15 downto 0);
			REF_VP3_vector	: out std_logic_vector(15 downto 0);
			REF_VP4_vector	: out std_logic_vector(15 downto 0);
			REF_VN1_vector	: out std_logic_vector(15 downto 0);
			REF_VN2_vector	: out std_logic_vector(15 downto 0);
			REF_VN3_vector	: out std_logic_vector(15 downto 0);
			REF_VN4_vector	: out std_logic_vector(15 downto 0)
			----------------
		
			);
	end component;


	component MMC_WDT is
		Port (
			CLK100		: in	std_logic;
			RESET		: in	std_logic;
			A			: in	std_logic_vector(7 downto 0);
			DW			: in	std_logic_vector(15 downto 0);
			DR			: out	std_logic_vector(15 downto 0);
			RD			: in	std_logic;
			WR			: in	std_logic;
			
			WDT_ON		: out	std_logic;
			ERROR		: out	std_logic
			);
	end component;


--	--ver2--
--	component control_top is
--		Port (
--			CLK100		: in	std_logic;									--100MNbN
--			Vc_UP1		: in	std_logic_vector(13 downto 0);		--RfTdVc
--			Vc_UP2		: in	std_logic_vector(13 downto 0);
--			Vc_UP3		: in	std_logic_vector(13 downto 0);
--			Vc_UP4		: in	std_logic_vector(13 downto 0);
--			Vc_UN1		: in	std_logic_vector(13 downto 0);
--			Vc_UN2		: in	std_logic_vector(13 downto 0);
--			Vc_UN3		: in	std_logic_vector(13 downto 0);
--			Vc_UN4		: in	std_logic_vector(13 downto 0);
--			Vc_VP1		: in	std_logic_vector(13 downto 0);
--			Vc_VP2		: in	std_logic_vector(13 downto 0);
--			Vc_VP3		: in	std_logic_vector(13 downto 0);
--			Vc_VP4		: in	std_logic_vector(13 downto 0);
--			Vc_VN1		: in	std_logic_vector(13 downto 0);
--			Vc_VN2		: in	std_logic_vector(13 downto 0);
--			Vc_VN3		: in	std_logic_vector(13 downto 0);
--			Vc_VN4		: in	std_logic_vector(13 downto 0);
--			Vc_WP1		: in	std_logic_vector(13 downto 0);
--			Vc_WP2		: in	std_logic_vector(13 downto 0);
--			Vc_WP3		: in	std_logic_vector(13 downto 0);
--			Vc_WP4		: in	std_logic_vector(13 downto 0);
--			Vc_WN1		: in	std_logic_vector(13 downto 0);
--			Vc_WN2		: in	std_logic_vector(13 downto 0);
--			Vc_WN3		: in	std_logic_vector(13 downto 0);
--			Vc_WN4		: in	std_logic_vector(13 downto 0);
--			Iz_UP			: in	std_logic_vector(13 downto 0);		--A[dIz
--			Iz_UN			: in	std_logic_vector(13 downto 0);
--			Iz_VP			: in	std_logic_vector(13 downto 0);
--			Iz_VN			: in	std_logic_vector(13 downto 0);
--			Iz_WP			: in	std_logic_vector(13 downto 0);
--			Iz_WN			: in	std_logic_vector(13 downto 0);
--			Vdc			: in	std_logic_vector(13 downto 0);			--ddVdc
--			Vuv_mmc		: in	std_logic_vector(13 downto 0);		--MMCodVo_mmc
--			Vwu_mmc		: in	std_logic_vector(13 downto 0);
--			Vvw_mmc		: in	std_logic_vector(13 downto 0);
--			Vwu_inv		: in	std_logic_vector(13 downto 0);		--Co[^odVo_inv
--			Vvw_inv		: in	std_logic_vector(13 downto 0);
--			Vuv_inv		: in	std_logic_vector(13 downto 0);
--			IU_inv		: in	std_logic_vector(13 downto 0);		--Co[^odIo_inv
--			IV_inv		: in	std_logic_vector(13 downto 0);
--			IW_inv		: in	std_logic_vector(13 downto 0);
--			MRate_UP1	: out	std_logic_vector(15 downto 0);		--UT
--			MRate_UP2	: out	std_logic_vector(15 downto 0);
--			MRate_UP3	: out	std_logic_vector(15 downto 0);
--			MRate_UP4	: out	std_logic_vector(15 downto 0);
--			MRate_UN1	: out	std_logic_vector(15 downto 0);
--			MRate_UN2	: out	std_logic_vector(15 downto 0);
--			MRate_UN3	: out	std_logic_vector(15 downto 0);
--			MRate_UN4	: out	std_logic_vector(15 downto 0);
--			MRate_VP1	: out	std_logic_vector(15 downto 0);		--VT
--			MRate_VP2	: out	std_logic_vector(15 downto 0);
--			MRate_VP3	: out	std_logic_vector(15 downto 0);
--			MRate_VP4	: out	std_logic_vector(15 downto 0);
--			MRate_VN1	: out	std_logic_vector(15 downto 0);
--			MRate_VN2	: out	std_logic_vector(15 downto 0);
--			MRate_VN3	: out	std_logic_vector(15 downto 0);
--			MRate_VN4	: out	std_logic_vector(15 downto 0);
--			MRate_WP1	: out	std_logic_vector(15 downto 0);		--WT
--			MRate_WP2	: out	std_logic_vector(15 downto 0);
--			MRate_WP3	: out	std_logic_vector(15 downto 0);
--			MRate_WP4	: out	std_logic_vector(15 downto 0);
--			MRate_WN1	: out	std_logic_vector(15 downto 0);
--			MRate_WN2	: out	std_logic_vector(15 downto 0);
--			MRate_WN3	: out	std_logic_vector(15 downto 0);
--			MRate_WN4	: out	std_logic_vector(15 downto 0);
--			--ver3.4.2--chipscopefobOp---------------------
--			VAu_out 		: out std_logic_vector(15 downto 0);
--			VAv_out 		: out std_logic_vector(15 downto 0);
--			vb_up1_out  : out std_logic_vector(15 downto 0);
--			vb_up2_out  : out std_logic_vector(15 downto 0);
--			vb_up3_out  : out std_logic_vector(15 downto 0);
--			vb_up4_out  : out std_logic_vector(15 downto 0);
--			vb_un1_out  : out std_logic_vector(15 downto 0);
--			vb_un2_out  : out std_logic_vector(15 downto 0);
--			vb_un3_out  : out std_logic_vector(15 downto 0);
--			vb_un4_out  : out std_logic_vector(15 downto 0);
--			vb_vp1_out  : out std_logic_vector(15 downto 0);
--			vb_vp2_out  : out std_logic_vector(15 downto 0);
--			vb_vp3_out  : out std_logic_vector(15 downto 0);
--			vb_vp4_out  : out std_logic_vector(15 downto 0);
--			vb_vn1_out  : out std_logic_vector(15 downto 0);
--			vb_vn2_out  : out std_logic_vector(15 downto 0);
--			vb_vn3_out  : out std_logic_vector(15 downto 0);
--			vb_vn4_out  : out std_logic_vector(15 downto 0);
--			-------------------------------------------------
--			
--			--ver3.7--lp
--			GATE_EN : in std_logic;
--			----------
--			--ver3.7.1--
--			VA_no_limit_U : out std_logic_vector(15 downto 0);
--			VA_no_limit_V : out std_logic_vector(15 downto 0);
--			----------
--			--ver3.7.5--
--			Voref_out	: out std_logic_vector(15 downto 0);
--			------------
--			--ver3.8.5(debug)--
--			Iz_UP_N_debug_out : out std_logic_vector(15 downto 0);
--			Iz_UP_N_out 		: out std_logic_vector(15 downto 0);
--			-------------------
----QMDB_ver1.1--------------------------------------------------------------------------------------
--		   dVoref_out       : OUT     std_logic_vector(15 downto 0);
--		   IC_N_out	        : OUT   	std_logic_vector(15 DOWNTO 0);
--	 	   Vo_N_out	        : OUT 	   std_logic_vector(15 DOWNTO 0);
--	 	   QMDB_OUTPUT_out  : OUT     std_logic_vector(15 DOWNTO 0);
--	 	   Vdc_N_out		  : OUT     std_logic_vector(15 DOWNTO 0);
--		  soft_cnt_out		: OUT		 std_logic_vector(15 DOWNTO 0);
--		  clk10k_B_out			: OUT 	 std_logic;
--		  	cnt80k_debug		: OUT		 std_logic_vector(15 DOWNTO 0)
-----------------------------------------------------------------------------------------------------
--			);
--	end component;
--	-----------
	
		--ver2.3--
	component MMC_PWM_FPGA is
		Port (
			CLK100		: in	std_logic;
			RESET		: in	std_logic;
			A			: in	std_logic_vector(7 downto 0);
			DW			: in	std_logic_vector(15 downto 0);
	--		DR			: out	std_logic_vector(15 downto 0)	:= "00000000" & "00000000";
	--		DR2		: out	std_logic_vector(15 downto 0)	:= "00000000" & "00000000";
	--		DR3		: out	std_logic_vector(15 downto 0)	:= "00000000" & "00000000";
			RD			: in	std_logic;
			WR			: in	std_logic;
			
			GB			: in  std_logic;
			
			MRate_UP1 : in std_logic_vector(15 downto 0);
			MRate_UP2 : in std_logic_vector(15 downto 0);
			MRate_UP3 : in std_logic_vector(15 downto 0);
			MRate_UP4 : in std_logic_vector(15 downto 0);
			MRate_UN1 : in std_logic_vector(15 downto 0);
			MRate_UN2 : in std_logic_vector(15 downto 0);
			MRate_UN3 : in std_logic_vector(15 downto 0);
			MRate_UN4 : in std_logic_vector(15 downto 0);
			MRate_VP1 : in std_logic_vector(15 downto 0);
			MRate_VP2 : in std_logic_vector(15 downto 0);
			MRate_VP3 : in std_logic_vector(15 downto 0);
			MRate_VP4 : in std_logic_vector(15 downto 0);
			MRate_VN1 : in std_logic_vector(15 downto 0);
			MRate_VN2 : in std_logic_vector(15 downto 0);
			MRate_VN3 : in std_logic_vector(15 downto 0);
			MRate_VN4 : in std_logic_vector(15 downto 0);
			
			--ver2.6.4--
			REF_UP1_vector	: in std_logic_vector(15 downto 0);
			REF_UP2_vector	: in std_logic_vector(15 downto 0);
			REF_UP3_vector	: in std_logic_vector(15 downto 0);
			REF_UP4_vector	: in std_logic_vector(15 downto 0);
			REF_UN1_vector	: in std_logic_vector(15 downto 0);
			REF_UN2_vector	: in std_logic_vector(15 downto 0);
			REF_UN3_vector	: in std_logic_vector(15 downto 0);
			REF_UN4_vector	: in std_logic_vector(15 downto 0);
			REF_VP1_vector	: in std_logic_vector(15 downto 0);
			REF_VP2_vector	: in std_logic_vector(15 downto 0);
			REF_VP3_vector	: in std_logic_vector(15 downto 0);
			REF_VP4_vector	: in std_logic_vector(15 downto 0);
			REF_VN1_vector	: in std_logic_vector(15 downto 0);
			REF_VN2_vector	: in std_logic_vector(15 downto 0);
			REF_VN3_vector	: in std_logic_vector(15 downto 0);
			REF_VN4_vector	: in std_logic_vector(15 downto 0);
			----------------
			
			--ver3.4.2--chipscopefobOp---------------------
			VAu 	  : in std_logic_vector(15 downto 0);
			VAv 	  : in std_logic_vector(15 downto 0);
			vb_up1  : in std_logic_vector(15 downto 0);
			vb_up2  : in std_logic_vector(15 downto 0);
			vb_up3  : in std_logic_vector(15 downto 0);
			vb_up4  : in std_logic_vector(15 downto 0);
			vb_un1  : in std_logic_vector(15 downto 0);
			vb_un2  : in std_logic_vector(15 downto 0);
			vb_un3  : in std_logic_vector(15 downto 0);
			vb_un4  : in std_logic_vector(15 downto 0);
			vb_vp1  : in std_logic_vector(15 downto 0);
			vb_vp2  : in std_logic_vector(15 downto 0);
			vb_vp3  : in std_logic_vector(15 downto 0);
			vb_vp4  : in std_logic_vector(15 downto 0);
			vb_vn1  : in std_logic_vector(15 downto 0);
			vb_vn2  : in std_logic_vector(15 downto 0);
			vb_vn3  : in std_logic_vector(15 downto 0);
			vb_vn4  : in std_logic_vector(15 downto 0);
			-------------------------------------------------
			
			nMMC_PWM1_FPGA	: out	std_logic_vector(47 downto 0)	:= "11111111" & "11111111" & "11111111" & "11111111" & "11111111" & "11111111";
			nMMC_PWM2_FPGA	: out	std_logic_vector(47 downto 0)	:= "11111111" & "11111111" & "11111111" & "11111111" & "11111111" & "11111111";
			nINT_SIG_FPGA	: out	std_logic := '1';
			
			PWM_ON_FPGA		: out	std_logic := '0';
			GB_FPGA			: in	std_logic;
			
			--ver3.7--lp
			GATE_EN_out : out std_logic;
			----------
			--ver3.7.1--
			VA_no_limit_U : in std_logic_vector(15 downto 0);
			VA_no_limit_V : in std_logic_vector(15 downto 0);
			----------
			--ver3.7.5--
			Voref_BB	: in std_logic_vector(15 downto 0);
			------------
			--ver3.8.5(debug)--
			Iz_UP_N_debug_out_BB : in std_logic_vector(15 downto 0);
			Iz_UP_N_out_BB 		: in std_logic_vector(15 downto 0);
			-------------------
--QMDB_ver1.1--------------------------------------------------------------------------------------
		   dVoref_out_BB        : IN     std_logic_vector(15 downto 0);
		   IC_N_out_BB	        : IN   	std_logic_vector(15 DOWNTO 0);
	 	   Vo_N_out_BB	        : IN 	   std_logic_vector(15 DOWNTO 0);
	 	   QMDB_OUTPUT_out_BB   : IN     std_logic_vector(15 DOWNTO 0);
	 	   Vdc_N_out_BB		    : IN     std_logic_vector(15 DOWNTO 0);
			soft_cnt_out_BB	      : IN     std_logic_vector(15 DOWNTO 0);
			clk10k_BBB           : IN     std_logic;
			cnt80k_debug_BB			: IN 		std_logic_vector(15 DOWNTO 0)
---------------------------------------------------------------------------------------------------

			);
	end component;
		-----------
		
		
			--20220303
		component Control_top2 is
  PORT( clk                               :   IN    std_logic;
        reset                             :   IN    std_logic;
        clk_enable                        :   IN    std_logic;
        V_uv                              :   IN    std_logic_vector(13 DOWNTO 0);  -- ufix14
        V_uw                              :   IN    std_logic_vector(13 DOWNTO 0);  -- ufix14
        iout_u                            :   IN    std_logic_vector(13 DOWNTO 0);  -- ufix14
        iout_v                            :   IN    std_logic_vector(13 DOWNTO 0);  -- ufix14
        iL_u                              :   IN    std_logic_vector(13 DOWNTO 0);  -- ufix14
        iL_v                              :   IN    std_logic_vector(13 DOWNTO 0);  -- ufix14
        switch_control                    :   IN    std_logic_vector(13 DOWNTO 0);  -- ufix14
        Up_current                        :   IN    std_logic_vector(13 DOWNTO 0);  -- ufix14
        Un_current                        :   IN    std_logic_vector(13 DOWNTO 0);  -- ufix14
        Vp_current                        :   IN    std_logic_vector(13 DOWNTO 0);  -- ufix14
        Vn_current                        :   IN    std_logic_vector(13 DOWNTO 0);  -- ufix14
        Wp_current                        :   IN    std_logic_vector(13 DOWNTO 0);  -- ufix14
        UW_current                        :   IN    std_logic_vector(13 DOWNTO 0);  -- ufix14
        CV_U1                             :   IN    std_logic_vector(13 DOWNTO 0);  -- ufix14
        CV_U2                             :   IN    std_logic_vector(13 DOWNTO 0);  -- ufix14
        CV_U3                             :   IN    std_logic_vector(13 DOWNTO 0);  -- ufix14
        CV_U4                             :   IN    std_logic_vector(13 DOWNTO 0);  -- ufix14
        CV_U5                             :   IN    std_logic_vector(13 DOWNTO 0);  -- ufix14
        CV_U6                             :   IN    std_logic_vector(13 DOWNTO 0);  -- ufix14
        CV_U7                             :   IN    std_logic_vector(13 DOWNTO 0);  -- ufix14
        CV_U8                             :   IN    std_logic_vector(13 DOWNTO 0);  -- ufix14
        CV_V1                             :   IN    std_logic_vector(13 DOWNTO 0);  -- ufix14
        CV_V2                             :   IN    std_logic_vector(13 DOWNTO 0);  -- ufix14
        CV_V3                             :   IN    std_logic_vector(13 DOWNTO 0);  -- ufix14
        CV_V4                             :   IN    std_logic_vector(13 DOWNTO 0);  -- ufix14
        CV_V5                             :   IN    std_logic_vector(13 DOWNTO 0);  -- ufix14
        CV_V6                             :   IN    std_logic_vector(13 DOWNTO 0);  -- ufix14
        CV_V7                             :   IN    std_logic_vector(13 DOWNTO 0);  -- ufix14
        CV_V8                             :   IN    std_logic_vector(13 DOWNTO 0);  -- ufix14
        CV_W1                             :   IN    std_logic_vector(13 DOWNTO 0);  -- ufix14
        CV_W2                             :   IN    std_logic_vector(13 DOWNTO 0);  -- ufix14
        CV_W3                             :   IN    std_logic_vector(13 DOWNTO 0);  -- ufix14
        CV_W4                             :   IN    std_logic_vector(13 DOWNTO 0);  -- ufix14
        CV_W5                             :   IN    std_logic_vector(13 DOWNTO 0);  -- ufix14
        CV_W6                             :   IN    std_logic_vector(13 DOWNTO 0);  -- ufix14
        CV_W7                             :   IN    std_logic_vector(13 DOWNTO 0);  -- ufix14
        CV_W8                             :   IN    std_logic_vector(13 DOWNTO 0);  -- ufix14
        ce_out                            :   OUT   std_logic;
        MMC_PWM_OUT1                      :   OUT   std_logic_vector(47 DOWNTO 0);  -- ufix48
        MMC_PWM_OUT2                      :   OUT   std_logic_vector(47 DOWNTO 0);  -- ufix48
		  
        Vuv                               :   OUT   std_logic_vector(19 DOWNTO 0);  -- ufix14
        Vvw                               :   OUT   std_logic_vector(19 DOWNTO 0);  -- ufix14
        Ioutu                             :   OUT   std_logic_vector(19 DOWNTO 0);  -- ufix14
        Ioutv                             :   OUT   std_logic_vector(19 DOWNTO 0);  -- ufix14
        ILu                               :   OUT   std_logic_vector(19 DOWNTO 0);  -- ufix14
        ILv                               :   OUT   std_logic_vector(19 DOWNTO 0);  -- ufix14
        Vrefa                             :   OUT   std_logic_vector(19 DOWNTO 0);  -- sfix14
        Vrefb                             :   OUT   std_logic_vector(19 DOWNTO 0);  -- sfix14
        deltaTa                           :   OUT   std_logic_vector(15 DOWNTO 0);  -- int16
        deltaTb                           :   OUT   std_logic_vector(15 DOWNTO 0);  -- int16
        deltaTa_hold                      :   OUT   std_logic_vector(15 DOWNTO 0);  -- int16
        deltaTb_hold                      :   OUT   std_logic_vector(15 DOWNTO 0);  -- int16
        switch_control1                   :   OUT   std_logic_vector(13 DOWNTO 0)  -- ufix14
        );
		end component;


begin

-- ====================================================
-- Generate Internal Clock (100MHz)
--   20MHzNbN5{{CLK100
--   {WbNBUS_CLK_INgp
-- ====================================================

	INBUF_CLK:IBUFG
		port map(
			O	=> CLK20B,
			I	=> CLK20_IN
			);

-- DCM_SP: Digital Clock Manager
-- Spartan-6
-- Xilinx HDL Libraries Guide, version 14.3
	DCM_SP_inst : DCM_SP
		generic map (
			CLKDV_DIVIDE 			=> 2.0, 					-- CLKDV divide value
																-- (1.5,2,2.5,3,3.5,4,4.5,5,5.5,6,6.5,7,7.5,8,9,10,11,12,13,14,15,16).
			CLKFX_DIVIDE 			=> 1, 						-- Divide value on CLKFX outputs - D - (1-32)
			CLKFX_MULTIPLY 			=> 5, 						-- Multiply value on CLKFX outputs - M - (2-32)
			CLKIN_DIVIDE_BY_2 		=> FALSE, 					-- CLKIN divide by two (TRUE/FALSE)
			CLKIN_PERIOD 			=> 50.0, 					-- Input clock period specified in nS
			CLKOUT_PHASE_SHIFT 		=> "NONE", 					-- Output phase shift (NONE, FIXED, VARIABLE)
			CLK_FEEDBACK 			=> "1X", 					-- Feedback source (NONE, 1X, 2X)
			DESKEW_ADJUST 			=> "SYSTEM_SYNCHRONOUS", 	-- SYSTEM_SYNCHRNOUS or SOURCE_SYNCHRONOUS
			PHASE_SHIFT 			=> 0, 						-- Amount of fixed phase shift (-255 to 255)
			STARTUP_WAIT 			=> FALSE 					-- Delay config DONE until DCM_SP LOCKED (TRUE/FALSE)
		)
		port map (
			CLK0 		=> CLK0, 						-- 1-bit output: 0 degree clock output
			CLK180 		=> open, 						-- 1-bit output: 180 degree clock output
			CLK270 		=> open, 						-- 1-bit output: 270 degree clock output
			CLK2X 		=> open, 						-- 1-bit output: 2X clock frequency clock output
			CLK2X180 	=> open, 						-- 1-bit output: 2X clock frequency, 180 degree clock output
			CLK90 		=> open, 						-- 1-bit output: 90 degree clock output
			CLKDV 		=> open, 						-- 1-bit output: Divided clock output
			CLKFX 		=> CLK100FX, 					-- 1-bit output: Digital Frequency Synthesizer output (DFS)
			CLKFX180 	=> open, 						-- 1-bit output: 180 degree CLKFX output
			LOCKED 		=> LOCKED_100, 					-- 1-bit output: DCM_SP Lock Output
			PSDONE 		=> open, 						-- 1-bit output: Phase shift done output
			STATUS 		=> LOCK_STATUS_100, 			-- 8-bit output: DCM_SP status output
			CLKFB 		=> CLKFB, 						-- 1-bit input: Clock feedback input
			CLKIN 		=> CLK20B, 						-- 1-bit input: Clock input
			DSSEN 		=> '0', 						-- 1-bit input: Unsupported, specify to GND.
			PSCLK 		=> '0', 						-- 1-bit input: Phase shift clock input
			PSEN 		=> '0', 						-- 1-bit input: Phase shift enable
			PSINCDEC 	=> '0', 						-- 1-bit input: Phase shift increment/decrement input
			RST 		=> RESET_DCM_100 				-- 1-bit input: Active high reset input
		);
-- End of DCM_SP_inst instantiation

	CLKF_BUF: BUFG
		port map(
			O => CLKFB,
			I => CLK0
			);

	CLK100OUT_BUF: BUFG
	port map(
			O	=> CLK100,
			I	=> CLK100FX
			);



	DCM_SP_inst2 : DCM_SP
		generic map (
			CLKDV_DIVIDE 			=> 2.0, 					-- CLKDV divide value
																-- (1.5,2,2.5,3,3.5,4,4.5,5,5.5,6,6.5,7,7.5,8,9,10,11,12,13,14,15,16).
			CLKFX_DIVIDE 			=> 2, 						-- Divide value on CLKFX outputs - D - (1-32)
			CLKFX_MULTIPLY 			=> 3, 						-- Multiply value on CLKFX outputs - M - (2-32)
			CLKIN_DIVIDE_BY_2 		=> FALSE, 					-- CLKIN divide by two (TRUE/FALSE)
			CLKIN_PERIOD 			=> 10.0, 					-- Input clock period specified in nS
			CLKOUT_PHASE_SHIFT 		=> "NONE", 					-- Output phase shift (NONE, FIXED, VARIABLE)
			CLK_FEEDBACK 			=> "1X", 					-- Feedback source (NONE, 1X, 2X)
			DESKEW_ADJUST 			=> "SYSTEM_SYNCHRONOUS", 	-- SYSTEM_SYNCHRNOUS or SOURCE_SYNCHRONOUS
			PHASE_SHIFT 			=> 0, 						-- Amount of fixed phase shift (-255 to 255)
			STARTUP_WAIT 			=> FALSE 					-- Delay config DONE until DCM_SP LOCKED (TRUE/FALSE)
		)
		port map (
			CLK0 		=> CLK0_2, 						-- 1-bit output: 0 degree clock output
			CLK180 		=> open, 						-- 1-bit output: 180 degree clock output
			CLK270 		=> open, 						-- 1-bit output: 270 degree clock output
			CLK2X 		=> open, 						-- 1-bit output: 2X clock frequency clock output
			CLK2X180 	=> open, 						-- 1-bit output: 2X clock frequency, 180 degree clock output
			CLK90 		=> open, 						-- 1-bit output: 90 degree clock output
			CLKDV 		=> open, 						-- 1-bit output: Divided clock output
			CLKFX 		=> CLK150FX, 					-- 1-bit output: Digital Frequency Synthesizer output (DFS)
			CLKFX180 	=> open, 						-- 1-bit output: 180 degree CLKFX output
			LOCKED 		=> LOCKED_150, 					-- 1-bit output: DCM_SP Lock Output
			PSDONE 		=> open, 						-- 1-bit output: Phase shift done output
			STATUS 		=> LOCK_STATUS_150, 			-- 8-bit output: DCM_SP status output
			CLKFB 		=> CLKFB_2,						-- 1-bit input: Clock feedback input
			CLKIN 		=> CLK100, 						-- 1-bit input: Clock input
			DSSEN 		=> '0', 						-- 1-bit input: Unsupported, specify to GND.
			PSCLK 		=> '0', 						-- 1-bit input: Phase shift clock input
			PSEN 		=> '0', 						-- 1-bit input: Phase shift enable
			PSINCDEC 	=> '0', 						-- 1-bit input: Phase shift increment/decrement input
			RST 		=> RESET_DCM_150 				-- 1-bit input: Active high reset input
		);
-- End of DCM_SP_inst instantiation

	CLKF_BUF_2: BUFG
		port map(
			O => CLKFB_2,
			I => CLK0_2
			);

	CLK150OUT_BUF: BUFG
	port map(
			O	=> CLK150,
			I	=> CLK150FX
			);



	process(CLK20B)
	begin
		if CLK20B'event and CLK20B = '1' then
			if LOCKED_100 = '1' and LOCK_STATUS_100(1) = '0' and LOCK_STATUS_100(2) = '0' and LOCK_STATUS_100(7) = '0' then
				LOCK_CNT_100	<= "000" & "00000000" & "00000000";
				RESET_DCM_100	<= '0';
				DCM_OK_100		<= '1';
			else
				DCM_OK_100		<= '0';
				LOCK_CNT_100	<= LOCK_CNT_100 + 1;
				if LOCK_CNT_100 =  "000" & "00000000" & "00000001" then
					RESET_DCM_100	<= '1';
				else
					RESET_DCM_100	<= '0';
				end if;
			end if;
		end if;
	end process;

	process(CLK20B)
	begin
		if CLK20B'event and CLK20B = '1' then
			if LOCKED_150 = '1' and LOCK_STATUS_150(1) = '0' and LOCK_STATUS_150(2) = '0' and LOCK_STATUS_150(7) = '0' then
				LOCK_CNT_150	<= "000" & "00000000" & "00000000";
				RESET_DCM_150	<= '0';
				DCM_OK_150		<= '1';
			else
				DCM_OK_150		<= '0';
				LOCK_CNT_150	<= LOCK_CNT_150 + 1;
				if LOCK_CNT_150 =  "000" & "00000000" & "00000001" then
					RESET_DCM_150	<= '1';
				else
					RESET_DCM_150	<= '0';
				end if;
			end if;
		end if;
	end process;

-- ====================================================
-- Generate Internal RESET
--   RESET    : ZbgM
--   RESET_BB : ZbgM(CLK100)
-- ====================================================

	RESET	<= '0' when nRESET_IN = '1' and DCM_OK_100 = '1' and RESET_DCM_100 = '0' and DCM_OK_150 = '1' and RESET_DCM_150 = '0' else '1';
	
	process(CLK100)
	begin
		if CLK100'event and CLK100 = '1' then
			RESET_B		<= RESET;
			RESET_BB	<= RESET_B;
		end if;
	end process;


-- ====================================================
-- Input Filter
--   tB^
--   tB^Kv
-- ====================================================
	LPF_DI0: BBB_LPF	generic map ( LPF_TIME => 5,	INVERT => 0 )	port map ( LPF_IN => DI_IN(0),		CLK => CLK100,		LPF_OUT => DI_LPF(0) ) ;
	LPF_DI1: BBB_LPF	generic map ( LPF_TIME => 5,	INVERT => 0 )	port map ( LPF_IN => DI_IN(1),		CLK => CLK100,		LPF_OUT => DI_LPF(1) ) ;
	LPF_DI2: BBB_LPF	generic map ( LPF_TIME => 5,	INVERT => 0 )	port map ( LPF_IN => DI_IN(2),		CLK => CLK100,		LPF_OUT => DI_LPF(2) ) ;
	LPF_DI3: BBB_LPF	generic map ( LPF_TIME => 5,	INVERT => 0 )	port map ( LPF_IN => DI_IN(3),		CLK => CLK100,		LPF_OUT => DI_LPF(3) ) ;
	LPF_DI4: BBB_LPF	generic map ( LPF_TIME => 5,	INVERT => 0 )	port map ( LPF_IN => DI_IN(4),		CLK => CLK100,		LPF_OUT => DI_LPF(4) ) ;
	LPF_DI5: BBB_LPF	generic map ( LPF_TIME => 5,	INVERT => 0 )	port map ( LPF_IN => DI_IN(5),		CLK => CLK100,		LPF_OUT => DI_LPF(5) ) ;
	LPF_DI6: BBB_LPF	generic map ( LPF_TIME => 5,	INVERT => 0 )	port map ( LPF_IN => DI_IN(6),		CLK => CLK100,		LPF_OUT => DI_LPF(6) ) ;
	LPF_DI7: BBB_LPF	generic map ( LPF_TIME => 5,	INVERT => 0 )	port map ( LPF_IN => DI_IN(7),		CLK => CLK100,		LPF_OUT => DI_LPF(7) ) ;
	LPF_DI8: BBB_LPF	generic map ( LPF_TIME => 5,	INVERT => 0 )	port map ( LPF_IN => DI_IN(8),		CLK => CLK100,		LPF_OUT => DI_LPF(8) ) ;
	LPF_DI9: BBB_LPF	generic map ( LPF_TIME => 5,	INVERT => 0 )	port map ( LPF_IN => DI_IN(9),		CLK => CLK100,		LPF_OUT => DI_LPF(9) ) ;
	LPF_DI10: BBB_LPF	generic map ( LPF_TIME => 5,	INVERT => 0 )	port map ( LPF_IN => DI_IN(10),		CLK => CLK100,		LPF_OUT => DI_LPF(10) ) ;
	LPF_DI11: BBB_LPF	generic map ( LPF_TIME => 5,	INVERT => 0 )	port map ( LPF_IN => DI_IN(11),		CLK => CLK100,		LPF_OUT => DI_LPF(11) ) ;
	LPF_DI12: BBB_LPF	generic map ( LPF_TIME => 5,	INVERT => 0 )	port map ( LPF_IN => DI_IN(12),		CLK => CLK100,		LPF_OUT => DI_LPF(12) ) ;
	LPF_DI13: BBB_LPF	generic map ( LPF_TIME => 5,	INVERT => 0 )	port map ( LPF_IN => DI_IN(13),		CLK => CLK100,		LPF_OUT => DI_LPF(13) ) ;
	LPF_DI14: BBB_LPF	generic map ( LPF_TIME => 5,	INVERT => 0 )	port map ( LPF_IN => DI_IN(14),		CLK => CLK100,		LPF_OUT => DI_LPF(14) ) ;
	LPF_DI15: BBB_LPF	generic map ( LPF_TIME => 5,	INVERT => 0 )	port map ( LPF_IN => DI_IN(15),		CLK => CLK100,		LPF_OUT => DI_LPF(15) ) ;

	LPF_DSW0: BBB_LPF	generic map ( LPF_TIME => 5,	INVERT => 0 )	port map ( LPF_IN => DSW_IN(0),		CLK => CLK100,		LPF_OUT => DSW_LPF(0) ) ;
	LPF_DSW1: BBB_LPF	generic map ( LPF_TIME => 5,	INVERT => 0 )	port map ( LPF_IN => DSW_IN(1),		CLK => CLK100,		LPF_OUT => DSW_LPF(1) ) ;
	LPF_DSW2: BBB_LPF	generic map ( LPF_TIME => 5,	INVERT => 0 )	port map ( LPF_IN => DSW_IN(2),		CLK => CLK100,		LPF_OUT => DSW_LPF(2) ) ;
	LPF_DSW3: BBB_LPF	generic map ( LPF_TIME => 5,	INVERT => 0 )	port map ( LPF_IN => DSW_IN(3),		CLK => CLK100,		LPF_OUT => DSW_LPF(3) ) ;
	LPF_DSW4: BBB_LPF	generic map ( LPF_TIME => 5,	INVERT => 0 )	port map ( LPF_IN => DSW_BOOT_IN(0),		CLK => CLK100,		LPF_OUT => DSW_BOOT_LPF(0) ) ;
	LPF_DSW5: BBB_LPF	generic map ( LPF_TIME => 5,	INVERT => 0 )	port map ( LPF_IN => DSW_BOOT_IN(1),		CLK => CLK100,		LPF_OUT => DSW_BOOT_LPF(1) ) ;

	LPF_DO0: BBB_LPF	generic map ( LPF_TIME => 5,	INVERT => 0 )	port map ( LPF_IN => DO_IN(0),		CLK => CLK100,		LPF_OUT => DO_LPF(0) ) ;
	LPF_DO1: BBB_LPF	generic map ( LPF_TIME => 5,	INVERT => 0 )	port map ( LPF_IN => DO_IN(1),		CLK => CLK100,		LPF_OUT => DO_LPF(1) ) ;
	LPF_DO2: BBB_LPF	generic map ( LPF_TIME => 5,	INVERT => 0 )	port map ( LPF_IN => DO_IN(2),		CLK => CLK100,		LPF_OUT => DO_LPF(2) ) ;
	LPF_DO3: BBB_LPF	generic map ( LPF_TIME => 5,	INVERT => 0 )	port map ( LPF_IN => DO_IN(3),		CLK => CLK100,		LPF_OUT => DO_LPF(3) ) ;
	LPF_DO4: BBB_LPF	generic map ( LPF_TIME => 5,	INVERT => 0 )	port map ( LPF_IN => DO_IN(4),		CLK => CLK100,		LPF_OUT => DO_LPF(4) ) ;
	LPF_DO5: BBB_LPF	generic map ( LPF_TIME => 5,	INVERT => 0 )	port map ( LPF_IN => DO_IN(5),		CLK => CLK100,		LPF_OUT => DO_LPF(5) ) ;
	LPF_DO6: BBB_LPF	generic map ( LPF_TIME => 5,	INVERT => 0 )	port map ( LPF_IN => DO_IN(6),		CLK => CLK100,		LPF_OUT => DO_LPF(6) ) ;
	LPF_DO7: BBB_LPF	generic map ( LPF_TIME => 5,	INVERT => 0 )	port map ( LPF_IN => DO_IN(7),		CLK => CLK100,		LPF_OUT => DO_LPF(7) ) ;
	LPF_DO8: BBB_LPF	generic map ( LPF_TIME => 5,	INVERT => 0 )	port map ( LPF_IN => DO_IN(8),		CLK => CLK100,		LPF_OUT => DO_LPF(8) ) ;
	LPF_DO9: BBB_LPF	generic map ( LPF_TIME => 5,	INVERT => 0 )	port map ( LPF_IN => DO_IN(9),		CLK => CLK100,		LPF_OUT => DO_LPF(9) ) ;
	LPF_DO10: BBB_LPF	generic map ( LPF_TIME => 5,	INVERT => 0 )	port map ( LPF_IN => DO_IN(10),		CLK => CLK100,		LPF_OUT => DO_LPF(10) ) ;
	LPF_DO11: BBB_LPF	generic map ( LPF_TIME => 5,	INVERT => 0 )	port map ( LPF_IN => DO_IN(11),		CLK => CLK100,		LPF_OUT => DO_LPF(11) ) ;
	LPF_DO12: BBB_LPF	generic map ( LPF_TIME => 5,	INVERT => 0 )	port map ( LPF_IN => DO_IN(12),		CLK => CLK100,		LPF_OUT => DO_LPF(12) ) ;
	LPF_DO13: BBB_LPF	generic map ( LPF_TIME => 5,	INVERT => 0 )	port map ( LPF_IN => DO_IN(13),		CLK => CLK100,		LPF_OUT => DO_LPF(13) ) ;
	LPF_DO14: BBB_LPF	generic map ( LPF_TIME => 5,	INVERT => 0 )	port map ( LPF_IN => DO_IN(14),		CLK => CLK100,		LPF_OUT => DO_LPF(14) ) ;
	LPF_DO15: BBB_LPF	generic map ( LPF_TIME => 5,	INVERT => 0 )	port map ( LPF_IN => DO_IN(15),		CLK => CLK100,		LPF_OUT => DO_LPF(15) ) ;



-- ====================================================
-- BUS
--   oXRnWMgp
--   /RD, /WRsgp
-- ====================================================

	process(CLK100)
	begin
		if CLK100'event and CLK100 = '1' then
			BUS_CS2_B		<=	not nBUS_CS2_IN;
			BUS_CS2_BB		<=	BUS_CS2_B;
			BUS_CS3_B		<=	not nBUS_CS3_IN;
			BUS_CS3_BB		<=	BUS_CS3_B;
			BUS_RD_B		<=	not nBUS_RD_IN;
			BUS_RD_BB		<=	BUS_RD_B;
			BUS_WR_B		<=	not nBUS_WR_IN;
			BUS_WR_BB		<=	BUS_WR_B;
			BUS_RnW_B		<=	BUS_RnW_IN;
			BUS_RnW_BB		<=	BUS_RnW_B;
			BUS_A_B			<=	BUS_A_IN & BUS_BA1_IN;
			BUS_A_BB		<=	BUS_A_B;
			BUS_DW_B		<=	BUS_D_IO;
			BUS_DW_BB		<=	BUS_DW_B;

			if RESET_BB = '1' then
				BUS_CNT	<= 0;
			elsif BUS_CS2_BB /= '1' then
				BUS_CNT	<= 0;
			elsif BUS_CNT = 20 then
				BUS_CNT	<= 20;
			else
				BUS_CNT	<= BUS_CNT +1;
			end if;
		end if;
	end process;
	
	
	BUS_RD_0		<= '1' when BUS_CS2_BB = '1' and BUS_RnW_BB = '1' and BUS_CNT = 1 and BUS_A_BB(15 downto 8) = x"00" else '0';
--	BUS_WR_0		<= '1' when BUS_CS2_BB = '1' and BUS_RnW_BB = '0' and BUS_CNT = 3 and BUS_A_BB(15 downto 8) = x"00" else '0';
	BUS_RD_1		<= '1' when BUS_CS2_BB = '1' and BUS_RnW_BB = '1' and BUS_CNT = 1 and BUS_A_BB(15 downto 8) = x"01" else '0';
	BUS_WR_1		<= '1' when BUS_CS2_BB = '1' and BUS_RnW_BB = '0' and BUS_CNT = 3 and BUS_A_BB(15 downto 8) = x"01" else '0';
	BUS_RD_2		<= '1' when BUS_CS2_BB = '1' and BUS_RnW_BB = '1' and BUS_CNT = 1 and BUS_A_BB(15 downto 8) = x"02" else '0';
	BUS_WR_2		<= '1' when BUS_CS2_BB = '1' and BUS_RnW_BB = '0' and BUS_CNT = 3 and BUS_A_BB(15 downto 8) = x"02" else '0';
	BUS_RD_3		<= '1' when BUS_CS2_BB = '1' and BUS_RnW_BB = '1' and BUS_CNT = 1 and BUS_A_BB(15 downto 8) = x"03" else '0';
	BUS_WR_3		<= '1' when BUS_CS2_BB = '1' and BUS_RnW_BB = '0' and BUS_CNT = 3 and BUS_A_BB(15 downto 8) = x"03" else '0';
	BUS_RD_FF		<= '1' when BUS_CS2_BB = '1' and BUS_RnW_BB = '1' and BUS_CNT = 1 and BUS_A_BB(15 downto 8) = x"FF" else '0';
	BUS_WR_FF		<= '1' when BUS_CS2_BB = '1' and BUS_RnW_BB = '0' and BUS_CNT = 3 and BUS_A_BB(15 downto 8) = x"FF" else '0';
	
	BUS_D_IO	<=			BUS_DR_0	when nBUS_CS2_IN = '0' and BUS_RnW_IN = '1' and BUS_A_BB(15 downto 8) = x"00"
					else	BUS_DR_1	when nBUS_CS2_IN = '0' and BUS_RnW_IN = '1' and BUS_A_BB(15 downto 8) = x"01" and BUS_A_BB(7 downto 5) = "000"
					else	BUS_DR_1_2	when nBUS_CS2_IN = '0' and BUS_RnW_IN = '1' and BUS_A_BB(15 downto 8) = x"01" and BUS_A_BB(7 downto 5) = "001"
					else	BUS_DR_1_3	when nBUS_CS2_IN = '0' and BUS_RnW_IN = '1' and BUS_A_BB(15 downto 8) = x"01"
					else	BUS_DR_2	when nBUS_CS2_IN = '0' and BUS_RnW_IN = '1' and BUS_A_BB(15 downto 8) = x"02" and BUS_A_BB(7 downto 6) = "00"
					else	BUS_DR_2_2	when nBUS_CS2_IN = '0' and BUS_RnW_IN = '1' and BUS_A_BB(15 downto 8) = x"02"
					else	BUS_DR_3	when nBUS_CS2_IN = '0' and BUS_RnW_IN = '1' and BUS_A_BB(15 downto 8) = x"03"
					else	BUS_DR_FF	when nBUS_CS2_IN = '0' and BUS_RnW_IN = '1' and BUS_A_BB(15 downto 8) = x"FF"
					else	"ZZZZZZZZ" & "ZZZZZZZZ";


	process(CLK100)
	begin
		if CLK100'event and CLK100 = '1' then
			if BUS_RD_0 = '1' then
				case BUS_A_BB(7 downto 0) is
					when x"00"	=> BUS_DR_0	<=	DI_LPF;
					when x"01"	=> BUS_DR_0	<=	DO_LPF(15 downto 0);
					when x"02"	=> BUS_DR_0	<=	"00000000" & "0000" & DSW_LPF(3 downto 0);
					when x"03"	=> BUS_DR_0	<=	"00000000" & "000000" & DSW_BOOT_LPF(1 downto 0);
--					when x"04"	=> BUS_DR_0	<=	"00000000" & "0000" & LED;

					when x"FF"	=> BUS_DR_0	<=	CONV_std_logic_vector (FPGA_VERSION_MAJOR, 4) 
													& CONV_std_logic_vector (FPGA_VERSION_MINOR, 4) 
													& CONV_std_logic_vector (FPGA_VERSION_BUILD, 8);
					when others => BUS_DR_0	<=	"00000000" & "00000000";
				end case;
			end if;

--			if RESET_BB = '1' then
--				LED		<= "0000";
--			elsif BUS_WR_0 = '1' then
--				case BUS_A_BB(7 downto 0) is
--					when x"04"	=> LED		<= BUS_DW_BB(3 downto 0);
--					when others	=> null;
--				end case;
--			else
--				null;
--			end if;
		end if;
	end process;



	process(CLK100)
	begin
		if CLK100'event and CLK100 = '1' then
			if BUS_RD_FF = '1' then
				case BUS_A_BB(7 downto 0) is
					when x"FE"	=> BUS_DR_FF	<=	"00000000"&"000" & DUMMY;
					when x"FF"	=> BUS_DR_FF	<=	SPR;
					when others => BUS_DR_FF	<=	"00000000" & "00000000";
				end case;
			end if;

			if RESET_BB = '1' then
				SPR		<= "00000000" & "00000000";
			elsif BUS_WR_FF = '1' then
				case BUS_A_BB(7 downto 0) is
					when x"FF"	=> SPR		<= BUS_DW_BB;
					when others	=> null;
				end case;
			else
				null;
			end if;
		end if;
	end process;



-- ====================================================
-- MMC_PWM
-- ====================================================

	MMC_PWM_GEN :MMC_PWM
		port map(
			CLK100		=> CLK100,
			RESET		=> RESET_BB,
			A			=> BUS_A_BB(7 downto 0),
			DW			=> BUS_DW_BB,
			DR			=> BUS_DR_1,
			DR2			=> BUS_DR_1_2,
			DR3			=> BUS_DR_1_3,
			RD			=> BUS_RD_1,
			WR			=> BUS_WR_1,
			
--			nMMC_PWM1	=> nMMC_PWM1_OUT,
--			nMMC_PWM2	=> nMMC_PWM2_OUT,
			--(20190716)
			nMMC_PWM1	=> nMMC_PWM1_buf,
			nMMC_PWM2	=> nMMC_PWM2_buf,
			nINT_SIG	=> nINT_SIG,
			
			PWM_ON		=> PWM_ON,
			GB			=> GB,
		
		--20190814_PFPGA_PIfobOp--
			MRate_UP1	=>	MRate_UP1_B,		--UT
			MRate_UP2	=>	MRate_UP2_B,
			MRate_UP3	=>	MRate_UP3_B,
			MRate_UP4	=>	MRate_UP4_B,
			MRate_UN1	=>	MRate_UN1_B,
			MRate_UN2	=>	MRate_UN2_B,
			MRate_UN3	=>	MRate_UN3_B,
			MRate_UN4	=>	MRate_UN4_B,
			MRate_VP1	=>	MRate_VP1_B,		--VT
			MRate_VP2	=>	MRate_VP2_B,
			MRate_VP3	=>	MRate_VP3_B,
			MRate_VP4	=>	MRate_VP4_B,
			MRate_VN1	=>	MRate_VN1_B,
			MRate_VN2	=>	MRate_VN2_B,
			MRate_VN3	=>	MRate_VN3_B,
			MRate_VN4	=>	MRate_VN4_B,
			MRate_WP1	=>	MRate_WP1_B,		--WT
			MRate_WP2	=>	MRate_WP2_B,
			MRate_WP3	=>	MRate_WP3_B,
			MRate_WP4	=>	MRate_WP4_B,
			MRate_WN1	=>	MRate_WN1_B,
			MRate_WN2	=>	MRate_WN2_B,
			MRate_WN3	=>	MRate_WN3_B,
			MRate_WN4	=>	MRate_WN4_B,
		-----------------------------------
			--ver2.6.4--
			REF_UP1_vector => REF_UP1_vector_B,
			REF_UP2_vector => REF_UP2_vector_B,
			REF_UP3_vector => REF_UP3_vector_B,
			REF_UP4_vector	=> REF_UP4_vector_B,
			REF_UN1_vector => REF_UN1_vector_B,
			REF_UN2_vector	=> REF_UN2_vector_B,
			REF_UN3_vector	=> REF_UN3_vector_B,
			REF_UN4_vector	=> REF_UN4_vector_B,
			REF_VP1_vector => REF_VP1_vector_B,	
			REF_VP2_vector	=> REF_VP2_vector_B,
			REF_VP3_vector	=> REF_VP3_vector_B,
			REF_VP4_vector	=> REF_VP4_vector_B,
			REF_VN1_vector	=> REF_VN1_vector_B,
			REF_VN2_vector	=> REF_VN2_vector_B,
			REF_VN3_vector	=> REF_VN3_vector_B,
			REF_VN4_vector	=> REF_VN4_vector_B
			----------------
			
			);


-- ====================================================
-- ADC
-- ====================================================

	process(CLK150)
	begin
		if CLK150'event and CLK150 = '1' then
			ADC_GR0_CS_B	<= not nADC_GR0_CS_OUT;
			ADC_GR1_CS_B	<= not nADC_GR1_CS_OUT;
			ADC_GR2_CS_B	<= not nADC_GR2_CS_OUT;
			ADC_GR3_CS_B	<= not nADC_GR3_CS_OUT;
			ADC_GR4_CS_B	<= not nADC_GR4_CS_OUT;
			ADC_GR5_CS_B	<= not nADC_GR5_CS_OUT;
			ADC_GR6_CS_B	<= not nADC_GR6_CS_OUT;
			ADC_GR7_CS_B	<= not nADC_GR7_CS_OUT;
			ADC_GR8_CS_B	<= not nADC_GR8_CS_OUT;
			ADC_GR9_CS_B	<= not nADC_GR9_CS_OUT;
			
			ADC_GR0_SCLK_B	<= ADC_GR0_SCLK_OUT;
			ADC_GR1_SCLK_B	<= ADC_GR1_SCLK_OUT;
			ADC_GR2_SCLK_B	<= ADC_GR2_SCLK_OUT;
			ADC_GR3_SCLK_B	<= ADC_GR3_SCLK_OUT;
			ADC_GR4_SCLK_B	<= ADC_GR4_SCLK_OUT;
			ADC_GR5_SCLK_B	<= ADC_GR5_SCLK_OUT;
			ADC_GR6_SCLK_B	<= ADC_GR6_SCLK_OUT;
			ADC_GR7_SCLK_B	<= ADC_GR7_SCLK_OUT;
			ADC_GR8_SCLK_B	<= ADC_GR8_SCLK_OUT;
			ADC_GR9_SCLK_B	<= ADC_GR9_SCLK_OUT;
			
			ADC_GR0_DATA_B	<= ADC_GR0_DATA_IN;
			ADC_GR1_DATA_B	<= ADC_GR1_DATA_IN;
			ADC_GR2_DATA_B	<= ADC_GR2_DATA_IN;
			ADC_GR3_DATA_B	<= ADC_GR3_DATA_IN;
			ADC_GR4_DATA_B	<= ADC_GR4_DATA_IN;
			ADC_GR5_DATA_B	<= ADC_GR5_DATA_IN;
			ADC_GR6_DATA_B	<= ADC_GR6_DATA_IN;
			ADC_GR7_DATA_B	<= ADC_GR7_DATA_IN;
			ADC_GR8_DATA_B	<= ADC_GR8_DATA_IN;
			ADC_GR9_DATA_B	<= ADC_GR9_DATA_IN;
		end if;
	end process;

	ADC :AD7357
		port map(
			CLK150		=> CLK150,
			CLK100		=> CLK100,
			RESET		=> RESET_BB,
			A			=> BUS_A_BB(7 downto 0),
			DW			=> BUS_DW_BB,
			DR			=> BUS_DR_2,
			DR2			=> BUS_DR_2_2,
			RD			=> BUS_RD_2,
			WR			=> BUS_WR_2,
			
			AD_CS		=> AD_CS,
			AD_CLK		=> AD_CLK,
			AD_CH_EN	=> AD_CH_EN,
			
			ADC_GR0_CS		=> ADC_GR0_CS_B,
			ADC_GR0_SCLK	=> ADC_GR0_SCLK_B,
			ADC_GR0_DATA	=> ADC_GR0_DATA_B,

			ADC_GR1_CS		=> ADC_GR1_CS_B,
			ADC_GR1_SCLK	=> ADC_GR1_SCLK_B,
			ADC_GR1_DATA	=> ADC_GR1_DATA_B,

			ADC_GR2_CS		=> ADC_GR2_CS_B,
			ADC_GR2_SCLK	=> ADC_GR2_SCLK_B,
			ADC_GR2_DATA	=> ADC_GR2_DATA_B,

			ADC_GR3_CS		=> ADC_GR3_CS_B,
			ADC_GR3_SCLK	=> ADC_GR3_SCLK_B,
			ADC_GR3_DATA	=> ADC_GR3_DATA_B,

			ADC_GR4_CS		=> ADC_GR4_CS_B,
			ADC_GR4_SCLK	=> ADC_GR4_SCLK_B,
			ADC_GR4_DATA	=> ADC_GR4_DATA_B,

			ADC_GR5_CS		=> ADC_GR5_CS_B,
			ADC_GR5_SCLK	=> ADC_GR5_SCLK_B,
			ADC_GR5_DATA	=> ADC_GR5_DATA_B,

			ADC_GR6_CS		=> ADC_GR6_CS_B,
			ADC_GR6_SCLK	=> ADC_GR6_SCLK_B,
			ADC_GR6_DATA	=> ADC_GR6_DATA_B,

			ADC_GR7_CS		=> ADC_GR7_CS_B,
			ADC_GR7_SCLK	=> ADC_GR7_SCLK_B,
			ADC_GR7_DATA	=> ADC_GR7_DATA_B,

			ADC_GR8_CS		=> ADC_GR8_CS_B,
			ADC_GR8_SCLK	=> ADC_GR8_SCLK_B,
			ADC_GR8_DATA	=> ADC_GR8_DATA_B,

			ADC_GR9_CS		=> ADC_GR9_CS_B,
			ADC_GR9_SCLK	=> ADC_GR9_SCLK_B,
			ADC_GR9_DATA	=> ADC_GR9_DATA_B,
			
			ADC_ERR			=> ADC_ERR,
			
			--ver2--
			Vc_UP1			=>	Vc_UP1_B,		--RfTdVc
			Vc_UP2			=>	Vc_UP2_B,
			Vc_UP3			=>	Vc_UP3_B,
			Vc_UP4			=>	Vc_UP4_B,
			Vc_UN1			=>	Vc_UN1_B,
			Vc_UN2			=>	Vc_UN2_B,
			Vc_UN3			=>	Vc_UN3_B,
			Vc_UN4			=>	Vc_UN4_B,
			Vc_VP1			=>	Vc_VP1_B,
			Vc_VP2			=>	Vc_VP2_B,
			Vc_VP3			=>	Vc_VP3_B,
			Vc_VP4			=>	Vc_VP4_B,
			Vc_VN1			=>	Vc_VN1_B,
			Vc_VN2			=>	Vc_VN2_B,
			Vc_VN3			=>	Vc_VN3_B,
			Vc_VN4			=>	Vc_VN4_B,
			Vc_WP1			=>	Vc_WP1_B,
			Vc_WP2			=>	Vc_WP2_B,
			Vc_WP3			=>	Vc_WP3_B,
			Vc_WP4			=>	Vc_WP4_B,
			Vc_WN1			=>	Vc_WN1_B,
			Vc_WN2			=>	Vc_WN2_B,
			Vc_WN3			=>	Vc_WN3_B,
			Vc_WN4			=>	Vc_WN4_B,
			Iz_UP				=>	Iz_UP_B,			--A[dIz
			Iz_UN				=>	Iz_UN_B,
			Iz_VP				=>	Iz_VP_B,
			Iz_VN				=>	Iz_VN_B,
			Iz_WP				=>	Iz_WP_B,
			Iz_WN				=>	Iz_WN_B,
			Vdc				=>	Vdc_B,			--ddVdc
			Vuv_mmc			=>	Vuv_mmc_B,		--MMCodVo_mmc
			Vwu_mmc			=>	Vwu_mmc_B,
			Vvw_mmc			=>	Vvw_mmc_B,
			Vwu_inv			=>	Vwu_inv_B,		--Co[^odVo_inv
			Vvw_inv			=>	Vvw_inv_B,
			Vuv_inv			=>	Vuv_inv_B,
			IU_inv			=>	IU_inv_B,		--Co[^odIo_inv
			IV_inv			=>	IV_inv_B,
			IW_inv			=>	IW_inv_B		
			-------
			);

	nADC_GR0_CS_OUT		<= not AD_CS when AD_CH_EN(0) = '1' else 'Z';
	nADC_GR1_CS_OUT		<= not AD_CS when AD_CH_EN(1) = '1' else 'Z';
	nADC_GR2_CS_OUT		<= not AD_CS when AD_CH_EN(2) = '1' else 'Z';
	nADC_GR3_CS_OUT		<= not AD_CS when AD_CH_EN(3) = '1' else 'Z';
	nADC_GR4_CS_OUT		<= not AD_CS when AD_CH_EN(4) = '1' else 'Z';
	nADC_GR5_CS_OUT		<= not AD_CS when AD_CH_EN(5) = '1' else 'Z';
	nADC_GR6_CS_OUT		<= not AD_CS when AD_CH_EN(6) = '1' else 'Z';
	nADC_GR7_CS_OUT		<= not AD_CS when AD_CH_EN(7) = '1' else 'Z';
	nADC_GR8_CS_OUT		<= not AD_CS when AD_CH_EN(8) = '1' else 'Z';
	nADC_GR9_CS_OUT		<= not AD_CS when AD_CH_EN(9) = '1' else 'Z';
	ADC_GR0_SCLK_OUT	<= AD_CLK when AD_CH_EN(0) = '1' else 'Z';
	ADC_GR1_SCLK_OUT	<= AD_CLK when AD_CH_EN(1) = '1' else 'Z';
	ADC_GR2_SCLK_OUT	<= AD_CLK when AD_CH_EN(2) = '1' else 'Z';
	ADC_GR3_SCLK_OUT	<= AD_CLK when AD_CH_EN(3) = '1' else 'Z';
	ADC_GR4_SCLK_OUT	<= AD_CLK when AD_CH_EN(4) = '1' else 'Z';
	ADC_GR5_SCLK_OUT	<= AD_CLK when AD_CH_EN(5) = '1' else 'Z';
	ADC_GR6_SCLK_OUT	<= AD_CLK when AD_CH_EN(6) = '1' else 'Z';
	ADC_GR7_SCLK_OUT	<= AD_CLK when AD_CH_EN(7) = '1' else 'Z';
	ADC_GR8_SCLK_OUT	<= AD_CLK when AD_CH_EN(8) = '1' else 'Z';
	ADC_GR9_SCLK_OUT	<= AD_CLK when AD_CH_EN(9) = '1' else 'Z';


-- ====================================================
-- WDT
-- ====================================================

	WDT :MMC_WDT
		port map(
			CLK100		=> CLK100,
			RESET		=> RESET_BB,
			A			=> BUS_A_BB(7 downto 0),
			DW			=> BUS_DW_BB,
			DR			=> BUS_DR_3,
			RD			=> BUS_RD_3,
			WR			=> BUS_WR_3,
			
			WDT_ON		=> WDT_ON,
			ERROR		=> WDT_ERR
			);




-- ====================================================
-- CPU Signals etc.
-- ====================================================
	GB	<= '1' when ADC_ERR = '1' or WDT_ERR = '1' else '0';

	CPU_FPGA_RES(0)				<= nINT_SIG;
	CPU_FPGA_RES(9 downto 1)	<= "111111111";
	FPGA_RES					<= "1111111111";


-- ====================================================
-- DIO
-- ====================================================
-- RgAEg(20190713)
--	DO_OUT					<= DO_IN(15 downto 0);			-- X[

--------------------------------------------------------------Q[gMoIII---------------------------------------------------------------------
-- (20190713):HeadspringDSPQ[gM
--	nMMC_PWM1_OUT <= nMMC_PWM1_buf;
--	nMMC_PWM2_OUT <= nMMC_PWM2_buf;
	
-- ver3.1:FPGAQ[gM
--	nMMC_PWM1_OUT <= nMMC_PWM1_FPGA_B;
--	nMMC_PWM2_OUT <= nMMC_PWM2_FPGA_B;

	nMMC_PWM1_OUT <= nMMC_PWM1_FPGA_buf_buf_buf;
	nMMC_PWM2_OUT <= nMMC_PWM2_FPGA_buf_buf_buf;
---------------------------------------------------------------------------------------------------------------------------------------------------


	DO_OUT					<= nMMC_PWM1_buf(0) & nMMC_PWM1_buf(32) & nMMC_PWM2_buf(16) & BUS_A_IN(15 downto 5) & BUS_A_IN(0) & BUS_RnW_IN ;
--ver3.7.3--DSPQ[gMo---
--	nINV_PWM_OUT			<= nMMC_PWM1_buf(44) & nMMC_PWM1_buf(40) & nMMC_PWM1_buf(36) & nMMC_PWM1_buf(32) & nMMC_PWM1_buf(28) & nMMC_PWM1_buf(24) & nMMC_PWM1_buf(20) & nMMC_PWM1_buf(16) & nMMC_PWM1_buf(12) & nMMC_PWM1_buf(8) & nMMC_PWM1_buf(4) & nMMC_PWM1_buf(0);
-------------------------------------

--ver3.7.3--FPGAQ[gMo---
	nINV_PWM_OUT			<= nMMC_PWM1_FPGA_B(44) & nMMC_PWM1_FPGA_B(40) & nMMC_PWM1_FPGA_B(36) & nMMC_PWM1_FPGA_B(32) & nMMC_PWM1_FPGA_B(28) & nMMC_PWM1_FPGA_B(24) & nMMC_PWM1_FPGA_B(20) & nMMC_PWM1_FPGA_B(16) & nMMC_PWM1_FPGA_B(12) & nMMC_PWM1_FPGA_B(8) & nMMC_PWM1_FPGA_B(4) & nMMC_PWM1_FPGA_B(0);
-------------------------------------

---------------------
	DI_OUT(15 downto 0)		<= DI_LPF;						-- tB^MX[

	nLED_OUT(0)				<= not PWM_ON;
	nLED_OUT(1)				<= not WDT_ON;
	nLED_OUT(2)				<= not WDT_ERR;
	nLED_OUT(3)				<= not ADC_ERR;



-- ====================================================
-- Gate Signal
-- ====================================================
--RgAEg
--vO
-- nINV_PWM_OUT
--	nINV_PWM_OUT	<= not INV_PWM_IN when GB = '0' and RESET = '0' else "111111111111";		-- GB



--ver2--
-- ====================================================
-- control_top
-- ====================================================
--
--	Control :control_top
--		port map(
--			CLK100		=> CLK100,									--100MNbN
--			Vc_UP1		=>	Vc_UP1_B,		--RfTdVc
--			Vc_UP2		=>	Vc_UP2_B,
--			Vc_UP3		=>	Vc_UP3_B,
--			Vc_UP4		=>	Vc_UP4_B,
--			Vc_UN1		=>	Vc_UN1_B,
--			Vc_UN2		=>	Vc_UN2_B,
--			Vc_UN3		=>	Vc_UN3_B,
--			Vc_UN4		=>	Vc_UN4_B,
--			Vc_VP1		=>	Vc_VP1_B,
--			Vc_VP2		=>	Vc_VP2_B,
--			Vc_VP3		=>	Vc_VP3_B,
--			Vc_VP4		=>	Vc_VP4_B,
--			Vc_VN1		=>	Vc_VN1_B,
--			Vc_VN2		=>	Vc_VN2_B,
--			Vc_VN3		=>	Vc_VN3_B,
--			Vc_VN4		=>	Vc_VN4_B,
--			Vc_WP1		=>	Vc_WP1_B,
--			Vc_WP2		=>	Vc_WP2_B,
--			Vc_WP3		=>	Vc_WP3_B,
--			Vc_WP4		=>	Vc_WP4_B,
--			Vc_WN1		=>	Vc_WN1_B,
--			Vc_WN2		=>	Vc_WN2_B,
--			Vc_WN3		=>	Vc_WN3_B,
--			Vc_WN4		=>	Vc_WN4_B,
--			Iz_UP			=>	Iz_UP_B,			--A[dIz
--			Iz_UN			=>	Iz_UN_B,
--			Iz_VP			=>	Iz_VP_B,
--			Iz_VN			=>	Iz_VN_B,
--			Iz_WP			=>	Iz_WP_B,
--			Iz_WN			=>	Iz_WN_B,
--			Vdc			=>	Vdc_B,			--ddVdc
--			Vuv_mmc		=>	Vuv_mmc_B,		--MMCodVo_mmc
--			Vwu_mmc		=>	Vwu_mmc_B,
--			Vvw_mmc		=>	Vvw_mmc_B,
--			Vwu_inv		=>	Vwu_inv_B,		--Co[^odVo_inv
--			Vvw_inv		=>	Vvw_inv_B,
--			Vuv_inv		=>	Vuv_inv_B,
--			IU_inv		=>	IU_inv_B,		--Co[^odIo_inv
--			IV_inv		=>	IV_inv_B,
--			IW_inv		=>	IW_inv_B,
--			MRate_UP1	=>	MRate_UP1_B,		--UT
--			MRate_UP2	=>	MRate_UP2_B,
--			MRate_UP3	=>	MRate_UP3_B,
--			MRate_UP4	=>	MRate_UP4_B,
--			MRate_UN1	=>	MRate_UN1_B,
--			MRate_UN2	=>	MRate_UN2_B,
--			MRate_UN3	=>	MRate_UN3_B,
--			MRate_UN4	=>	MRate_UN4_B,
--			MRate_VP1	=>	MRate_VP1_B,		--VT
--			MRate_VP2	=>	MRate_VP2_B,
--			MRate_VP3	=>	MRate_VP3_B,
--			MRate_VP4	=>	MRate_VP4_B,
--			MRate_VN1	=>	MRate_VN1_B,
--			MRate_VN2	=>	MRate_VN2_B,
--			MRate_VN3	=>	MRate_VN3_B,
--			MRate_VN4	=>	MRate_VN4_B,
--			MRate_WP1	=>	MRate_WP1_B,		--WT
--			MRate_WP2	=>	MRate_WP2_B,
--			MRate_WP3	=>	MRate_WP3_B,
--			MRate_WP4	=>	MRate_WP4_B,
--			MRate_WN1	=>	MRate_WN1_B,
--			MRate_WN2	=>	MRate_WN2_B,
--			MRate_WN3	=>	MRate_WN3_B,
--			MRate_WN4	=>	MRate_WN4_B,
--			------
--			--ver3.4.2--chipscopefobOp---------------------
--			VAu_out 		=>	VAu_B,
--			VAv_out 		=>	VAv_B,
--			vb_up1_out  =>	vb_up1_B,
--			vb_up2_out  =>	vb_up2_B,
--			vb_up3_out  =>	vb_up3_B,
--			vb_up4_out  =>	vb_up4_B,
--			vb_un1_out  =>	vb_un1_B,
--			vb_un2_out  =>	vb_un2_B,
--			vb_un3_out  =>	vb_un3_B,
--			vb_un4_out  =>	vb_un4_B,
--			vb_vp1_out  =>	vb_vp1_B,
--			vb_vp2_out  =>	vb_vp2_B,
--			vb_vp3_out  =>	vb_vp3_B,
--			vb_vp4_out  =>	vb_vp4_B,
--			vb_vn1_out  =>	vb_vn1_B,
--			vb_vn2_out  =>	vb_vn2_B,
--			vb_vn3_out  =>	vb_vn3_B,
--			vb_vn4_out  =>	vb_vn4_B,
--			-------------------------------------------------
--			--ver3.7--lp
--			GATE_EN => GATE_EN_B,
--			----------
--			--ver3.7.1--
--			VA_no_limit_U  =>	VA_no_limit_U_B,
--			VA_no_limit_V  =>	VA_no_limit_V_B,
--			----------
--			--ver3.7.5--
--			Voref_out  =>	Voref_B,
--			------------
--			--ver3.8.5(debug)--
--			Iz_UP_N_debug_out => Iz_UP_N_debug_out_B,
--			Iz_UP_N_out => Iz_UP_N_out_B,
--			-------------------
--			
--			----ver1.1(QMDB)----			
--			dVoref_out => dVoref_out_B,
--			IC_N_out   => IC_N_out_B,
--			Vo_N_out   => Vo_N_out_B,
--			QMDB_OUTPUT_out => QMDB_OUTPUT_out_B,
--			Vdc_N_out  => Vdc_N_out_B,
--			soft_cnt_out	 => soft_cnt_out_B,
--			clk10k_B_out => clk10k_BB,
--			cnt80k_debug => cnt80k_debug_B
--			-------------------
--			);
--
-------------


--ver2.3--
-- ====================================================
-- MMC_PWM_FPGA
-- ====================================================

	MMC_PWM_FPGA_GEN : MMC_PWM_FPGA
		port map(
			CLK100		=> CLK100,
			RESET		=> RESET_BB,
			A			=> BUS_A_BB(7 downto 0),
			DW			=> BUS_DW_BB,
--			DR			=> BUS_DR_1,
--			DR2			=> BUS_DR_1_2,
--			DR3			=> BUS_DR_1_3,
			RD			=> BUS_RD_1,
			WR			=> BUS_WR_1,
			
			GB       => GB,
			
			MRate_UP1 => X"0000",
			MRate_UP2 => X"0000",
			MRate_UP3 => X"0000",
			MRate_UP4 => X"0000",
			MRate_UN1 => X"0000",
			MRate_UN2 => X"0000",
			MRate_UN3 => X"0000",
			MRate_UN4 => X"0000",
			MRate_VP1 => X"0000",
			MRate_VP2 => X"0000",
			MRate_VP3 => X"0000",
			MRate_VP4 => X"0000",
			MRate_VN1 => X"0000",
			MRate_VN2 => X"0000",
			MRate_VN3 => X"0000",
			MRate_VN4 => X"0000",
		
			--ver2.6.4--
			REF_UP1_vector => X"0000",
			REF_UP2_vector => X"0000",
			REF_UP3_vector => X"0000",
			REF_UP4_vector	=> X"0000",
			REF_UN1_vector => X"0000",
			REF_UN2_vector	=> X"0000",
			REF_UN3_vector	=> X"0000",
			REF_UN4_vector	=> X"0000",
			REF_VP1_vector => X"0000",
			REF_VP2_vector	=> X"0000",
			REF_VP3_vector	=> X"0000",
			REF_VP4_vector	=> X"0000",
			REF_VN1_vector	=> X"0000",
			REF_VN2_vector	=> X"0000",
			REF_VN3_vector	=> X"0000",
			REF_VN4_vector	=> X"0000",
			----------------
			--ver3.4.2--chipscopefobOp---------------------
			VAu 	  =>	X"0000",
			VAv 	  =>	X"0000",
			vb_up1  =>	X"0000",
			vb_up2  =>	X"0000",
			vb_up3  =>	X"0000",
			vb_up4  =>	X"0000",
			vb_un1  =>	X"0000",
			vb_un2  =>	X"0000",
			vb_un3  =>	X"0000",
			vb_un4  =>	X"0000",
			vb_vp1  =>	X"0000",
			vb_vp2  =>	X"0000",
			vb_vp3  =>	X"0000",
			vb_vp4  =>	X"0000",
			vb_vn1  =>	X"0000",
			vb_vn2  =>	X"0000",
			vb_vn3  =>	X"0000",
			vb_vn4  =>	X"0000",

			nMMC_PWM1_FPGA	=> nMMC_PWM1_FPGA_B,
			nMMC_PWM2_FPGA	=> nMMC_PWM2_FPGA_B,
			nINT_SIG_FPGA	=> nINT_SIG_FPGA_B,
			
			PWM_ON_FPGA		=> PWM_ON_FPGA_B,
			GB_FPGA			=> GB_FPGA_B,

			--ver3.7--lp
--			GATE_EN_out => GATE_EN_B,
			----------
			--ver3.7.1--
			VA_no_limit_U  =>	X"0000",
			VA_no_limit_V  =>	X"0000",
			----------
			--ver3.7.5--
			Voref_BB  =>	X"0000",
			------------
			--ver3.8.5(debug)--
			Iz_UP_N_debug_out_BB => X"0000",
			Iz_UP_N_out_BB => X"0000",
			-------------------
			----ver1.1(QMDB)----			
			dVoref_out_BB      => X"0000",
			IC_N_out_BB        => X"0000",
			Vo_N_out_BB        => X"0000",
			QMDB_OUTPUT_out_BB => X"0000",
			Vdc_N_out_BB       => X"0000",
			soft_cnt_out_BB	 => X"0000",
			clk10k_BBB => '1',
			cnt80k_debug_BB =>X"0000"
			);

--	SVPWM_Control--20220303
		SVPWM_Control :control_top2
		port map(
			CLK		=> CLK100,									--100MNbN
			reset 		=> RESET,
			clk_enable => '1',
			V_uv		=>	Vuv_mmc_B,
         V_uw		=>	Vvw_mmc_B,
         iout_u	=>	Vuv_inv_B,
         iout_v 	=>	Vvw_inv_B,
         iL_u 	   =>	IU_inv_B,
         iL_v     =>	IV_inv_B,
			switch_control		=>	Vwu_inv_B,
			Up_current			=>	Iz_UP_B,			--A[dIz
         Un_current			=>	Iz_UN_B,
         Vp_current 			=>	Iz_VP_B,
         Vn_current			=>	Iz_VN_B,
         Wp_current			=>	Iz_WP_B,
         UW_current			=>	Iz_WN_B,
			CV_U1		=>	Vc_UP1_B,		--RfTdVc
			CV_U2		=>	Vc_UP2_B,
			CV_U3		=>	Vc_UP3_B,
			CV_U4		=>	Vc_UP4_B,
			CV_U5		=>	Vc_UN1_B,
			CV_U6		=>	Vc_UN2_B,
			CV_U7		=>	Vc_UN3_B,
			CV_U8		=>	Vc_UN4_B,
			CV_V1		=>	Vc_VP1_B,
			CV_V2		=>	Vc_VP2_B,
			CV_V3		=>	Vc_VP3_B,
			CV_V4		=>	Vc_VP4_B,
			CV_V5		=>	Vc_VN1_B,
			CV_V6		=>	Vc_VN2_B,
			CV_V7		=>	Vc_VN3_B,
			CV_V8		=>	Vc_VN4_B,
			CV_W1		=>	Vc_WP1_B,
			CV_W2		=>	Vc_WP2_B,
			CV_W3		=>	Vc_WP3_B,
			CV_W4		=>	Vc_WP4_B,
			CV_W5		=>	Vc_WN1_B,
			CV_W6		=>	Vc_WN2_B,
			CV_W7		=>	Vc_WN3_B,
			CV_W8		=>	Vc_WN4_B,
			----ce_out_0			=>	,	
			----ce_out_1			=>	,
			MMC_PWM_OUT1	=> nMMC_PWM1_FPGA_buf_buf_buf,
			MMC_PWM_OUT2	=> nMMC_PWM2_FPGA_buf_buf_buf,
			
			Vuv					=> Vuv_ILA,
			Vvw					=> Vvw_ILA,
			Ioutu					=> Ioutu_ILA,
			Ioutv					=> Ioutv_ILA,
			ILu					=> ILu_ILA,
			ILv					=> ILv_ILA,
			Vrefa					=> Vrefa_ILA,
			Vrefb					=> Vrefb_ILA,
			deltaTa				=> deltaTa_ILA,
			deltaTb				=> deltaTb_ILA,
			deltaTa_hold		=> deltaTa_hold_ILA,
			deltaTb_hold		=> deltaTb_hold_ILA,
			switch_control1	=> switche_contol_ILA
			);

-----------

--	count_1MHz--
cnt_1MHz0 :count_1M
		port map(
			CLK		=> CLK100,									--100MNbN
			cnt_1M 	=> cnt_1MHz
	);

--	count_50kHz--
cnt_50kHz0 :count_50k
		port map(
			CLK		=> cnt_1MHz,
			cnt_50k 	=> cnt_50kHz
	);
	
---icon--------
icon_srdcdb0 : icon_srdcdb
	port map(
			CONTROL0 => CONTROL0
	);

---ila---------
ila_srdcdb0 : ila_srdcdb
	port map(
			CONTROL => CONTROL0,
			CLK => cnt_50kHz,
			TRIG0 => Vuv_ILA,
			TRIG1 => Vvw_ILA,
			TRIG2 => Ioutu_ILA,
			TRIG3 => Ioutv_ILA,
			TRIG4 => ILu_ILA,
			TRIG5 => ILv_ILA,
			TRIG6 => Vrefa_ILA,
			TRIG7 => Vrefb_ILA,
			TRIG8 => deltaTa_ILA,
			TRIG9 => deltaTb_ILA,
			TRIG10 => deltaTa_hold_ILA,
			TRIG11 => deltaTb_hold_ILA,
			TRIG12 => switche_contol_ILA
	);
---------------


-- ====================================================
-- Dummy Signal for unused pins
-- ====================================================
	process(CLK100)
	begin
		if CLK100'event and CLK100 = '1' then
			DUMMY	<= BUS_CLK_IN & BUS_CS3_BB & BUS_RD_BB & BUS_WR_BB & BUS_A_BB(16);
		end if;
	end process;
	
	nBUS_WAIT_OUT	<= '1';


end Behavioral;



\end{mylisting}
%%Root main.tex
%---------------------------------------------------------------------------
%付録A　付録
%---------------------------------------------------------------------------
%\lstnewenvironment{mylisting}[1][]
%    {\lstset{
%        frame=single,
%        numbersep=10pt,
%        tabsize=2,
%        #1
%    }
%}{}
\begin{mylisting}[language=VHDL,caption=ゲート信号出力モジュール,label=MMC_PWM_GEN_module_VHDL]


----------------------------------------------------------------------------------
-- Company: 
-- Engineer: 
-- 
-- Create Date:    22:06:24 08/08/2019 
-- Design Name: 
-- Module Name:    MMC_PWM_FPGA - Behavioral 
-- Project Name: 
-- Target Devices: 
-- Tool versions: 
-- Description: 
--
-- Dependencies: 
--
-- Revision: 
-- Revision 0.01 - File Created
-- Additional Comments: 
--
----------------------------------------------------------------------------------
library IEEE;
use IEEE.STD_LOGIC_1164.ALL;
use IEEE.STD_LOGIC_arith.all;
use IEEE.STD_LOGIC_SIGNED.ALL;

-- Uncomment the following library declaration if using
-- arithmetic functions with Signed or Unsigned values
--use IEEE.NUMERIC_STD.ALL;

-- Uncomment the following library declaration if instantiating
-- any Xilinx primitives in this code.
--library UNISIM;
--use UNISIM.VComponents.all;

entity MMC_PWM_FPGA is

	generic (

--ver2.5.4--
		--三角波刻み10kHz
--		mrate_max				: integer := 40960000;	
	--	tri_cnt_max				: integer := 2500;		
	--	tri_cnt_max_half		: integer := 1250
------------				--三角波刻み5kHz
		mrate_max				: integer := 20480000;	
		tri_cnt_max				: integer := 5000;		
		tri_cnt_max_half		: integer := 2500
------------		


	);
	
	Port (
		CLK100		: in	std_logic;
		RESET		: in	std_logic;
		A			: in	std_logic_vector(7 downto 0);
		DW			: in	std_logic_vector(15 downto 0);
--		DR			: out	std_logic_vector(15 downto 0)	:= "00000000" & "00000000";
--		DR2		: out	std_logic_vector(15 downto 0)	:= "00000000" & "00000000";
--		DR3		: out	std_logic_vector(15 downto 0)	:= "00000000" & "00000000";
		RD			: in	std_logic;
		WR			: in	std_logic;
		
		GB			: in  std_logic;
		
		MRate_UP1 : in std_logic_vector(15 downto 0);
		MRate_UP2 : in std_logic_vector(15 downto 0);
		MRate_UP3 : in std_logic_vector(15 downto 0);
		MRate_UP4 : in std_logic_vector(15 downto 0);
		MRate_UN1 : in std_logic_vector(15 downto 0);
		MRate_UN2 : in std_logic_vector(15 downto 0);
		MRate_UN3 : in std_logic_vector(15 downto 0);
		MRate_UN4 : in std_logic_vector(15 downto 0);
		MRate_VP1 : in std_logic_vector(15 downto 0);
		MRate_VP2 : in std_logic_vector(15 downto 0);
		MRate_VP3 : in std_logic_vector(15 downto 0);
		MRate_VP4 : in std_logic_vector(15 downto 0);
		MRate_VN1 : in std_logic_vector(15 downto 0);
		MRate_VN2 : in std_logic_vector(15 downto 0);
		MRate_VN3 : in std_logic_vector(15 downto 0);
		MRate_VN4 : in std_logic_vector(15 downto 0);
		
		--ver2.6.4追加--
		REF_UP1_vector	: in std_logic_vector(15 downto 0);
		REF_UP2_vector	: in std_logic_vector(15 downto 0);
		REF_UP3_vector	: in std_logic_vector(15 downto 0);
		REF_UP4_vector	: in std_logic_vector(15 downto 0);
		REF_UN1_vector	: in std_logic_vector(15 downto 0);
		REF_UN2_vector	: in std_logic_vector(15 downto 0);
		REF_UN3_vector	: in std_logic_vector(15 downto 0);
		REF_UN4_vector	: in std_logic_vector(15 downto 0);
		REF_VP1_vector	: in std_logic_vector(15 downto 0);
		REF_VP2_vector	: in std_logic_vector(15 downto 0);
		REF_VP3_vector	: in std_logic_vector(15 downto 0);
		REF_VP4_vector	: in std_logic_vector(15 downto 0);
		REF_VN1_vector	: in std_logic_vector(15 downto 0);
		REF_VN2_vector	: in std_logic_vector(15 downto 0);
		REF_VN3_vector	: in std_logic_vector(15 downto 0);
		REF_VN4_vector	: in std_logic_vector(15 downto 0);
		----------------
		
		--ver3.4.2--chipscopeデバッグ用---------------------
		VAu 	  : in std_logic_vector(15 downto 0);
		VAv 	  : in std_logic_vector(15 downto 0);
		vb_up1  : in std_logic_vector(15 downto 0);
		vb_up2  : in std_logic_vector(15 downto 0);
		vb_up3  : in std_logic_vector(15 downto 0);
		vb_up4  : in std_logic_vector(15 downto 0);
		vb_un1  : in std_logic_vector(15 downto 0);
		vb_un2  : in std_logic_vector(15 downto 0);
		vb_un3  : in std_logic_vector(15 downto 0);
		vb_un4  : in std_logic_vector(15 downto 0);
		vb_vp1  : in std_logic_vector(15 downto 0);
		vb_vp2  : in std_logic_vector(15 downto 0);
		vb_vp3  : in std_logic_vector(15 downto 0);
		vb_vp4  : in std_logic_vector(15 downto 0);
		vb_vn1  : in std_logic_vector(15 downto 0);
		vb_vn2  : in std_logic_vector(15 downto 0);
		vb_vn3  : in std_logic_vector(15 downto 0);
		vb_vn4  : in std_logic_vector(15 downto 0);
		-------------------------------------------------
		
		nMMC_PWM1_FPGA	: out	std_logic_vector(47 downto 0);
		nMMC_PWM2_FPGA	: out	std_logic_vector(47 downto 0);
		nINT_SIG_FPGA	: out	std_logic := '1';
		
		PWM_ON_FPGA		: out	std_logic := '0';
		GB_FPGA			: in	std_logic;
		
		--ver3.7--平均値制御の偏差積分初期化用
		GATE_EN_out : out std_logic;
		----------
		--ver3.7.1--
		VA_no_limit_U : in std_logic_vector(15 downto 0);
		VA_no_limit_V : in std_logic_vector(15 downto 0);
		----------
		--ver3.7.5--
		Voref_BB	: in std_logic_vector(15 downto 0);
		------------
		--ver3.8.5(debug)--
		Iz_UP_N_debug_out_BB : in std_logic_vector(15 downto 0);
		Iz_UP_N_out_BB 		: in std_logic_vector(15 downto 0);
		-------------------
--QMDB_ver1.1--------------------------------------------------------------------------------------
		dVoref_out_BB        : IN     std_logic_vector(15 downto 0);
		IC_N_out_BB	        : IN   	std_logic_vector(15 DOWNTO 0);
		Vo_N_out_BB	        : IN 	   std_logic_vector(15 DOWNTO 0);
		QMDB_OUTPUT_out_BB   : IN     std_logic_vector(15 DOWNTO 0);
		Vdc_N_out_BB		    : IN     std_logic_vector(15 DOWNTO 0);
		soft_cnt_out_BB	      : IN     std_logic_vector(15 DOWNTO 0);
		clk10k_BBB           : IN     std_logic;
		cnt80k_debug_BB		: IN 		std_logic_vector(15 DOWNTO 0)
---------------------------------------------------------------------------------------------------
		);
end MMC_PWM_FPGA;

architecture Behavioral of MMC_PWM_FPGA is

--	constant Carrier_4k_count_max			: integer range 0 to 25000						:= 25000;--CLK100M
	
--	constant Carrier_10k_count_max		: integer range 0 to 10000						:= 10000;--CLK100M
--	constant Carrier_10k_count_max		: integer range 0 to 200						:= 200;	--CLK2M
--	constant Carrier_10k_count_max		: integer range 0 to 40							:= 40;	--CLK400k

	signal 	MRate_UP1_C						: std_logic_vector(31 downto 0):=X"00000000";
	signal 	MRate_UP2_C						: std_logic_vector(31 downto 0):=X"00000000";
	signal 	MRate_UP3_C						: std_logic_vector(31 downto 0):=X"00000000";
	signal 	MRate_UP4_C						: std_logic_vector(31 downto 0):=X"00000000";
	signal 	MRate_UN1_C						: std_logic_vector(31 downto 0):=X"00000000";
	signal 	MRate_UN2_C						: std_logic_vector(31 downto 0):=X"00000000";
	signal 	MRate_UN3_C						: std_logic_vector(31 downto 0):=X"00000000";
	signal 	MRate_UN4_C						: std_logic_vector(31 downto 0):=X"00000000";
	signal 	MRate_VP1_C						: std_logic_vector(31 downto 0):=X"00000000";
	signal 	MRate_VP2_C						: std_logic_vector(31 downto 0):=X"00000000";
	signal 	MRate_VP3_C						: std_logic_vector(31 downto 0):=X"00000000";
	signal 	MRate_VP4_C						: std_logic_vector(31 downto 0):=X"00000000";
	signal 	MRate_VN1_C						: std_logic_vector(31 downto 0):=X"00000000";
	signal 	MRate_VN2_C						: std_logic_vector(31 downto 0):=X"00000000";
	signal 	MRate_VN3_C						: std_logic_vector(31 downto 0):=X"00000000";
	signal 	MRate_VN4_C						: std_logic_vector(31 downto 0):=X"00000000";
	
	signal 	MRate_UP1_C_2					: std_logic_vector(15 downto 0):=X"0000";
	signal 	MRate_UP2_C_2					: std_logic_vector(15 downto 0):=X"0000";
	signal 	MRate_UP3_C_2					: std_logic_vector(15 downto 0):=X"0000";
	signal 	MRate_UP4_C_2					: std_logic_vector(15 downto 0):=X"0000";
	signal 	MRate_UN1_C_2					: std_logic_vector(15 downto 0):=X"0000";
	signal 	MRate_UN2_C_2					: std_logic_vector(15 downto 0):=X"0000";
	signal 	MRate_UN3_C_2					: std_logic_vector(15 downto 0):=X"0000";
	signal 	MRate_UN4_C_2					: std_logic_vector(15 downto 0):=X"0000";
	signal 	MRate_VP1_C_2					: std_logic_vector(15 downto 0):=X"0000";
	signal 	MRate_VP2_C_2					: std_logic_vector(15 downto 0):=X"0000";
	signal 	MRate_VP3_C_2					: std_logic_vector(15 downto 0):=X"0000";
	signal 	MRate_VP4_C_2					: std_logic_vector(15 downto 0):=X"0000";
	signal 	MRate_VN1_C_2					: std_logic_vector(15 downto 0):=X"0000";
	signal 	MRate_VN2_C_2					: std_logic_vector(15 downto 0):=X"0000";
	signal 	MRate_VN3_C_2					: std_logic_vector(15 downto 0):=X"0000";
	signal 	MRate_VN4_C_2					: std_logic_vector(15 downto 0):=X"0000";

	signal 	MRate_UP1_C_3					: std_logic_vector(15 downto 0):=X"0000";
	signal 	MRate_UP2_C_3					: std_logic_vector(15 downto 0):=X"0000";
	signal 	MRate_UP3_C_3					: std_logic_vector(15 downto 0):=X"0000";
	signal 	MRate_UP4_C_3					: std_logic_vector(15 downto 0):=X"0000";
	signal 	MRate_UN1_C_3					: std_logic_vector(15 downto 0):=X"0000";
	signal 	MRate_UN2_C_3					: std_logic_vector(15 downto 0):=X"0000";
	signal 	MRate_UN3_C_3					: std_logic_vector(15 downto 0):=X"0000";
	signal 	MRate_UN4_C_3					: std_logic_vector(15 downto 0):=X"0000";
	signal 	MRate_VP1_C_3					: std_logic_vector(15 downto 0):=X"0000";
	signal 	MRate_VP2_C_3					: std_logic_vector(15 downto 0):=X"0000";
	signal 	MRate_VP3_C_3					: std_logic_vector(15 downto 0):=X"0000";
	signal 	MRate_VP4_C_3					: std_logic_vector(15 downto 0):=X"0000";
	signal 	MRate_VN1_C_3					: std_logic_vector(15 downto 0):=X"0000";
	signal 	MRate_VN2_C_3					: std_logic_vector(15 downto 0):=X"0000";
	signal 	MRate_VN3_C_3					: std_logic_vector(15 downto 0):=X"0000";
	signal 	MRate_VN4_C_3					: std_logic_vector(15 downto 0):=X"0000";

	signal 	MRate_UP1_C_4					: std_logic_vector(15 downto 0):=X"0000";
	signal 	MRate_UP2_C_4					: std_logic_vector(15 downto 0):=X"0000";
	signal 	MRate_UP3_C_4					: std_logic_vector(15 downto 0):=X"0000";
	signal 	MRate_UP4_C_4					: std_logic_vector(15 downto 0):=X"0000";
	signal 	MRate_UN1_C_4					: std_logic_vector(15 downto 0):=X"0000";
	signal 	MRate_UN2_C_4					: std_logic_vector(15 downto 0):=X"0000";
	signal 	MRate_UN3_C_4					: std_logic_vector(15 downto 0):=X"0000";
	signal 	MRate_UN4_C_4					: std_logic_vector(15 downto 0):=X"0000";
	signal 	MRate_VP1_C_4					: std_logic_vector(15 downto 0):=X"0000";
	signal 	MRate_VP2_C_4					: std_logic_vector(15 downto 0):=X"0000";
	signal 	MRate_VP3_C_4					: std_logic_vector(15 downto 0):=X"0000";
	signal 	MRate_VP4_C_4					: std_logic_vector(15 downto 0):=X"0000";
	signal 	MRate_VN1_C_4					: std_logic_vector(15 downto 0):=X"0000";
	signal 	MRate_VN2_C_4					: std_logic_vector(15 downto 0):=X"0000";
	signal 	MRate_VN3_C_4					: std_logic_vector(15 downto 0):=X"0000";
	signal 	MRate_VN4_C_4					: std_logic_vector(15 downto 0):=X"0000";

	signal 	MRate_UP1_C_4_hold					: std_logic_vector(15 downto 0):=X"0000";
	signal 	MRate_UP2_C_4_hold					: std_logic_vector(15 downto 0):=X"0000";
	signal 	MRate_UP3_C_4_hold					: std_logic_vector(15 downto 0):=X"0000";
	signal 	MRate_UP4_C_4_hold					: std_logic_vector(15 downto 0):=X"0000";
	signal 	MRate_UN1_C_4_hold					: std_logic_vector(15 downto 0):=X"0000";
	signal 	MRate_UN2_C_4_hold					: std_logic_vector(15 downto 0):=X"0000";
	signal 	MRate_UN3_C_4_hold					: std_logic_vector(15 downto 0):=X"0000";
	signal 	MRate_UN4_C_4_hold					: std_logic_vector(15 downto 0):=X"0000";
	signal 	MRate_VP1_C_4_hold					: std_logic_vector(15 downto 0):=X"0000";
	signal 	MRate_VP2_C_4_hold					: std_logic_vector(15 downto 0):=X"0000";
	signal 	MRate_VP3_C_4_hold					: std_logic_vector(15 downto 0):=X"0000";
	signal 	MRate_VP4_C_4_hold					: std_logic_vector(15 downto 0):=X"0000";
	signal 	MRate_VN1_C_4_hold					: std_logic_vector(15 downto 0):=X"0000";
	signal 	MRate_VN2_C_4_hold					: std_logic_vector(15 downto 0):=X"0000";
	signal 	MRate_VN3_C_4_hold					: std_logic_vector(15 downto 0):=X"0000";
	signal 	MRate_VN4_C_4_hold					: std_logic_vector(15 downto 0):=X"0000";

	signal 	CNT_out_UP1						: integer range -32767 to 32767;
	signal 	CNT_out_UP2						: integer range -32767 to 32767;
	signal 	CNT_out_UP3						: integer range -32767 to 32767;
	signal 	CNT_out_UP4						: integer range -32767 to 32767;
	signal 	CNT_out_UN1						: integer range -32767 to 32767;
	signal 	CNT_out_UN2						: integer range -32767 to 32767;
	signal 	CNT_out_UN3						: integer range -32767 to 32767;
	signal 	CNT_out_UN4						: integer range -32767 to 32767;
	signal 	CNT_out_VP1						: integer range -32767 to 32767;
	signal 	CNT_out_VP2						: integer range -32767 to 32767;
	signal 	CNT_out_VP3						: integer range -32767 to 32767;
	signal 	CNT_out_VP4						: integer range -32767 to 32767;
	signal 	CNT_out_VN1						: integer range -32767 to 32767;
	signal 	CNT_out_VN2						: integer range -32767 to 32767;
	signal 	CNT_out_VN3						: integer range -32767 to 32767;
	signal 	CNT_out_VN4						: integer range -32767 to 32767;

 	signal 	CNT_out_UP1_vector			: std_logic_vector(15 downto 0);
 	signal 	CNT_out_UP2_vector			: std_logic_vector(15 downto 0);
 	signal 	CNT_out_UP3_vector			: std_logic_vector(15 downto 0);
 	signal 	CNT_out_UP4_vector			: std_logic_vector(15 downto 0);
 	signal 	CNT_out_UN1_vector			: std_logic_vector(15 downto 0);
 	signal 	CNT_out_UN2_vector			: std_logic_vector(15 downto 0);
 	signal 	CNT_out_UN3_vector			: std_logic_vector(15 downto 0);
 	signal 	CNT_out_UN4_vector			: std_logic_vector(15 downto 0);
 	signal 	CNT_out_VP1_vector			: std_logic_vector(15 downto 0);
 	signal 	CNT_out_VP2_vector			: std_logic_vector(15 downto 0);
 	signal 	CNT_out_VP3_vector			: std_logic_vector(15 downto 0);
 	signal 	CNT_out_VP4_vector			: std_logic_vector(15 downto 0);
 	signal 	CNT_out_VN1_vector			: std_logic_vector(15 downto 0);
 	signal 	CNT_out_VN2_vector			: std_logic_vector(15 downto 0);
 	signal 	CNT_out_VN3_vector			: std_logic_vector(15 downto 0);
 	signal 	CNT_out_VN4_vector			: std_logic_vector(15 downto 0);
	


	--ver7.2(分解能向上)----------------------------------------------------------------------------------
	constant pwm_assist						: std_logic_vector(15 downto 0) := X"1388"; --5000
	constant pwm_assist_half				: std_logic_vector(15 downto 0) := X"09C4"; --2500	
--	constant pwm_assist						: std_logic_vector(15 downto 0) := "0000100111000100"; --2500
--	constant pwm_assist_half				: std_logic_vector(15 downto 0) := "0000010011100010"; --1250

	--追加ver2.6--
	--https://note.cman.jp/convert/bit/
	--04E2 +1250 -- +100%
	--04A4 +1188 -- +95%
	--FB5C -1188 -- -95%
	--FB1E -1250 -- -100%
	
	--09C4 +2500 -- +100%
	--0947 +2375 -- +95%
	--F6B9 -2375 -- -95%
	--F63C -2500 -- -100%
	
	signal	upper_limit : std_logic_vector(15 downto 0):=X"0947";
	signal	lower_limit : std_logic_vector(15 downto 0):=X"F6B9";

   -----------------------------------------------------------------------------------------------------

	constant count_max_default				: integer range 0 to 65535			:= 65535; 

--　0~5000の場合
--	signal 	cnt_phase_0						: integer range 0 to tri_cnt_max		:= tri_cnt_max_half;	--三角波初期カウント値
--	signal 	cnt_phase_90					: integer range 0 to tri_cnt_max		:= 0;						--三角波初期カウント値		
--	signal 	cnt_phase_180					: integer range 0 to tri_cnt_max		:= tri_cnt_max_half;	--三角波初期カウント値		
--	signal 	cnt_phase_270					: integer range 0 to tri_cnt_max		:= tri_cnt_max;		--三角波初期カウント値	

--ver7.3(位相ずれで9level実現)----------------------------------------------------------------------------------------------------------
--45度位相ずれの場合------------------------
	signal 	cnt_phase_0						: integer range -tri_cnt_max to tri_cnt_max		:= 0;					   --三角波初期カウント値
	signal 	cnt_phase_45					: integer range -tri_cnt_max to tri_cnt_max		:= -tri_cnt_max_half;	--三角波初期カウント値
	signal 	cnt_phase_90					: integer range -tri_cnt_max to tri_cnt_max		:= -tri_cnt_max;	   --三角波初期カウント値		
	signal 	cnt_phase_135					: integer range -tri_cnt_max to tri_cnt_max		:= -tri_cnt_max_half;	--三角波初期カウント値		
	signal 	cnt_phase_180					: integer range -tri_cnt_max to tri_cnt_max		:= 0;					   --三角波初期カウント値		
	signal 	cnt_phase_225					: integer range -tri_cnt_max to tri_cnt_max		:= tri_cnt_max_half;	--三角波初期カウント値		
	signal 	cnt_phase_270					: integer range -tri_cnt_max to tri_cnt_max		:= tri_cnt_max;	   --三角波初期カウント値	
	signal 	cnt_phase_315					: integer range -tri_cnt_max to tri_cnt_max		:= tri_cnt_max_half;	--三角波初期カウント値
	
	signal 	cnt_phase_UD_0					: std_logic			:= '0';		--三角波UP/DOWNフラグ
	signal 	cnt_phase_UD_45				: std_logic			:= '0';		--三角波UP/DOWNフラグ
	signal 	cnt_phase_UD_90				: std_logic			:= '0';		--三角波UP/DOWNフラグ		
	signal 	cnt_phase_UD_135				: std_logic			:= '1';		--三角波UP/DOWNフラグ		
	signal 	cnt_phase_UD_180				: std_logic			:= '1';		--三角波UP/DOWNフラグ		
	signal 	cnt_phase_UD_225				: std_logic			:= '1';		--三角波UP/DOWNフラグ		
	signal 	cnt_phase_UD_270				: std_logic			:= '1';		--三角波UP/DOWNフラグ		
	signal 	cnt_phase_UD_315				: std_logic			:= '0';		--三角波UP/DOWNフラグ		
----------------------------------------
--90度位相ずれの場合------------------------
--	signal 	cnt_phase_0						: integer range -tri_cnt_max to tri_cnt_max		:= 0;					   --三角波初期カウント値
--	signal 	cnt_phase_45					: integer range -tri_cnt_max to tri_cnt_max		:= 0;						--三角波初期カウント値
--	signal 	cnt_phase_90					: integer range -tri_cnt_max to tri_cnt_max		:= -tri_cnt_max;	   --三角波初期カウント値		
--	signal 	cnt_phase_135					: integer range -tri_cnt_max to tri_cnt_max		:= -tri_cnt_max;		--三角波初期カウント値		
--	signal 	cnt_phase_180					: integer range -tri_cnt_max to tri_cnt_max		:= 0;					   --三角波初期カウント値		
--	signal 	cnt_phase_225					: integer range -tri_cnt_max to tri_cnt_max		:= 0;						--三角波初期カウント値		
--	signal 	cnt_phase_270					: integer range -tri_cnt_max to tri_cnt_max		:= tri_cnt_max;	   --三角波初期カウント値	
--	signal 	cnt_phase_315					: integer range -tri_cnt_max to tri_cnt_max		:= tri_cnt_max;		--三角波初期カウント値
--
--	signal 	cnt_phase_UD_0					: std_logic			:= '0';		--三角波UP/DOWNフラグ
--	signal 	cnt_phase_UD_45				: std_logic			:= '0';		--三角波UP/DOWNフラグ
--	signal 	cnt_phase_UD_90				: std_logic			:= '0';		--三角波UP/DOWNフラグ		
--	signal 	cnt_phase_UD_135				: std_logic			:= '0';		--三角波UP/DOWNフラグ		
--	signal 	cnt_phase_UD_180				: std_logic			:= '1';		--三角波UP/DOWNフラグ		
--	signal 	cnt_phase_UD_225				: std_logic			:= '1';		--三角波UP/DOWNフラグ		
--	signal 	cnt_phase_UD_270				: std_logic			:= '1';		--三角波UP/DOWNフラグ		
--	signal 	cnt_phase_UD_315				: std_logic			:= '1';		--三角波UP/DOWNフラグ		
----------------------------------------
-----------------------------------------------------------------------------------------------------------------------------------



	signal nMMC_PWM1_FPGA_buf		: std_logic_vector(47 downto 0)	:= "11111111" & "11111111" & "11111111" & "11111111" & "11111111" & "11111111";
	signal nMMC_PWM2_FPGA_buf		: std_logic_vector(47 downto 0)	:= "11111111" & "11111111" & "11111111" & "11111111" & "11111111" & "11111111";
	signal nMMC_PWM1_FPGA_buf2		: std_logic_vector(47 downto 0)	:= "11111111" & "11111111" & "11111111" & "11111111" & "11111111" & "11111111";
	signal nMMC_PWM2_FPGA_buf2		: std_logic_vector(47 downto 0)	:= "11111111" & "11111111" & "11111111" & "11111111" & "11111111" & "11111111";


	--追加ver2.5--
	signal cnt_400k_reg    : std_logic_vector(15 downto 0):=(others => '0');
	signal cnt_2M_reg      : std_logic_vector(15 downto 0):=(others => '0');
	constant cnt_400k_max  : std_logic_vector(15 downto 0):="0000000001111100";
	constant cnt_2M_max    : std_logic_vector(15 downto 0):="0000000000011000";
	signal clk2M           : std_logic			:= '0';	
	signal clk400k         : std_logic			:= '0';	
	signal tri_clk_B : std_logic			:= '0';	
	-------------
	
	
	signal	cnt_10k_B : std_logic_vector(15 downto 0);--出力：cnt10k
	signal	cnt_20k_B : std_logic_vector(15 downto 0);--出力：cnt20k
	signal	cnt_100k_B : std_logic_vector(15 downto 0);--出力：cnt100k
	signal	cnt_1M_B : std_logic_vector(15 downto 0);--出力：cnt1M
	signal	cnt_soft_B : std_logic_vector(3 downto 0); -- 10k or 20k or 100k --
	signal	cnt_10k_ref_B : std_logic_vector(7 downto 0); --出力：10k_ref
	signal	cnt_20k_ref_B : std_logic_vector(8 downto 0);--出力：20k_ref
	signal	cnt_100k_ref_B : std_logic_vector(10 downto 0);--出力：100k_ref
	signal	counter_1M_B : std_logic_vector(15 downto 0);
	
	signal	clk100k_B : std_logic := '0' ;
	signal	clk500k_B : std_logic := '0' ;
	--20220303追加
	signal	clk1M_B : std_logic := '0' ;
	
	--追加ver2.5.6--	
	signal	CONTROL0 : std_logic_vector(35 downto 0);
	---------------
	


	----------------------
	--デフォルトのMMC_PWM_GENにある変数
	----------------------
	signal SPR					: std_logic_vector(15 downto 0)		:= "00000000" & "00000000";

	signal INT_EN				: std_logic_vector(3 downto 0)		:= "0000";

	signal CNT_MAX_SET			: integer range 0 to 32767			:= 0;
	signal DT_SET					: integer range 0 to 1023			:= 1023;
	signal CNT_PHASE_SET_UP1	: integer range -32767 to 32767		:= 0;
	signal CNT_PHASE_SET_UP2	: integer range -32767 to 32767		:= 0;
	signal CNT_PHASE_SET_UP3	: integer range -32767 to 32767		:= 0;
	signal CNT_PHASE_SET_UP4	: integer range -32767 to 32767		:= 0;
	signal CNT_PHASE_SET_UN1	: integer range -32767 to 32767		:= 0;
	signal CNT_PHASE_SET_UN2	: integer range -32767 to 32767		:= 0;
	signal CNT_PHASE_SET_UN3	: integer range -32767 to 32767		:= 0;
	signal CNT_PHASE_SET_UN4	: integer range -32767 to 32767		:= 0;
	signal CNT_PHASE_SET_VP1	: integer range -32767 to 32767		:= 0;
	signal CNT_PHASE_SET_VP2	: integer range -32767 to 32767		:= 0;
	signal CNT_PHASE_SET_VP3	: integer range -32767 to 32767		:= 0;
	signal CNT_PHASE_SET_VP4	: integer range -32767 to 32767		:= 0;
	signal CNT_PHASE_SET_VN1	: integer range -32767 to 32767		:= 0;
	signal CNT_PHASE_SET_VN2	: integer range -32767 to 32767		:= 0;
	signal CNT_PHASE_SET_VN3	: integer range -32767 to 32767		:= 0;
	signal CNT_PHASE_SET_VN4	: integer range -32767 to 32767		:= 0;
	signal CNT_PHASE_SET_WP1	: integer range -32767 to 32767		:= 0;
	signal CNT_PHASE_SET_WP2	: integer range -32767 to 32767		:= 0;
	signal CNT_PHASE_SET_WP3	: integer range -32767 to 32767		:= 0;
	signal CNT_PHASE_SET_WP4	: integer range -32767 to 32767		:= 0;
	signal CNT_PHASE_SET_WN1	: integer range -32767 to 32767		:= 0;
	signal CNT_PHASE_SET_WN2	: integer range -32767 to 32767		:= 0;
	signal CNT_PHASE_SET_WN3	: integer range -32767 to 32767		:= 0;
	signal CNT_PHASE_SET_WN4	: integer range -32767 to 32767		:= 0;
	signal CNT_PHASE_UD_SET_UP1	: std_logic							:= '0';
	signal CNT_PHASE_UD_SET_UP2	: std_logic							:= '0';
	signal CNT_PHASE_UD_SET_UP3	: std_logic							:= '0';
	signal CNT_PHASE_UD_SET_UP4	: std_logic							:= '0';
	signal CNT_PHASE_UD_SET_UN1	: std_logic							:= '0';
	signal CNT_PHASE_UD_SET_UN2	: std_logic							:= '0';
	signal CNT_PHASE_UD_SET_UN3	: std_logic							:= '0';
	signal CNT_PHASE_UD_SET_UN4	: std_logic							:= '0';
	signal CNT_PHASE_UD_SET_VP1	: std_logic							:= '0';
	signal CNT_PHASE_UD_SET_VP2	: std_logic							:= '0';
	signal CNT_PHASE_UD_SET_VP3	: std_logic							:= '0';
	signal CNT_PHASE_UD_SET_VP4	: std_logic							:= '0';
	signal CNT_PHASE_UD_SET_VN1	: std_logic							:= '0';
	signal CNT_PHASE_UD_SET_VN2	: std_logic							:= '0';
	signal CNT_PHASE_UD_SET_VN3	: std_logic							:= '0';
	signal CNT_PHASE_UD_SET_VN4	: std_logic							:= '0';
	signal CNT_PHASE_UD_SET_WP1	: std_logic							:= '0';
	signal CNT_PHASE_UD_SET_WP2	: std_logic							:= '0';
	signal CNT_PHASE_UD_SET_WP3	: std_logic							:= '0';
	signal CNT_PHASE_UD_SET_WP4	: std_logic							:= '0';
	signal CNT_PHASE_UD_SET_WN1	: std_logic							:= '0';
	signal CNT_PHASE_UD_SET_WN2	: std_logic							:= '0';
	signal CNT_PHASE_UD_SET_WN3	: std_logic							:= '0';
	signal CNT_PHASE_UD_SET_WN4	: std_logic							:= '0';
	signal CNT_EN				: std_logic							:= '0';
	signal GATE_EN				: std_logic							:= '0';
	signal REF_UP1				: integer range -32767 to 32767		:= 0;
	signal REF_UP2				: integer range -32767 to 32767		:= 0;
	signal REF_UP3				: integer range -32767 to 32767		:= 0;
	signal REF_UP4				: integer range -32767 to 32767		:= 0;
	signal REF_UN1				: integer range -32767 to 32767		:= 0;
	signal REF_UN2				: integer range -32767 to 32767		:= 0;
	signal REF_UN3				: integer range -32767 to 32767		:= 0;
	signal REF_UN4				: integer range -32767 to 32767		:= 0;
	signal REF_VP1				: integer range -32767 to 32767		:= 0;
	signal REF_VP2				: integer range -32767 to 32767		:= 0;
	signal REF_VP3				: integer range -32767 to 32767		:= 0;
	signal REF_VP4				: integer range -32767 to 32767		:= 0;
	signal REF_VN1				: integer range -32767 to 32767		:= 0;
	signal REF_VN2				: integer range -32767 to 32767		:= 0;
	signal REF_VN3				: integer range -32767 to 32767		:= 0;
	signal REF_VN4				: integer range -32767 to 32767		:= 0;
	signal REF_WP1				: integer range -32767 to 32767		:= 0;
	signal REF_WP2				: integer range -32767 to 32767		:= 0;
	signal REF_WP3				: integer range -32767 to 32767		:= 0;
	signal REF_WP4				: integer range -32767 to 32767		:= 0;
	signal REF_WN1				: integer range -32767 to 32767		:= 0;
	signal REF_WN2				: integer range -32767 to 32767		:= 0;
	signal REF_WN3				: integer range -32767 to 32767		:= 0;
	signal REF_WN4				: integer range -32767 to 32767		:= 0;
	

	-----------------------------
	--ver3で追加
	signal	cnt_10k_0_B : std_logic_vector(15 downto 0);
	signal	cnt_10k_90_B : std_logic_vector(15 downto 0);
	signal	cnt_10k_180_B : std_logic_vector(15 downto 0);
	signal	cnt_10k_270_B : std_logic_vector(15 downto 0);
	
	signal 	MRate_UP1_C_4_int_hold			: integer range -32767 to 32767		:= 0;
	signal 	MRate_UP2_C_4_int_hold			: integer range -32767 to 32767		:= 0;
	signal 	MRate_UP3_C_4_int_hold			: integer range -32767 to 32767		:= 0;
	signal 	MRate_UP4_C_4_int_hold			: integer range -32767 to 32767		:= 0;
	signal 	MRate_UN1_C_4_int_hold			: integer range -32767 to 32767		:= 0;
	signal 	MRate_UN2_C_4_int_hold			: integer range -32767 to 32767		:= 0;
	signal 	MRate_UN3_C_4_int_hold			: integer range -32767 to 32767		:= 0;
	signal 	MRate_UN4_C_4_int_hold			: integer range -32767 to 32767		:= 0;
	signal 	MRate_VP1_C_4_int_hold			: integer range -32767 to 32767		:= 0;
	signal 	MRate_VP2_C_4_int_hold			: integer range -32767 to 32767		:= 0;
	signal 	MRate_VP3_C_4_int_hold			: integer range -32767 to 32767		:= 0;
	signal 	MRate_VP4_C_4_int_hold			: integer range -32767 to 32767		:= 0;
	signal 	MRate_VN1_C_4_int_hold			: integer range -32767 to 32767		:= 0;
	signal 	MRate_VN2_C_4_int_hold			: integer range -32767 to 32767		:= 0;
	signal 	MRate_VN3_C_4_int_hold			: integer range -32767 to 32767		:= 0;
	signal 	MRate_VN4_C_4_int_hold			: integer range -32767 to 32767		:= 0;
	
	--ver3.3で追加
	signal 	REF_hold_out_UP1			: integer range -32767 to 32767		:= 0;
	signal 	REF_hold_out_UP2			: integer range -32767 to 32767		:= 0;
	signal 	REF_hold_out_UP3			: integer range -32767 to 32767		:= 0;
	signal 	REF_hold_out_UP4			: integer range -32767 to 32767		:= 0;
	signal 	REF_hold_out_UN1			: integer range -32767 to 32767		:= 0;
	signal 	REF_hold_out_UN2			: integer range -32767 to 32767		:= 0;
	signal 	REF_hold_out_UN3			: integer range -32767 to 32767		:= 0;
	signal 	REF_hold_out_UN4			: integer range -32767 to 32767		:= 0;
	signal 	REF_hold_out_VP1			: integer range -32767 to 32767		:= 0;
	signal 	REF_hold_out_VP2			: integer range -32767 to 32767		:= 0;
	signal 	REF_hold_out_VP3			: integer range -32767 to 32767		:= 0;
	signal 	REF_hold_out_VP4			: integer range -32767 to 32767		:= 0;
	signal 	REF_hold_out_VN1			: integer range -32767 to 32767		:= 0;
	signal 	REF_hold_out_VN2			: integer range -32767 to 32767		:= 0;
	signal 	REF_hold_out_VN3			: integer range -32767 to 32767		:= 0;
	signal 	REF_hold_out_VN4			: integer range -32767 to 32767		:= 0;
	
	signal 	REF_hold_out_UP1_vector					: std_logic_vector(15 downto 0):=X"0000";
	signal 	REF_hold_out_UP2_vector					: std_logic_vector(15 downto 0):=X"0000";
	signal 	REF_hold_out_UP3_vector					: std_logic_vector(15 downto 0):=X"0000";
	signal 	REF_hold_out_UP4_vector					: std_logic_vector(15 downto 0):=X"0000";
	signal 	REF_hold_out_UN1_vector					: std_logic_vector(15 downto 0):=X"0000";
	signal 	REF_hold_out_UN2_vector					: std_logic_vector(15 downto 0):=X"0000";
	signal 	REF_hold_out_UN3_vector					: std_logic_vector(15 downto 0):=X"0000";
	signal 	REF_hold_out_UN4_vector					: std_logic_vector(15 downto 0):=X"0000";
	signal 	REF_hold_out_VP1_vector					: std_logic_vector(15 downto 0):=X"0000";
	signal 	REF_hold_out_VP2_vector					: std_logic_vector(15 downto 0):=X"0000";
	signal 	REF_hold_out_VP3_vector					: std_logic_vector(15 downto 0):=X"0000";
	signal 	REF_hold_out_VP4_vector					: std_logic_vector(15 downto 0):=X"0000";
	signal 	REF_hold_out_VN1_vector					: std_logic_vector(15 downto 0):=X"0000";
	signal 	REF_hold_out_VN2_vector					: std_logic_vector(15 downto 0):=X"0000";
	signal 	REF_hold_out_VN3_vector					: std_logic_vector(15 downto 0):=X"0000";
	signal 	REF_hold_out_VN4_vector					: std_logic_vector(15 downto 0):=X"0000";
	
	signal   nMMC_PWM1_FPGA_buf2_UP1 : std_logic_vector(0 downto 0);
	
	--追加ver2.3--
		component MMC_PWM_FPGA_SUB is
		Port (
		CLK100				: in  std_logic;
		RESET					: in	std_logic;
		
		CNT_MAX_SET			: in	integer range 0 to 32767;
		CNT_PHASE_SET		: in	integer range -32767 to 32767;
		CNT_PHASE_UD_SET	: in	std_logic;
		DT_SET				: in	integer range 0 to 1023;
		CNT_EN				: in	std_logic;
		GATE_EN				: in	std_logic;
		REF					: in	integer range -32767 to 32767;
		nGATE					: out	std_logic_vector(3 downto 0)	:= "1111";
		
		--追加--
		REF_hold_out			: out integer range -32767 to 32767;
		CNT_out					: out integer range -32767 to 32767
--		CNT_hold1				: out integer range -32767 to 32767;
--		CNT_hold2				: out integer range -32767 to 32767;
--		CNT_hold3				: out integer range -32767 to 32767;
--		CNT_hold4				: out integer range -32767 to 32767
		-------

			);
		end component;
		-----------

	component cnt_clk_3 is
		Port (
			CLK100 : in std_logic;--入力：クロック
			cnt_10k : out std_logic_vector(15 downto 0);--出力：cnt10k
			cnt_20k : out std_logic_vector(15 downto 0);--出力：cnt20k
			cnt_100k : out std_logic_vector(15 downto 0);--出力：cnt100k
			cnt_1M : out std_logic_vector(15 downto 0);--出力：cnt1M
			cnt_soft : out std_logic_vector(3 downto 0); -- 10k or 20k or 100k --
			cnt_10k_ref : out std_logic_vector(7 downto 0); --出力：10k_ref
			cnt_20k_ref : out std_logic_vector(8 downto 0);--出力：20k_ref
			cnt_100k_ref : out std_logic_vector(10 downto 0);--出力：100k_ref
			counter_1M : out std_logic_vector(15 downto 0);
			clk100k : out std_logic;
			clk500k : out std_logic;
			clk1M : out std_logic;
			CNT_EN : in std_logic;
			cnt_10k_0 : out std_logic_vector(15 downto 0);
			cnt_10k_90 : out std_logic_vector(15 downto 0);
			cnt_10k_180 : out std_logic_vector(15 downto 0);
			cnt_10k_270 : out std_logic_vector(15 downto 0)
			);
	end component;



	--追加ver2.5.6--
--		component icon1 is
--			Port (
--			  CONTROL0                          :   INOUT   std_logic_vector(35 DOWNTO 0)
--			);
--		end component;
--
--
--	--見る信号の変数のみコメントアウトを外す
--	--このとき、最後となる変数のセミコロンを外すのを忘れないこと
--		component ila1 is
--			Port (
--			  CONTROL                           :   IN   std_logic_vector(35 DOWNTO 0);
--			  CLK		                           :   IN   std_logic;
--			  
--			  TRIG0	                           :   IN   std_logic_vector(15 DOWNTO 0); --制御量u(正規化あと:UP1)
----			  TRIG1	                           : 	 IN   std_logic_vector(15 DOWNTO 0); --制御量u(headspring:UP1)
----			  TRIG2	                           : 	 IN   std_logic_vector(15 DOWNTO 0); --三角波:UP1			  
--			  TRIG1	                           : 	 IN   std_logic_vector(15 DOWNTO 0); --平均値制御出力(u相)
--			  TRIG2	                           : 	 IN   std_logic_vector(15 DOWNTO 0); --バランス制御出力(up1)
--			  TRIG3	                           : 	 IN   std_logic_vector(15 DOWNTO 0); --Iup_10kHz			  
--			  TRIG4	                           : 	 IN   std_logic_vector(15 DOWNTO 0); --Voref
--			  TRIG5	                           : 	 IN   std_logic_vector(15 DOWNTO 0); --dVoref
--			  TRIG6	                           : 	 IN   std_logic_vector(15 DOWNTO 0); --Vo              
--			  TRIG7	                           : 	 IN   std_logic_vector(15 DOWNTO 0); --IC
--			  TRIG8	                   		   : 	 IN   std_logic_vector(15 DOWNTO 0); --QMDBout
----			  TRIG9	                   		   : 	 IN   std_logic_vector(15 DOWNTO 0); --Vdc_N
--			  TRIG9	                   		   : 	 IN   std_logic_vector(15 DOWNTO 0); --soft_cnt_out_BB,cnt80k_debug_BB
--			  TRIG10	                   		   : 	 IN   std_logic_vector(0 downto 0) --clk10k_BBB
--			  
----			  TRIG0	                           :   IN   std_logic_vector(15 DOWNTO 0); --制御量u(正規化前:UP1)
----			  TRIG1	                           :   IN   std_logic_vector(15 DOWNTO 0); --制御量u(正規化あと_hold後:UP1)
----			  TRIG4	                           : 	 IN   std_logic_vector(0 DOWNTO 0);--ゲート信号:UP1
----			  TRIG7	                           : 	 IN   std_logic_vector(15 DOWNTO 0)--平均値制御出力リミッタなし(u相)
----			  TRIG8	                           : 	 IN   std_logic_vector(15 DOWNTO 0);--Iup_100MHz
--
--
----			  TRIG1									   :   IN   std_logic_vector(15 DOWNTO 0); --制御量u(正規化前:VP1)
----			  TRIG3	                           :   IN   std_logic_vector(13 DOWNTO 0); --制御量u(正規化あと:VP1)			  
----			  TRIG5	                           : 	 IN   std_logic_vector(13 DOWNTO 0)  --制御量u(三角波:VP1)
--
----			  TRIG6	                           : 	 IN   integer range 0 to 16384; 
----			  TRIG7	                           : 	 IN   integer range 0 to 16384; 
----			  TRIG8	                           : 	 IN   integer range -mrate_max to mrate_max;  
----			  TRIG9	                           : 	 IN   integer range -mrate_max to mrate_max;  
----			  TRIG10	                           : 	 IN   integer range -mrate_max to mrate_max;  
----			  TRIG11	                           : 	 IN   integer range -mrate_max to mrate_max;  
----			  TRIG12	                           :   IN   integer range 0 to 2500;  
----			  TRIG13	                           :   IN   integer range 0 to 2500;  
----			  TRIG14	                           :   IN   integer range 0 to 2500;  
----			  TRIG15	                           :   IN   integer range 0 to 2500
--
--			);
--		end component;
	---------------
	
begin

----ver2.6.1

	process(CLK100)
	begin
		if CLK100'event and CLK100 = '1' then
		
----------------
		--cnt80k=X"003D"
		--制御量ΔTを三角波に合わせて5000倍
		--5000倍した制御量ΔTのmaxは----となる
			MRate_UP1_C <= MRate_UP1 * pwm_assist;
			MRate_UP2_C <= MRate_UP2 * pwm_assist;
			MRate_UP3_C <= MRate_UP3 * pwm_assist;
			MRate_UP4_C <= MRate_UP4 * pwm_assist;
			
			MRate_UN1_C <= MRate_UN1 * pwm_assist;
			MRate_UN2_C <= MRate_UN2 * pwm_assist;
			MRate_UN3_C <= MRate_UN3 * pwm_assist;
			MRate_UN4_C <= MRate_UN4 * pwm_assist;

			MRate_VP1_C <= MRate_VP1 * pwm_assist;
			MRate_VP2_C <= MRate_VP2 * pwm_assist;
			MRate_VP3_C <= MRate_VP3 * pwm_assist;
			MRate_VP4_C <= MRate_VP4 * pwm_assist;
			
			MRate_VN1_C <= MRate_VN1 * pwm_assist;
			MRate_VN2_C <= MRate_VN2 * pwm_assist;
			MRate_VN3_C <= MRate_VN3 * pwm_assist;
			MRate_VN4_C <= MRate_VN4 * pwm_assist;
			
		--cnt80k=X"003E"
		--これを14bit(16384)シフトする(整数部の取得)で0~5000にする。
			MRate_UP1_C_2 <= MRate_UP1_C(29 downto 14);
			MRate_UP2_C_2 <= MRate_UP2_C(29 downto 14);
			MRate_UP3_C_2 <= MRate_UP3_C(29 downto 14);
			MRate_UP4_C_2 <= MRate_UP4_C(29 downto 14);
			
			MRate_UN1_C_2 <= MRate_UN1_C(29 downto 14);
			MRate_UN2_C_2 <= MRate_UN2_C(29 downto 14);
			MRate_UN3_C_2 <= MRate_UN3_C(29 downto 14);
			MRate_UN4_C_2 <= MRate_UN4_C(29 downto 14);

			MRate_VP1_C_2 <= MRate_VP1_C(29 downto 14);
			MRate_VP2_C_2 <= MRate_VP2_C(29 downto 14);
			MRate_VP3_C_2 <= MRate_VP3_C(29 downto 14);
			MRate_VP4_C_2 <= MRate_VP4_C(29 downto 14);
			
			MRate_VN1_C_2 <= MRate_VN1_C(29 downto 14);
			MRate_VN2_C_2 <= MRate_VN2_C(29 downto 14);
			MRate_VN3_C_2 <= MRate_VN3_C(29 downto 14);
			MRate_VN4_C_2 <= MRate_VN4_C(29 downto 14);

		--cnt80k=X"003F"
		--オフセット-2500してheadspringの三角波比較(10kHz)に合わせる
			MRate_UP1_C_3 <= MRate_UP1_C_2 - pwm_assist_half;
			MRate_UP2_C_3 <= MRate_UP2_C_2 - pwm_assist_half;
			MRate_UP3_C_3 <= MRate_UP3_C_2 - pwm_assist_half;
			MRate_UP4_C_3 <= MRate_UP4_C_2 - pwm_assist_half;

			MRate_UN1_C_3 <= MRate_UN1_C_2 - pwm_assist_half;
			MRate_UN2_C_3 <= MRate_UN2_C_2 - pwm_assist_half;
			MRate_UN3_C_3 <= MRate_UN3_C_2 - pwm_assist_half;
			MRate_UN4_C_3 <= MRate_UN4_C_2 - pwm_assist_half;
			
			MRate_VP1_C_3 <= MRate_VP1_C_2 - pwm_assist_half;
			MRate_VP2_C_3 <= MRate_VP2_C_2 - pwm_assist_half;
			MRate_VP3_C_3 <= MRate_VP3_C_2 - pwm_assist_half;
			MRate_VP4_C_3 <= MRate_VP4_C_2 - pwm_assist_half;

			MRate_VN1_C_3 <= MRate_VN1_C_2 - pwm_assist_half;
			MRate_VN2_C_3 <= MRate_VN2_C_2 - pwm_assist_half;
			MRate_VN3_C_3 <= MRate_VN3_C_2 - pwm_assist_half;
			MRate_VN4_C_3 <= MRate_VN4_C_2 - pwm_assist_half;

		--ver2.6

		-------------------------------
		--リミッタ処理なし
		-------------------------------
		--	MRate_UP1_C_4 <= MRate_UP1_C_3;
		--	MRate_UP2_C_4 <= MRate_UP2_C_3;
		--	MRate_UP3_C_4 <= MRate_UP3_C_3;
		--	MRate_UP4_C_4 <= MRate_UP4_C_3;
		--	
		--	MRate_UN1_C_4 <= MRate_UN1_C_3;
		--	MRate_UN2_C_4 <= MRate_UN2_C_3;
		--	MRate_UN3_C_4 <= MRate_UN3_C_3;
		--	MRate_UN4_C_4 <= MRate_UN4_C_3;
		--
		--	MRate_VP1_C_4 <= MRate_VP1_C_3;
		--	MRate_VP2_C_4 <= MRate_VP2_C_3;
		--	MRate_VP3_C_4 <= MRate_VP3_C_3;
		--	MRate_VP4_C_4 <= MRate_VP4_C_3;
		--	
		--	MRate_VN1_C_4 <= MRate_VN1_C_3;
		--	MRate_VN2_C_4 <= MRate_VN2_C_3;
		--	MRate_VN3_C_4 <= MRate_VN3_C_3;
		--	MRate_VN4_C_4 <= MRate_VN4_C_3;

		-------------------------------
		--リミッタ処理あり
		-------------------------------

		--cnt80k=X"0040"
		--UP1
			if( MRate_UP1_C_3 > upper_limit ) then
				MRate_UP1_C_4 <= upper_limit;
			elsif( MRate_UP1_C_3	< lower_limit ) then
				MRate_UP1_C_4 <= lower_limit;
			else
				MRate_UP1_C_4 <= MRate_UP1_C_3;
			end if;
		--UP2
			if( MRate_UP2_C_3 > upper_limit ) then
				MRate_UP2_C_4 <= upper_limit;
			elsif( MRate_UP2_C_3	< lower_limit ) then
				MRate_UP2_C_4 <= lower_limit;
			else
				MRate_UP2_C_4 <= MRate_UP2_C_3;
			end if;
		--UP3
			if( MRate_UP3_C_3 > upper_limit ) then
				MRate_UP3_C_4 <= upper_limit;
			elsif( MRate_UP3_C_3	< lower_limit ) then
				MRate_UP3_C_4 <= lower_limit;
			else
				MRate_UP3_C_4 <= MRate_UP3_C_3;
			end if;
		--UP4
			if( MRate_UP4_C_3 > upper_limit ) then
				MRate_UP4_C_4 <= upper_limit;
			elsif( MRate_UP4_C_3	< lower_limit ) then
				MRate_UP4_C_4 <= lower_limit;
			else
				MRate_UP4_C_4 <= MRate_UP4_C_3;
			end if;
		--UN1
			if( MRate_UN1_C_3 > upper_limit ) then
				MRate_UN1_C_4 <= upper_limit;
			elsif( MRate_UN1_C_3	< lower_limit ) then
				MRate_UN1_C_4 <= lower_limit;
			else
				MRate_UN1_C_4 <= MRate_UN1_C_3;
			end if;
		--UN2
			if( MRate_UN2_C_3 > upper_limit ) then
				MRate_UN2_C_4 <= upper_limit;
			elsif( MRate_UN2_C_3	< lower_limit ) then
				MRate_UN2_C_4 <= lower_limit;
			else
				MRate_UN2_C_4 <= MRate_UN2_C_3;
			end if;
		--UN3
			if( MRate_UN3_C_3 > upper_limit ) then
				MRate_UN3_C_4 <= upper_limit;
			elsif( MRate_UN3_C_3	< lower_limit ) then
				MRate_UN3_C_4 <= lower_limit;
			else
				MRate_UN3_C_4 <= MRate_UN3_C_3;
			end if;
		--UN4
			if( MRate_UN4_C_3 > upper_limit ) then
				MRate_UN4_C_4 <= upper_limit;
			elsif( MRate_UN4_C_3	< lower_limit ) then
				MRate_UN4_C_4 <= lower_limit;
			else
				MRate_UN4_C_4 <= MRate_UN4_C_3;
			end if;

		--VP1
			if( MRate_VP1_C_3 > upper_limit ) then
				MRate_VP1_C_4 <= upper_limit;
			elsif( MRate_VP1_C_3	< lower_limit ) then
				MRate_VP1_C_4 <= lower_limit;
			else
				MRate_VP1_C_4 <= MRate_VP1_C_3;
			end if;
		--VP2
			if( MRate_VP2_C_3 > upper_limit ) then
				MRate_VP2_C_4 <= upper_limit;
			elsif( MRate_VP2_C_3	< lower_limit ) then
				MRate_VP2_C_4 <= lower_limit;
			else
				MRate_VP2_C_4 <= MRate_VP2_C_3;
			end if;
		--VP3
			if( MRate_VP3_C_3 > upper_limit ) then
				MRate_VP3_C_4 <= upper_limit;
			elsif( MRate_VP3_C_3	< lower_limit ) then
				MRate_VP3_C_4 <= lower_limit;
			else
				MRate_VP3_C_4 <= MRate_VP3_C_3;
			end if;
		--VP4
			if( MRate_VP4_C_3 > upper_limit ) then
				MRate_VP4_C_4 <= upper_limit;
			elsif( MRate_VP4_C_3	< lower_limit ) then
				MRate_VP4_C_4 <= lower_limit;
			else
				MRate_VP4_C_4 <= MRate_VP4_C_3;
			end if;
		--VN1
			if( MRate_VN1_C_3 > upper_limit ) then
				MRate_VN1_C_4 <= upper_limit;
			elsif( MRate_VN1_C_3	< lower_limit ) then
				MRate_VN1_C_4 <= lower_limit;
			else
				MRate_VN1_C_4 <= MRate_VN1_C_3;
			end if;
		--VN2
			if( MRate_VN2_C_3 > upper_limit ) then
				MRate_VN2_C_4 <= upper_limit;
			elsif( MRate_VN2_C_3	< lower_limit ) then
				MRate_VN2_C_4 <= lower_limit;
			else
				MRate_VN2_C_4 <= MRate_VN2_C_3;
			end if;
		--VN3
			if( MRate_VN3_C_3 > upper_limit ) then
				MRate_VN3_C_4 <= upper_limit;
			elsif( MRate_VN3_C_3	< lower_limit ) then
				MRate_VN3_C_4 <= lower_limit;
			else
				MRate_VN3_C_4 <= MRate_VN3_C_3;
			end if;
		--VN4
			if( MRate_VN4_C_3 > upper_limit ) then
				MRate_VN4_C_4 <= upper_limit;
			elsif( MRate_VN4_C_3	< lower_limit ) then
				MRate_VN4_C_4 <= lower_limit;
			else
				MRate_VN4_C_4 <= MRate_VN4_C_3;
			end if;

------------

		--	--hold処理 -> 下の三角波モジュールで行う--
		----	if(cnt_10k_0_B = X"0000")then
		--		MRate_UP1_C_4_hold <= MRate_UP1_C_4;
		--		MRate_UN1_C_4_hold <= MRate_UN1_C_4;
		--		MRate_VP1_C_4_hold <= MRate_VP1_C_4;
		--		MRate_VN1_C_4_hold <= MRate_VN1_C_4;
		----	elsif(cnt_10k_90_B = X"0000")then
		--		MRate_UP2_C_4_hold <= MRate_UP2_C_4;
		--		MRate_UN2_C_4_hold <= MRate_UN2_C_4;
		--		MRate_VP2_C_4_hold <= MRate_VP2_C_4;
		--		MRate_VN2_C_4_hold <= MRate_VN2_C_4;
		----	elsif(cnt_10k_180_B = X"0000")then		
		--		MRate_UP3_C_4_hold <= MRate_UP3_C_4;
		--		MRate_UN3_C_4_hold <= MRate_UN3_C_4;
		--		MRate_VP3_C_4_hold <= MRate_VP3_C_4;
		--		MRate_VN3_C_4_hold <= MRate_VN3_C_4;
		----	elsif(cnt_10k_270_B = X"0000")then		
		--		MRate_UP4_C_4_hold <= MRate_UP4_C_4;
		--		MRate_UN4_C_4_hold <= MRate_UN4_C_4;
		--		MRate_VP4_C_4_hold <= MRate_VP4_C_4;
		--		MRate_VN4_C_4_hold <= MRate_VN4_C_4;
		----	end if;
		--	
		--	--制御入力をvector型に変更して下のモジュールへ
		--	MRate_UP1_C_4_int_hold <= CONV_INTEGER(MRate_UP1_C_4_hold);
		--	MRate_UP2_C_4_int_hold <= CONV_INTEGER(MRate_UP2_C_4_hold);
		--	MRate_UP3_C_4_int_hold <= CONV_INTEGER(MRate_UP3_C_4_hold);
		--	MRate_UP4_C_4_int_hold <= CONV_INTEGER(MRate_UP4_C_4_hold);
		--
		--	MRate_UN1_C_4_int_hold <= CONV_INTEGER(MRate_UN1_C_4_hold);
		--	MRate_UN2_C_4_int_hold <= CONV_INTEGER(MRate_UN2_C_4_hold);
		--	MRate_UN3_C_4_int_hold <= CONV_INTEGER(MRate_UN3_C_4_hold);
		--	MRate_UN4_C_4_int_hold <= CONV_INTEGER(MRate_UN4_C_4_hold);
		--
		--	MRate_VP1_C_4_int_hold <= CONV_INTEGER(MRate_VP1_C_4_hold);
		--	MRate_VP2_C_4_int_hold <= CONV_INTEGER(MRate_VP2_C_4_hold);
		--	MRate_VP3_C_4_int_hold <= CONV_INTEGER(MRate_VP3_C_4_hold);
		--	MRate_VP4_C_4_int_hold <= CONV_INTEGER(MRate_VP4_C_4_hold);
		--
		--	MRate_VN1_C_4_int_hold <= CONV_INTEGER(MRate_VN1_C_4_hold);
		--	MRate_VN2_C_4_int_hold <= CONV_INTEGER(MRate_VN2_C_4_hold);
		--	MRate_VN3_C_4_int_hold <= CONV_INTEGER(MRate_VN3_C_4_hold);
		--	MRate_VN4_C_4_int_hold <= CONV_INTEGER(MRate_VN4_C_4_hold);
			
		--cnt80k=X"0041"
		--制御入力をvector型に変更して下のモジュールへ
			MRate_UP1_C_4_int_hold <= CONV_INTEGER(MRate_UP1_C_4);
			MRate_UP2_C_4_int_hold <= CONV_INTEGER(MRate_UP2_C_4);
			MRate_UP3_C_4_int_hold <= CONV_INTEGER(MRate_UP3_C_4);
			MRate_UP4_C_4_int_hold <= CONV_INTEGER(MRate_UP4_C_4);

			MRate_UN1_C_4_int_hold <= CONV_INTEGER(MRate_UN1_C_4);
			MRate_UN2_C_4_int_hold <= CONV_INTEGER(MRate_UN2_C_4);
			MRate_UN3_C_4_int_hold <= CONV_INTEGER(MRate_UN3_C_4);
			MRate_UN4_C_4_int_hold <= CONV_INTEGER(MRate_UN4_C_4);

			MRate_VP1_C_4_int_hold <= CONV_INTEGER(MRate_VP1_C_4);
			MRate_VP2_C_4_int_hold <= CONV_INTEGER(MRate_VP2_C_4);
			MRate_VP3_C_4_int_hold <= CONV_INTEGER(MRate_VP3_C_4);
			MRate_VP4_C_4_int_hold <= CONV_INTEGER(MRate_VP4_C_4);

			MRate_VN1_C_4_int_hold <= CONV_INTEGER(MRate_VN1_C_4);
			MRate_VN2_C_4_int_hold <= CONV_INTEGER(MRate_VN2_C_4);
			MRate_VN3_C_4_int_hold <= CONV_INTEGER(MRate_VN3_C_4);
			MRate_VN4_C_4_int_hold <= CONV_INTEGER(MRate_VN4_C_4);
			
		--三角波をint型に変更してchipscopeでデバッグ------------------------------------------
			CNT_out_UP1_vector	 <= CONV_STD_LOGIC_VECTOR(CNT_out_UP1,16);
			CNT_out_UP2_vector	 <= CONV_STD_LOGIC_VECTOR(CNT_out_UP2,16);
			CNT_out_UP3_vector	 <= CONV_STD_LOGIC_VECTOR(CNT_out_UP3,16);
			CNT_out_UP4_vector	 <= CONV_STD_LOGIC_VECTOR(CNT_out_UP4,16);
			CNT_out_UN1_vector	 <= CONV_STD_LOGIC_VECTOR(CNT_out_UN1,16);
			CNT_out_UN2_vector	 <= CONV_STD_LOGIC_VECTOR(CNT_out_UN2,16);
			CNT_out_UN3_vector	 <= CONV_STD_LOGIC_VECTOR(CNT_out_UN3,16);
			CNT_out_UN4_vector	 <= CONV_STD_LOGIC_VECTOR(CNT_out_UN4,16);
			CNT_out_VP1_vector	 <= CONV_STD_LOGIC_VECTOR(CNT_out_VP1,16);
			CNT_out_VP2_vector	 <= CONV_STD_LOGIC_VECTOR(CNT_out_VP2,16);
			CNT_out_VP3_vector	 <= CONV_STD_LOGIC_VECTOR(CNT_out_VP3,16);
			CNT_out_VP4_vector	 <= CONV_STD_LOGIC_VECTOR(CNT_out_VP4,16);
			CNT_out_VN1_vector	 <= CONV_STD_LOGIC_VECTOR(CNT_out_VN1,16);
			CNT_out_VN2_vector	 <= CONV_STD_LOGIC_VECTOR(CNT_out_VN2,16);
			CNT_out_VN3_vector	 <= CONV_STD_LOGIC_VECTOR(CNT_out_VN3,16);
			CNT_out_VN4_vector	 <= CONV_STD_LOGIC_VECTOR(CNT_out_VN4,16);
			
			REF_hold_out_UP1_vector	 <= CONV_STD_LOGIC_VECTOR(REF_hold_out_UP1,16);
			REF_hold_out_UP2_vector	 <= CONV_STD_LOGIC_VECTOR(REF_hold_out_UP2,16);
			REF_hold_out_UP3_vector	 <= CONV_STD_LOGIC_VECTOR(REF_hold_out_UP3,16);
			REF_hold_out_UP4_vector	 <= CONV_STD_LOGIC_VECTOR(REF_hold_out_UP4,16);
			REF_hold_out_UN1_vector	 <= CONV_STD_LOGIC_VECTOR(REF_hold_out_UN1,16);
			REF_hold_out_UN2_vector	 <= CONV_STD_LOGIC_VECTOR(REF_hold_out_UN2,16);
			REF_hold_out_UN3_vector	 <= CONV_STD_LOGIC_VECTOR(REF_hold_out_UN3,16);
			REF_hold_out_UN4_vector	 <= CONV_STD_LOGIC_VECTOR(REF_hold_out_UN4,16);
			REF_hold_out_VP1_vector	 <= CONV_STD_LOGIC_VECTOR(REF_hold_out_VP1,16);
			REF_hold_out_VP2_vector	 <= CONV_STD_LOGIC_VECTOR(REF_hold_out_VP2,16);
			REF_hold_out_VP3_vector	 <= CONV_STD_LOGIC_VECTOR(REF_hold_out_VP3,16);
			REF_hold_out_VP4_vector	 <= CONV_STD_LOGIC_VECTOR(REF_hold_out_VP4,16);
			REF_hold_out_VN1_vector	 <= CONV_STD_LOGIC_VECTOR(REF_hold_out_VN1,16);
			REF_hold_out_VN2_vector	 <= CONV_STD_LOGIC_VECTOR(REF_hold_out_VN2,16);
			REF_hold_out_VN3_vector	 <= CONV_STD_LOGIC_VECTOR(REF_hold_out_VN3,16);
			REF_hold_out_VN4_vector	 <= CONV_STD_LOGIC_VECTOR(REF_hold_out_VN4,16);
		-----------------------------------------------------------------------------
		end if;
	end process;

	----------------------
	--デフォルトのMMC_PWM_GENにある変数
	----------------------
	
	process(CLK100)
	begin
		if CLK100'event and CLK100 = '1' then

			if RESET = '1' then
				SPR				<= "00000000" & "00000000";
				CNT_MAX_SET			<= 32767;
				DT_SET				<= 1023;
				CNT_PHASE_SET_UP1	<= 0;
				CNT_PHASE_SET_UP2	<= 0;
				CNT_PHASE_SET_UP3	<= 0;
				CNT_PHASE_SET_UP4	<= 0;
				CNT_PHASE_SET_UN1	<= 0;
				CNT_PHASE_SET_UN2	<= 0;
				CNT_PHASE_SET_UN3	<= 0;
				CNT_PHASE_SET_UN4	<= 0;
				CNT_PHASE_SET_VP1	<= 0;
				CNT_PHASE_SET_VP2	<= 0;
				CNT_PHASE_SET_VP3	<= 0;
				CNT_PHASE_SET_VP4	<= 0;
				CNT_PHASE_SET_VN1	<= 0;
				CNT_PHASE_SET_VN2	<= 0;
				CNT_PHASE_SET_VN3	<= 0;
				CNT_PHASE_SET_VN4	<= 0;
				CNT_PHASE_SET_WP1	<= 0;
				CNT_PHASE_SET_WP2	<= 0;
				CNT_PHASE_SET_WP3	<= 0;
				CNT_PHASE_SET_WP4	<= 0;
				CNT_PHASE_SET_WN1	<= 0;
				CNT_PHASE_SET_WN2	<= 0;
				CNT_PHASE_SET_WN3	<= 0;
				CNT_PHASE_SET_WN4	<= 0;
				CNT_PHASE_UD_SET_UP1	<= '0';
				CNT_PHASE_UD_SET_UP2	<= '0';
				CNT_PHASE_UD_SET_UP3	<= '0';
				CNT_PHASE_UD_SET_UP4	<= '0';
				CNT_PHASE_UD_SET_UN1	<= '0';
				CNT_PHASE_UD_SET_UN2	<= '0';
				CNT_PHASE_UD_SET_UN3	<= '0';
				CNT_PHASE_UD_SET_UN4	<= '0';
				CNT_PHASE_UD_SET_VP1	<= '0';
				CNT_PHASE_UD_SET_VP2	<= '0';
				CNT_PHASE_UD_SET_VP3	<= '0';
				CNT_PHASE_UD_SET_VP4	<= '0';
				CNT_PHASE_UD_SET_VN1	<= '0';
				CNT_PHASE_UD_SET_VN2	<= '0';
				CNT_PHASE_UD_SET_VN3	<= '0';
				CNT_PHASE_UD_SET_VN4	<= '0';
				CNT_PHASE_UD_SET_WP1	<= '0';
				CNT_PHASE_UD_SET_WP2	<= '0';
				CNT_PHASE_UD_SET_WP3	<= '0';
				CNT_PHASE_UD_SET_WP4	<= '0';
				CNT_PHASE_UD_SET_WN1	<= '0';
				CNT_PHASE_UD_SET_WN2	<= '0';
				CNT_PHASE_UD_SET_WN3	<= '0';
				CNT_PHASE_UD_SET_WN4	<= '0';
				CNT_EN					<= '0';
				GATE_EN					<= '0';
				REF_UP1				<= 0;
				REF_UP2				<= 0;
				REF_UP3				<= 0;
				REF_UP4				<= 0;
				REF_UN1				<= 0;
				REF_UN2				<= 0;
				REF_UN3				<= 0;
				REF_UN4				<= 0;
				REF_VP1				<= 0;
				REF_VP2				<= 0;
				REF_VP3				<= 0;
				REF_VP4				<= 0;
				REF_VN1				<= 0;
				REF_VN2				<= 0;
				REF_VN3				<= 0;
				REF_VN4				<= 0;
				REF_WP1				<= 0;
				REF_WP2				<= 0;
				REF_WP3				<= 0;
				REF_WP4				<= 0;
				REF_WN1				<= 0;
				REF_WN2				<= 0;
				REF_WN3				<= 0;
				REF_WN4				<= 0;
				INT_EN				<= "0000";
				
			elsif WR = '1' then
				case A is
					when x"00"	=>	CNT_MAX_SET			<= CONV_INTEGER(DW(15 downto 0));
					when x"01"	=>	DT_SET				<= CONV_INTEGER(DW(10 downto 0));
					when x"02"	=>	CNT_PHASE_SET_UP1	<= CONV_INTEGER(DW(15 downto 0));
					when x"03"	=>	CNT_PHASE_SET_UP2	<= CONV_INTEGER(DW(15 downto 0));
					when x"04"	=>	CNT_PHASE_SET_UP3	<= CONV_INTEGER(DW(15 downto 0));
					when x"05"	=>	CNT_PHASE_SET_UP4	<= CONV_INTEGER(DW(15 downto 0));
					when x"06"	=>	CNT_PHASE_SET_UN1	<= CONV_INTEGER(DW(15 downto 0));
					when x"07"	=>	CNT_PHASE_SET_UN2	<= CONV_INTEGER(DW(15 downto 0));
					when x"08"	=>	CNT_PHASE_SET_UN3	<= CONV_INTEGER(DW(15 downto 0));
					when x"09"	=>	CNT_PHASE_SET_UN4	<= CONV_INTEGER(DW(15 downto 0));
					when x"0A"	=>	CNT_PHASE_SET_VP1	<= CONV_INTEGER(DW(15 downto 0));
					when x"0B"	=>	CNT_PHASE_SET_VP2	<= CONV_INTEGER(DW(15 downto 0));
					when x"0C"	=>	CNT_PHASE_SET_VP3	<= CONV_INTEGER(DW(15 downto 0));
					when x"0D"	=>	CNT_PHASE_SET_VP4	<= CONV_INTEGER(DW(15 downto 0));
					when x"0E"	=>	CNT_PHASE_SET_VN1	<= CONV_INTEGER(DW(15 downto 0));
					when x"0F"	=>	CNT_PHASE_SET_VN2	<= CONV_INTEGER(DW(15 downto 0));
					when x"10"	=>	CNT_PHASE_SET_VN3	<= CONV_INTEGER(DW(15 downto 0));
					when x"11"	=>	CNT_PHASE_SET_VN4	<= CONV_INTEGER(DW(15 downto 0));
					when x"12"	=>	CNT_PHASE_SET_WP1	<= CONV_INTEGER(DW(15 downto 0));
					when x"13"	=>	CNT_PHASE_SET_WP2	<= CONV_INTEGER(DW(15 downto 0));
					when x"14"	=>	CNT_PHASE_SET_WP3	<= CONV_INTEGER(DW(15 downto 0));
					when x"15"	=>	CNT_PHASE_SET_WP4	<= CONV_INTEGER(DW(15 downto 0));
					when x"16"	=>	CNT_PHASE_SET_WN1	<= CONV_INTEGER(DW(15 downto 0));
					when x"17"	=>	CNT_PHASE_SET_WN2	<= CONV_INTEGER(DW(15 downto 0));
					when x"18"	=>	CNT_PHASE_SET_WN3	<= CONV_INTEGER(DW(15 downto 0));
					when x"19"	=>	CNT_PHASE_SET_WN4	<= CONV_INTEGER(DW(15 downto 0));
					when x"1A"	=>	CNT_PHASE_UD_SET_UP1	<= DW(0);
					when x"1B"	=>	CNT_PHASE_UD_SET_UP2	<= DW(0);
					when x"1C"	=>	CNT_PHASE_UD_SET_UP3	<= DW(0);
					when x"1D"	=>	CNT_PHASE_UD_SET_UP4	<= DW(0);
					when x"1E"	=>	CNT_PHASE_UD_SET_UN1	<= DW(0);
					when x"1F"	=>	CNT_PHASE_UD_SET_UN2	<= DW(0);
					when x"20"	=>	CNT_PHASE_UD_SET_UN3	<= DW(0);
					when x"21"	=>	CNT_PHASE_UD_SET_UN4	<= DW(0);
					when x"22"	=>	CNT_PHASE_UD_SET_VP1	<= DW(0);
					when x"23"	=>	CNT_PHASE_UD_SET_VP2	<= DW(0);
					when x"24"	=>	CNT_PHASE_UD_SET_VP3	<= DW(0);
					when x"25"	=>	CNT_PHASE_UD_SET_VP4	<= DW(0);
					when x"26"	=>	CNT_PHASE_UD_SET_VN1	<= DW(0);
					when x"27"	=>	CNT_PHASE_UD_SET_VN2	<= DW(0);
					when x"28"	=>	CNT_PHASE_UD_SET_VN3	<= DW(0);
					when x"29"	=>	CNT_PHASE_UD_SET_VN4	<= DW(0);
					when x"2A"	=>	CNT_PHASE_UD_SET_WP1	<= DW(0);
					when x"2B"	=>	CNT_PHASE_UD_SET_WP2	<= DW(0);
					when x"2C"	=>	CNT_PHASE_UD_SET_WP3	<= DW(0);
					when x"2D"	=>	CNT_PHASE_UD_SET_WP4	<= DW(0);
					when x"2E"	=>	CNT_PHASE_UD_SET_WN1	<= DW(0);
					when x"2F"	=>	CNT_PHASE_UD_SET_WN2	<= DW(0);
					when x"30"	=>	CNT_PHASE_UD_SET_WN3	<= DW(0);
					when x"31"	=>	CNT_PHASE_UD_SET_WN4	<= DW(0);
					when x"32"	=>	CNT_EN					<= DW(0);
					when x"33"	=>	GATE_EN					<= DW(0);
					when x"34"	=>	REF_UP1				<= CONV_INTEGER(DW(15 downto 0));
					when x"35"	=>	REF_UP2				<= CONV_INTEGER(DW(15 downto 0));
					when x"36"	=>	REF_UP3				<= CONV_INTEGER(DW(15 downto 0));
					when x"37"	=>	REF_UP4				<= CONV_INTEGER(DW(15 downto 0));
					when x"38"	=>	REF_UN1				<= CONV_INTEGER(DW(15 downto 0));
					when x"39"	=>	REF_UN2				<= CONV_INTEGER(DW(15 downto 0));
					when x"3A"	=>	REF_UN3				<= CONV_INTEGER(DW(15 downto 0));
					when x"3B"	=>	REF_UN4				<= CONV_INTEGER(DW(15 downto 0));
					when x"3C"	=>	REF_VP1				<= CONV_INTEGER(DW(15 downto 0));
					when x"3D"	=>	REF_VP2				<= CONV_INTEGER(DW(15 downto 0));
					when x"3E"	=>	REF_VP3				<= CONV_INTEGER(DW(15 downto 0));
					when x"3F"	=>	REF_VP4				<= CONV_INTEGER(DW(15 downto 0));
					when x"40"	=>	REF_VN1				<= CONV_INTEGER(DW(15 downto 0));
					when x"41"	=>	REF_VN2				<= CONV_INTEGER(DW(15 downto 0));
					when x"42"	=>	REF_VN3				<= CONV_INTEGER(DW(15 downto 0));
					when x"43"	=>	REF_VN4				<= CONV_INTEGER(DW(15 downto 0));
					when x"44"	=>	REF_WP1				<= CONV_INTEGER(DW(15 downto 0));
					when x"45"	=>	REF_WP2				<= CONV_INTEGER(DW(15 downto 0));
					when x"46"	=>	REF_WP3				<= CONV_INTEGER(DW(15 downto 0));
					when x"47"	=>	REF_WP4				<= CONV_INTEGER(DW(15 downto 0));
					when x"48"	=>	REF_WN1				<= CONV_INTEGER(DW(15 downto 0));
					when x"49"	=>	REF_WN2				<= CONV_INTEGER(DW(15 downto 0));
					when x"4A"	=>	REF_WN3				<= CONV_INTEGER(DW(15 downto 0));
					when x"4B"	=>	REF_WN4				<= CONV_INTEGER(DW(15 downto 0));
					when x"4C"	=>	INT_EN				<= DW(3 downto 0);
					
				
					when x"FF"	=>	SPR			<= DW;
					when others		=> null;
				end case;
			else
				if GB = '1' then
					GATE_EN	<= '0';
				end if;
			end if;
			
		end if;
	end process;


-- ====================================================
-- MMC_PWM_FPGA_UP1
-- ====================================================

	MMCPWM_FPGA_UP1 : MMC_PWM_FPGA_SUB
		port map(
		CLK100				=> CLK100,
		RESET					=> RESET,
		
		CNT_MAX_SET			=> tri_cnt_max,
		CNT_PHASE_SET		=> cnt_phase_0,
		CNT_PHASE_UD_SET	=> cnt_phase_UD_0,
		DT_SET				=> DT_SET,
		CNT_EN				=> CNT_EN,
		GATE_EN				=> GATE_EN,
		REF					=> MRate_UP1_C_4_int_hold,
		REF_hold_out		=> REF_hold_out_UP1,
		nGATE					=> nMMC_PWM1_FPGA_buf(3 downto 0),
		CNT_out				=> CNT_out_UP1
		
			);
			
-- ====================================================
-- MMC_PWM_FPGA_UP2
-- ====================================================

	MMCPWM_FPGA_UP2 : MMC_PWM_FPGA_SUB
		port map(
		CLK100				=> CLK100,
		RESET					=> RESET,
		
		CNT_MAX_SET			=> tri_cnt_max,
		CNT_PHASE_SET		=> cnt_phase_90,
		CNT_PHASE_UD_SET	=> cnt_phase_UD_90,
		DT_SET				=> DT_SET,
		CNT_EN				=> CNT_EN,
		GATE_EN				=> GATE_EN,
		REF					=> MRate_UP2_C_4_int_hold,
		REF_hold_out		=> REF_hold_out_UP2,
		nGATE					=> nMMC_PWM1_FPGA_buf(7 downto 4),
		CNT_out				=> CNT_out_UP2

			);


-- ====================================================
-- MMC_PWM_FPGA_UP3
-- ====================================================

	MMCPWM_FPGA_UP3 : MMC_PWM_FPGA_SUB
		port map(
		CLK100				=> CLK100,
		RESET					=> RESET,
		
		CNT_MAX_SET			=> tri_cnt_max,
		CNT_PHASE_SET		=> cnt_phase_180,
		CNT_PHASE_UD_SET	=> cnt_phase_UD_180,
		DT_SET				=> DT_SET,
		CNT_EN				=> CNT_EN,
		GATE_EN				=> GATE_EN,
		REF					=> MRate_UP3_C_4_int_hold,
		REF_hold_out		=> REF_hold_out_UP3,
		nGATE					=> nMMC_PWM1_FPGA_buf(11 downto 8),
		CNT_out				=> CNT_out_UP3

			);

-- ====================================================
-- MMC_PWM_FPGA_UP4
-- ====================================================

	MMCPWM_FPGA_UP4 : MMC_PWM_FPGA_SUB
		port map(
		CLK100				=> CLK100,
		RESET					=> RESET,
		
		CNT_MAX_SET			=> tri_cnt_max,
		CNT_PHASE_SET		=> cnt_phase_270,
		CNT_PHASE_UD_SET	=> cnt_phase_UD_270,
		DT_SET				=> DT_SET,
		CNT_EN				=> CNT_EN,
		GATE_EN				=> GATE_EN,
		REF					=> MRate_UP4_C_4_int_hold,
		REF_hold_out		=> REF_hold_out_UP4,
		nGATE					=> nMMC_PWM1_FPGA_buf(15 downto 12),
		CNT_out				=> CNT_out_UP4

			);

-- ====================================================
-- MMC_PWM_FPGA_UN1
-- ====================================================

	MMCPWM_FPGA_UN1 : MMC_PWM_FPGA_SUB
		port map(
		CLK100				=> CLK100,
		RESET					=> RESET,
		
		CNT_MAX_SET			=> tri_cnt_max,
		CNT_PHASE_SET		=> cnt_phase_45,
		CNT_PHASE_UD_SET	=> cnt_phase_UD_45,
		DT_SET				=> DT_SET,
		CNT_EN				=> CNT_EN,
		GATE_EN				=> GATE_EN,
		REF					=> MRate_UN1_C_4_int_hold,
		REF_hold_out		=> REF_hold_out_UN1,
		nGATE					=> nMMC_PWM1_FPGA_buf(19 downto 16),
		CNT_out				=> CNT_out_UN1

			);
			
-- ====================================================
-- MMC_PWM_FPGA_UN2
-- ====================================================

	MMCPWM_FPGA_UN2 : MMC_PWM_FPGA_SUB
		port map(
		CLK100				=> CLK100,
		RESET					=> RESET,
		
		CNT_MAX_SET			=> tri_cnt_max,
		CNT_PHASE_SET		=> cnt_phase_135,
		CNT_PHASE_UD_SET	=> cnt_phase_UD_135,
		DT_SET				=> DT_SET,
		CNT_EN				=> CNT_EN,
		GATE_EN				=> GATE_EN,
		REF					=> MRate_UN2_C_4_int_hold,
		REF_hold_out		=> REF_hold_out_UN2,
		nGATE					=> nMMC_PWM1_FPGA_buf(23 downto 20),
		CNT_out				=> CNT_out_UN2

			);


-- ====================================================
-- MMC_PWM_FPGA_UN3
-- ====================================================

	MMCPWM_FPGA_UN3 : MMC_PWM_FPGA_SUB
		port map(
		CLK100				=> CLK100,
		RESET					=> RESET,
		
		CNT_MAX_SET			=> tri_cnt_max,
		CNT_PHASE_SET		=> cnt_phase_225,
		CNT_PHASE_UD_SET	=> cnt_phase_UD_225,
		DT_SET				=> DT_SET,
		CNT_EN				=> CNT_EN,
		GATE_EN				=> GATE_EN,
		REF					=> MRate_UN3_C_4_int_hold,
		REF_hold_out		=> REF_hold_out_UN3,
		nGATE					=> nMMC_PWM1_FPGA_buf(27 downto 24),
		CNT_out				=> CNT_out_UN3

			);

-- ====================================================
-- MMC_PWM_FPGA_UN4
-- ====================================================

	MMCPWM_FPGA_UN4 : MMC_PWM_FPGA_SUB
		port map(
		CLK100				=> CLK100,
		RESET					=> RESET,
		
		CNT_MAX_SET			=> tri_cnt_max,
		CNT_PHASE_SET		=> cnt_phase_315,
		CNT_PHASE_UD_SET	=> cnt_phase_UD_315,
		DT_SET				=> DT_SET,
		CNT_EN				=> CNT_EN,
		GATE_EN				=> GATE_EN,
		REF					=> MRate_UN4_C_4_int_hold,
		REF_hold_out		=> REF_hold_out_UN4,
		nGATE					=> nMMC_PWM1_FPGA_buf(31 downto 28),
		CNT_out				=> CNT_out_UN4

			);



-- ====================================================
-- MMC_PWM_FPGA_VP1
-- ====================================================

	MMCPWM_FPGA_VP1 : MMC_PWM_FPGA_SUB
		port map(
		CLK100				=> CLK100,
		RESET					=> RESET,
		
		CNT_MAX_SET			=> tri_cnt_max,
		CNT_PHASE_SET		=> cnt_phase_0,
		CNT_PHASE_UD_SET	=> cnt_phase_UD_0,
		DT_SET				=> DT_SET,
		CNT_EN				=> CNT_EN,
		GATE_EN				=> GATE_EN,
		REF					=> MRate_VP1_C_4_int_hold,
		REF_hold_out		=> REF_hold_out_VP1,
		nGATE					=> nMMC_PWM1_FPGA_buf(35 downto 32),
		CNT_out				=> CNT_out_VP1

			);


-- ====================================================
-- MMC_PWM_FPGA_VP2
-- ====================================================

	MMCPWM_FPGA_VP2 : MMC_PWM_FPGA_SUB
		port map(
		CLK100				=> CLK100,
		RESET					=> RESET,
		
		CNT_MAX_SET			=> tri_cnt_max,
		CNT_PHASE_SET		=> cnt_phase_90,
		CNT_PHASE_UD_SET	=> cnt_phase_UD_90,
		DT_SET				=> DT_SET,
		CNT_EN				=> CNT_EN,
		GATE_EN				=> GATE_EN,
		REF					=> MRate_VP2_C_4_int_hold,
		REF_hold_out		=> REF_hold_out_VP2,
		nGATE					=> nMMC_PWM1_FPGA_buf(39 downto 36),
		CNT_out				=> CNT_out_VP2

			);


-- ====================================================
-- MMC_PWM_FPGA_VP3
-- ====================================================

	MMCPWM_FPGA_VP3 : MMC_PWM_FPGA_SUB
		port map(
		CLK100				=> CLK100,
		RESET					=> RESET,
		
		CNT_MAX_SET			=> tri_cnt_max,
		CNT_PHASE_SET		=> cnt_phase_180,
		CNT_PHASE_UD_SET	=> cnt_phase_UD_180,
		DT_SET				=> DT_SET,
		CNT_EN				=> CNT_EN,
		GATE_EN				=> GATE_EN,
		REF					=> MRate_VP3_C_4_int_hold,
		REF_hold_out		=> REF_hold_out_VP3,
		nGATE					=> nMMC_PWM1_FPGA_buf(43 downto 40),
		CNT_out				=> CNT_out_VP3

			);

-- ====================================================
-- MMC_PWM_FPGA_VP4
-- ====================================================

	MMCPWM_FPGA_VP4 : MMC_PWM_FPGA_SUB
		port map(
		CLK100				=> CLK100,
		RESET					=> RESET,
		
		CNT_MAX_SET			=> tri_cnt_max,
		CNT_PHASE_SET		=> cnt_phase_270,
		CNT_PHASE_UD_SET	=> cnt_phase_UD_270,
		DT_SET				=> DT_SET,
		CNT_EN				=> CNT_EN,
		GATE_EN				=> GATE_EN,
		REF					=> MRate_VP4_C_4_int_hold,
		REF_hold_out		=> REF_hold_out_VP4,
		nGATE					=> nMMC_PWM1_FPGA_buf(47 downto 44),
		CNT_out				=> CNT_out_VP4

			);

-- ====================================================
-- MMC_PWM_FPGA_VN1
-- ====================================================

	MMCPWM_FPGA_VN1 : MMC_PWM_FPGA_SUB
		port map(
		CLK100				=> CLK100,
		RESET					=> RESET,
		
		CNT_MAX_SET			=> tri_cnt_max,
		CNT_PHASE_SET		=> cnt_phase_45,
		CNT_PHASE_UD_SET	=> cnt_phase_UD_45,
		DT_SET				=> DT_SET,
		CNT_EN				=> CNT_EN,
		GATE_EN				=> GATE_EN,
		REF					=> MRate_VN1_C_4_int_hold,
		REF_hold_out		=> REF_hold_out_VN1,
		nGATE					=> nMMC_PWM2_FPGA_buf(3 downto 0),
		CNT_out				=> CNT_out_VN1

			);
			
-- ====================================================
-- MMC_PWM_FPGA_VN2
-- ====================================================

	MMCPWM_FPGA_VN2 : MMC_PWM_FPGA_SUB
		port map(
		CLK100				=> CLK100,
		RESET					=> RESET,
		
		CNT_MAX_SET			=> tri_cnt_max,
		CNT_PHASE_SET		=> cnt_phase_135,
		CNT_PHASE_UD_SET	=> cnt_phase_UD_135,
		DT_SET				=> DT_SET,
		CNT_EN				=> CNT_EN,
		GATE_EN				=> GATE_EN,
		REF					=> MRate_VN2_C_4_int_hold,
		REF_hold_out		=> REF_hold_out_VN2,
		nGATE					=> nMMC_PWM2_FPGA_buf(7 downto 4),
		CNT_out				=> CNT_out_VN2

			);


-- ====================================================
-- MMC_PWM_FPGA_VN3
-- ====================================================

	MMCPWM_FPGA_VN3 : MMC_PWM_FPGA_SUB
		port map(
		CLK100				=> CLK100,
		RESET					=> RESET,
		
		CNT_MAX_SET			=> tri_cnt_max,
		CNT_PHASE_SET		=> cnt_phase_225,
		CNT_PHASE_UD_SET	=> cnt_phase_UD_225,
		DT_SET				=> DT_SET,
		CNT_EN				=> CNT_EN,
		GATE_EN				=> GATE_EN,
		REF					=> MRate_VN3_C_4_int_hold,
		REF_hold_out		=> REF_hold_out_VN3,
		nGATE					=> nMMC_PWM2_FPGA_buf(11 downto 8),
		CNT_out				=> CNT_out_VN3

			);

-- ====================================================
-- MMC_PWM_FPGA_VN4
-- ====================================================

	MMCPWM_FPGA_VN4 : MMC_PWM_FPGA_SUB
		port map(
		CLK100				=> CLK100,
		RESET					=> RESET,
		
		CNT_MAX_SET			=> tri_cnt_max,
		CNT_PHASE_SET		=> cnt_phase_315,
		CNT_PHASE_UD_SET	=> cnt_phase_UD_315,
		DT_SET				=> DT_SET,
		CNT_EN				=> CNT_EN,
		GATE_EN				=> GATE_EN,
		REF					=> MRate_VN4_C_4_int_hold,
		REF_hold_out		=> REF_hold_out_VN4,
		nGATE					=> nMMC_PWM2_FPGA_buf(15 downto 12),
		CNT_out				=> CNT_out_VN4

			);

-- ====================================================
-- cnt_clk
-- ====================================================

	count : cnt_clk_3
		port map(
			CLK100 => CLK100,
			cnt_10k => cnt_10k_B,
			cnt_20k => cnt_20k_B,
			cnt_100k => cnt_100k_B,
			cnt_1M => cnt_1M_B,
			cnt_soft => cnt_soft_B,
			cnt_10k_ref => cnt_10k_ref_B,
			cnt_20k_ref => cnt_20k_ref_B,
			cnt_100k_ref => cnt_100k_ref_B,
			counter_1M => counter_1M_B,
			clk100k => clk100k_B,
			clk500k => clk500k_B,
			clk1M => clk1M_B,
			cnt_EN => cnt_EN,
			cnt_10k_0 => cnt_10k_0_B,
			cnt_10k_90 => cnt_10k_90_B ,
			cnt_10k_180 => cnt_10k_180_B,
			cnt_10k_270 => cnt_10k_270_B
			);
			
--追加ver2.5.6--
---- ====================================================
---- icon
---- ====================================================
--
--	icon_pwm_fpga_gen : icon1
--		port map(
--			CONTROL0			=>	CONTROL0
--			);
--
---- ====================================================
---- ila
---- ====================================================
--
--	ila_pwm_fpga_gen : ila1
--		port map(
--        CONTROL         =>	CONTROL0,
----			CLK		         =>	CLK100,  	--100Mクロック
----			CLK		         =>	CLK100k_B,  --100kクロック
--			CLK		         =>	CLK500k_B,  --500kクロック
--
--        TRIG0	         =>	REF_hold_out_UP1_vector,--制御量u(正規化あと_hold後:UP1)
----		  TRIG1	         =>	REF_UP1_vector,    		--制御量u(headspring:UP1)
----		  TRIG2	         =>	CNT_out_UP1_vector,     --三角波:UP1   
--		  TRIG1				=> VAu,                    --平均値制御出力(u相)          -- sfix16_En6
--		  TRIG2				=> vb_up1,                 --バランス制御出力(up1)          -- sfix16_En6
--		  TRIG3				=> Iz_UP_N_out_BB,         --Iup_10kHz                 -- sfix16_En9
--		  TRIG4				=> Voref_BB,         	   --Voref                     -- sfix16_En6
--		  TRIG5				=> dVoref_out_BB,         	--dVoref                    -- sfix16_En6
--		  TRIG6				=> Vo_N_out_BB,				--Vo                        -- sfix16_En6
----		  TRIG7				=> IC_N_out_BB,         	--IC                        -- sfix16_En9
--		  TRIG7				=> QMDB_OUTPUT_out_BB,     --QMDBout                   -- sfix16_En6
----		  TRIG9				=> Vdc_N_out_BB         	--Vdc_N                     -- sfix16_En6
--		  TRIG8				=> soft_cnt_out_BB,         --soft_cnt_out_BB,cnt80k_debug_BB           
--		  TRIG9	         =>	MRate_UP1,              --制御量u(正規化前:UP1)
--		  TRIG10(0)			=> clk10k_BBB         	   --clk10k_BBB                
--
--
--
------    TRIG1	         =>	MRate_UP1_C_4,
----      TRIG1	         =>	REF_hold_out_UP1_vector,--制御量u(正規化あと_hold後:UP1)
----  	  TRIG2	         =>	CNT_out_UP1_vector,     --三角波:UP1
----  	  TRIG3	         =>	REF_UP1_vector,         --制御量u(headspring:UP1)
----		  TRIG4				=> nMMC_PWM1_FPGA_buf2_UP1,--ゲート信号:UP1
----		  TRIG5				=> VAu,                    --平均値制御出力(u相)
----		  TRIG6				=> vb_up1,                 --バランス制御出力(up1)
------	  TRIG7				=> VA_no_limit_U           --VAと同じ
----		  TRIG7				=> Voref_BB,           --FPGA版Voref
----		  TRIG8				=> Iz_UP_N_debug_out_BB,   --Iup_debug_100MHz
----		  TRIG9				=> Iz_UP_N_out_BB          --Iup_debug_10kHz
--		  
----        TRIG4	         =>	MRate_UP1_int_2,
----        TRIG5	         =>	MRate_UN1_int_2,
----        TRIG6	         =>	MRate_VP1_int_2,
----        TRIG7	         =>	MRate_VN1_int_2,
----        TRIG8	         =>	MRate_UP1_C,
----        TRIG9	         =>	MRate_UN1_C,
----        TRIG10	         =>	MRate_VP1_C,
----        TRIG11	         =>	MRate_VN1_C,
----        TRIG12	         =>	MRate_UP1_C_2,
----        TRIG13	         =>	MRate_UN1_C_2,
----        TRIG14	         => MRate_VP1_C_2,	
----        TRIG15	         => MRate_VN1_C_2	
--		  );

-------------

	--単相
	--nMMC_PWM1_FPGA_buf2 <= nMMC_PWM1_FPGA_buf and "11111111" & "11111111" & "11111111" & "11111111" & "11111111" & "11111111" ; 
	--nMMC_PWM2_FPGA_buf2 <= nMMC_PWM2_FPGA_buf and "11111111" & "11111111" & "11111111" & "11111111" & "11111111" & "11111111" ; 
--三相モード--
	--単相モード(w相を0)--
	nMMC_PWM1_FPGA_buf2 <= nMMC_PWM1_FPGA_buf and "11111111" & "11111111" & "11111111" & "11111111" & "11111111" & "11111111" ; 
	nMMC_PWM2_FPGA_buf2 <= nMMC_PWM2_FPGA_buf and "00000000" & "00000000" & "00000000" & "00000000" & "11111111" & "11111111" ; 


	
	nMMC_PWM1_FPGA <= nMMC_PWM1_FPGA_buf2;
	nMMC_PWM2_FPGA <= nMMC_PWM2_FPGA_buf2; 			

	nMMC_PWM1_FPGA_buf2_UP1(0) <= nMMC_PWM1_FPGA_buf2(0);
	
	GATE_EN_out <= GATE_EN;

end Behavioral;






\end{mylisting}

\newpage
%Root main.tex
%---------------------------------------------------------------------------
%付録A　付録
%---------------------------------------------------------------------------

%\lstnewenvironment{mylisting}[1][]
%    {\lstset{
%        frame=single,
%        numbersep=10pt,
%        tabsize=2,
%        #1
%    }
%}{}
\begin{mylisting}[language=VHDL,caption=AD変換モジュール,label=appF_ADC_module_VHDL]

----------------------------------------------------------------------------------
-- FPGA Logic for MMC Controller Board
-- AD7357
--
--		Logic No.:		HSP0056-DA0-0001
--		Target Board:	HSP0056-B00-0001
--		Target FPGA:	XC6SLX75-2FGG676C
--
--		Version:		1.0 Build1
--		Date:			2016/09/05
--		Copyright:		Headspring Inc.
--
----------------------------------------------------------------------------------
library IEEE;
Library UNISIM;
use IEEE.STD_LOGIC_1164.ALL;
use IEEE.STD_LOGIC_ARITH.ALL;
use IEEE.STD_LOGIC_UNSIGNED.ALL;
use UNISIM.vcomponents.all;

entity AD7357 is
	Port (
		CLK150		: in	std_logic;
		CLK100		: in	std_logic;
		RESET		: in	std_logic;
		A		: in	std_logic_vector(7 downto 0);
		DW		: in	std_logic_vector(15 downto 0);
		DR		: out	std_logic_vector(15 downto 0)	:= "00000000" & "00000000";
		DR2	: out	std_logic_vector(15 downto 0)	:= "00000000" & "00000000";
		RD		: in	std_logic;
		WR		: in	std_logic;
		
		AD_CS		: out	std_logic							:= '0';
		AD_CLK		: out	std_logic							:= '1';
		AD_CH_EN	: out	std_logic_vector(9 downto 0)		:= "1111111111";

		ADC_GR0_CS		: in	std_logic;
		ADC_GR0_SCLK	: in	std_logic;
		ADC_GR0_DATA	: in	std_logic_vector(3 downto 0);

		ADC_GR1_CS		: in	std_logic;
		ADC_GR1_SCLK	: in	std_logic;
		ADC_GR1_DATA	: in	std_logic_vector(3 downto 0);

		ADC_GR2_CS		: in	std_logic;
		ADC_GR2_SCLK	: in	std_logic;
		ADC_GR2_DATA	: in	std_logic_vector(3 downto 0);

		ADC_GR3_CS		: in	std_logic;
		ADC_GR3_SCLK	: in	std_logic;
		ADC_GR3_DATA	: in	std_logic_vector(3 downto 0);

		ADC_GR4_CS		: in	std_logic;
		ADC_GR4_SCLK	: in	std_logic;
		ADC_GR4_DATA	: in	std_logic_vector(3 downto 0);

		ADC_GR5_CS		: in	std_logic;
		ADC_GR5_SCLK	: in	std_logic;
		ADC_GR5_DATA	: in	std_logic_vector(3 downto 0);

		ADC_GR6_CS		: in	std_logic;
		ADC_GR6_SCLK	: in	std_logic;
		ADC_GR6_DATA	: in	std_logic_vector(3 downto 0);

		ADC_GR7_CS		: in	std_logic;
		ADC_GR7_SCLK	: in	std_logic;
		ADC_GR7_DATA	: in	std_logic_vector(3 downto 0);

		ADC_GR8_CS		: in	std_logic;
		ADC_GR8_SCLK	: in	std_logic;
		ADC_GR8_DATA	: in	std_logic_vector(3 downto 0);

		ADC_GR9_CS		: in	std_logic;
		ADC_GR9_SCLK	: in	std_logic;
		ADC_GR9_DATA	: in	std_logic_vector(3 downto 0);
		
		ADC_ERR			: out	std_logic						:= '0';
		
		--追加ver2--
		Vc_UP1		: out	std_logic_vector(13 downto 0);		--コンデンサ電圧Vc
		Vc_UP2		: out	std_logic_vector(13 downto 0);
		Vc_UP3		: out	std_logic_vector(13 downto 0);
		Vc_UP4		: out	std_logic_vector(13 downto 0);
		Vc_UN1		: out	std_logic_vector(13 downto 0);
		Vc_UN2		: out	std_logic_vector(13 downto 0);
		Vc_UN3		: out	std_logic_vector(13 downto 0);
		Vc_UN4		: out	std_logic_vector(13 downto 0);
		Vc_VP1		: out	std_logic_vector(13 downto 0);
		Vc_VP2		: out	std_logic_vector(13 downto 0);
		Vc_VP3		: out	std_logic_vector(13 downto 0);
		Vc_VP4		: out	std_logic_vector(13 downto 0);
		Vc_VN1		: out	std_logic_vector(13 downto 0);
		Vc_VN2		: out	std_logic_vector(13 downto 0);
		Vc_VN3		: out	std_logic_vector(13 downto 0);
		Vc_VN4		: out	std_logic_vector(13 downto 0);
		Vc_WP1		: out	std_logic_vector(13 downto 0);
		Vc_WP2		: out	std_logic_vector(13 downto 0);
		Vc_WP3		: out	std_logic_vector(13 downto 0);
		Vc_WP4		: out	std_logic_vector(13 downto 0);
		Vc_WN1		: out	std_logic_vector(13 downto 0);
		Vc_WN2		: out	std_logic_vector(13 downto 0);
		Vc_WN3		: out	std_logic_vector(13 downto 0);
		Vc_WN4		: out	std_logic_vector(13 downto 0);
		Iz_UP			: out	std_logic_vector(13 downto 0);		--上下アーム電流Iz
		Iz_UN			: out	std_logic_vector(13 downto 0);
		Iz_VP			: out	std_logic_vector(13 downto 0);
		Iz_VN			: out	std_logic_vector(13 downto 0);
		Iz_WP			: out	std_logic_vector(13 downto 0);
		Iz_WN			: out	std_logic_vector(13 downto 0);
		Vdc			: out	std_logic_vector(13 downto 0);			--直流電源電圧Vdc
		Vuv_mmc		: out	std_logic_vector(13 downto 0);		--MMC出力線間電圧Vo_mmc
		Vwu_mmc		: out	std_logic_vector(13 downto 0);
		Vvw_mmc		: out	std_logic_vector(13 downto 0);
		Vwu_inv		: out	std_logic_vector(13 downto 0);		--インバータ出力線間電圧Vo_inv
		Vvw_inv		: out	std_logic_vector(13 downto 0);
		Vuv_inv		: out	std_logic_vector(13 downto 0);
		IU_inv		: out	std_logic_vector(13 downto 0);		--インバータ出力電流Io_inv
		IV_inv		: out	std_logic_vector(13 downto 0);
		IW_inv		: out	std_logic_vector(13 downto 0)
		-------		
		);
end AD7357;

architecture Behavioral of AD7357 is

	signal SPR				: std_logic_vector(15 downto 0)		:= "00000000" & "00000000";

	signal ADC0_DATA_B		: std_logic_vector(3 downto 0)		:= "0000";
	signal ADC1_DATA_B		: std_logic_vector(3 downto 0)		:= "0000";
	signal ADC2_DATA_B		: std_logic_vector(3 downto 0)		:= "0000";
	signal ADC3_DATA_B		: std_logic_vector(3 downto 0)		:= "0000";
	signal ADC4_DATA_B		: std_logic_vector(3 downto 0)		:= "0000";
	signal ADC5_DATA_B		: std_logic_vector(3 downto 0)		:= "0000";
	signal ADC6_DATA_B		: std_logic_vector(3 downto 0)		:= "0000";
	signal ADC7_DATA_B		: std_logic_vector(3 downto 0)		:= "0000";
	signal ADC8_DATA_B		: std_logic_vector(3 downto 0)		:= "0000";
	signal ADC9_DATA_B		: std_logic_vector(3 downto 0)		:= "0000";

	signal CNT_ADPHASE		: integer range 0 to 70				:= 0;

	signal AD0_CS_B			: std_logic							:= '0';
	signal AD0_CS_BB		: std_logic							:= '0';
	signal AD0_CS_BBB		: std_logic							:= '0';
	signal AD1_CS_B			: std_logic							:= '0';
	signal AD1_CS_BB		: std_logic							:= '0';
	signal AD1_CS_BBB		: std_logic							:= '0';
	signal AD2_CS_B			: std_logic							:= '0';
	signal AD2_CS_BB		: std_logic							:= '0';
	signal AD2_CS_BBB		: std_logic							:= '0';
	signal AD3_CS_B			: std_logic							:= '0';
	signal AD3_CS_BB		: std_logic							:= '0';
	signal AD3_CS_BBB		: std_logic							:= '0';
	signal AD4_CS_B			: std_logic							:= '0';
	signal AD4_CS_BB		: std_logic							:= '0';
	signal AD4_CS_BBB		: std_logic							:= '0';
	signal AD5_CS_B			: std_logic							:= '0';
	signal AD5_CS_BB		: std_logic							:= '0';
	signal AD5_CS_BBB		: std_logic							:= '0';
	signal AD6_CS_B			: std_logic							:= '0';
	signal AD6_CS_BB		: std_logic							:= '0';
	signal AD6_CS_BBB		: std_logic							:= '0';
	signal AD7_CS_B			: std_logic							:= '0';
	signal AD7_CS_BB		: std_logic							:= '0';
	signal AD7_CS_BBB		: std_logic							:= '0';
	signal AD8_CS_B			: std_logic							:= '0';
	signal AD8_CS_BB		: std_logic							:= '0';
	signal AD8_CS_BBB		: std_logic							:= '0';
	signal AD9_CS_B			: std_logic							:= '0';
	signal AD9_CS_BB		: std_logic							:= '0';
	signal AD9_CS_BBB		: std_logic							:= '0';

	signal AD0_SCLK_B		: std_logic							:= '0';
	signal AD0_SCLK_BB		: std_logic							:= '0';
	signal AD0_SCLK_BBB		: std_logic							:= '0';
	signal AD1_SCLK_B		: std_logic							:= '0';
	signal AD1_SCLK_BB		: std_logic							:= '0';
	signal AD1_SCLK_BBB		: std_logic							:= '0';
	signal AD2_SCLK_B		: std_logic							:= '0';
	signal AD2_SCLK_BB		: std_logic							:= '0';
	signal AD2_SCLK_BBB		: std_logic							:= '0';
	signal AD3_SCLK_B		: std_logic							:= '0';
	signal AD3_SCLK_BB		: std_logic							:= '0';
	signal AD3_SCLK_BBB		: std_logic							:= '0';
	signal AD4_SCLK_B		: std_logic							:= '0';
	signal AD4_SCLK_BB		: std_logic							:= '0';
	signal AD4_SCLK_BBB		: std_logic							:= '0';
	signal AD5_SCLK_B		: std_logic							:= '0';
	signal AD5_SCLK_BB		: std_logic							:= '0';
	signal AD5_SCLK_BBB		: std_logic							:= '0';
	signal AD6_SCLK_B		: std_logic							:= '0';
	signal AD6_SCLK_BB		: std_logic							:= '0';
	signal AD6_SCLK_BBB		: std_logic							:= '0';
	signal AD7_SCLK_B		: std_logic							:= '0';
	signal AD7_SCLK_BB		: std_logic							:= '0';
	signal AD7_SCLK_BBB		: std_logic							:= '0';
	signal AD8_SCLK_B		: std_logic							:= '0';
	signal AD8_SCLK_BB		: std_logic							:= '0';
	signal AD8_SCLK_BBB		: std_logic							:= '0';
	signal AD9_SCLK_B		: std_logic							:= '0';
	signal AD9_SCLK_BB		: std_logic							:= '0';
	signal AD9_SCLK_BBB		: std_logic							:= '0';

	signal CONV_CNT0		: integer range 0 to 16				:= 0;
	signal CONV_CNT1		: integer range 0 to 16				:= 0;
	signal CONV_CNT2		: integer range 0 to 16				:= 0;
	signal CONV_CNT3		: integer range 0 to 16				:= 0;
	signal CONV_CNT4		: integer range 0 to 16				:= 0;
	signal CONV_CNT5		: integer range 0 to 16				:= 0;
	signal CONV_CNT6		: integer range 0 to 16				:= 0;
	signal CONV_CNT7		: integer range 0 to 16				:= 0;
	signal CONV_CNT8		: integer range 0 to 16				:= 0;
	signal CONV_CNT9		: integer range 0 to 16				:= 0;

	signal AD0_SR0			: std_logic_vector(13 downto 0)		:= "000000" & "00000000";
	signal AD0_SR1			: std_logic_vector(13 downto 0)		:= "000000" & "00000000";
	signal AD0_SR2			: std_logic_vector(13 downto 0)		:= "000000" & "00000000";
	signal AD0_SR3			: std_logic_vector(13 downto 0)		:= "000000" & "00000000";
	signal AD1_SR0			: std_logic_vector(13 downto 0)		:= "000000" & "00000000";
	signal AD1_SR1			: std_logic_vector(13 downto 0)		:= "000000" & "00000000";
	signal AD1_SR2			: std_logic_vector(13 downto 0)		:= "000000" & "00000000";
	signal AD1_SR3			: std_logic_vector(13 downto 0)		:= "000000" & "00000000";
	signal AD2_SR0			: std_logic_vector(13 downto 0)		:= "000000" & "00000000";
	signal AD2_SR1			: std_logic_vector(13 downto 0)		:= "000000" & "00000000";
	signal AD2_SR2			: std_logic_vector(13 downto 0)		:= "000000" & "00000000";
	signal AD2_SR3			: std_logic_vector(13 downto 0)		:= "000000" & "00000000";
	signal AD3_SR0			: std_logic_vector(13 downto 0)		:= "000000" & "00000000";
	signal AD3_SR1			: std_logic_vector(13 downto 0)		:= "000000" & "00000000";
	signal AD3_SR2			: std_logic_vector(13 downto 0)		:= "000000" & "00000000";
	signal AD3_SR3			: std_logic_vector(13 downto 0)		:= "000000" & "00000000";
	signal AD4_SR0			: std_logic_vector(13 downto 0)		:= "000000" & "00000000";
	signal AD4_SR1			: std_logic_vector(13 downto 0)		:= "000000" & "00000000";
	signal AD4_SR2			: std_logic_vector(13 downto 0)		:= "000000" & "00000000";
	signal AD4_SR3			: std_logic_vector(13 downto 0)		:= "000000" & "00000000";
	signal AD5_SR0			: std_logic_vector(13 downto 0)		:= "000000" & "00000000";
	signal AD5_SR1			: std_logic_vector(13 downto 0)		:= "000000" & "00000000";
	signal AD5_SR2			: std_logic_vector(13 downto 0)		:= "000000" & "00000000";
	signal AD5_SR3			: std_logic_vector(13 downto 0)		:= "000000" & "00000000";
	signal AD6_SR0			: std_logic_vector(13 downto 0)		:= "000000" & "00000000";
	signal AD6_SR1			: std_logic_vector(13 downto 0)		:= "000000" & "00000000";
	signal AD6_SR2			: std_logic_vector(13 downto 0)		:= "000000" & "00000000";
	signal AD6_SR3			: std_logic_vector(13 downto 0)		:= "000000" & "00000000";
	signal AD7_SR0			: std_logic_vector(13 downto 0)		:= "000000" & "00000000";
	signal AD7_SR1			: std_logic_vector(13 downto 0)		:= "000000" & "00000000";
	signal AD7_SR2			: std_logic_vector(13 downto 0)		:= "000000" & "00000000";
	signal AD7_SR3			: std_logic_vector(13 downto 0)		:= "000000" & "00000000";
	signal AD8_SR0			: std_logic_vector(13 downto 0)		:= "000000" & "00000000";
	signal AD8_SR1			: std_logic_vector(13 downto 0)		:= "000000" & "00000000";
	signal AD8_SR2			: std_logic_vector(13 downto 0)		:= "000000" & "00000000";
	signal AD8_SR3			: std_logic_vector(13 downto 0)		:= "000000" & "00000000";
	signal AD9_SR0			: std_logic_vector(13 downto 0)		:= "000000" & "00000000";
	signal AD9_SR1			: std_logic_vector(13 downto 0)		:= "000000" & "00000000";
	signal AD9_SR2			: std_logic_vector(13 downto 0)		:= "000000" & "00000000";
	signal AD9_SR3			: std_logic_vector(13 downto 0)		:= "000000" & "00000000";

	signal ADC_GR0_0		: std_logic_vector(13 downto 0)		:= "000000" & "00000000";
	signal ADC_GR0_1		: std_logic_vector(13 downto 0)		:= "000000" & "00000000";
	signal ADC_GR0_2		: std_logic_vector(13 downto 0)		:= "000000" & "00000000";
	signal ADC_GR0_3		: std_logic_vector(13 downto 0)		:= "000000" & "00000000";
	signal ADC_GR1_0		: std_logic_vector(13 downto 0)		:= "000000" & "00000000";
	signal ADC_GR1_1		: std_logic_vector(13 downto 0)		:= "000000" & "00000000";
	signal ADC_GR1_2		: std_logic_vector(13 downto 0)		:= "000000" & "00000000";
	signal ADC_GR1_3		: std_logic_vector(13 downto 0)		:= "000000" & "00000000";
	signal ADC_GR2_0		: std_logic_vector(13 downto 0)		:= "000000" & "00000000";
	signal ADC_GR2_1		: std_logic_vector(13 downto 0)		:= "000000" & "00000000";
	signal ADC_GR2_2		: std_logic_vector(13 downto 0)		:= "000000" & "00000000";
	signal ADC_GR2_3		: std_logic_vector(13 downto 0)		:= "000000" & "00000000";
	signal ADC_GR3_0		: std_logic_vector(13 downto 0)		:= "000000" & "00000000";
	signal ADC_GR3_1		: std_logic_vector(13 downto 0)		:= "000000" & "00000000";
	signal ADC_GR3_2		: std_logic_vector(13 downto 0)		:= "000000" & "00000000";
	signal ADC_GR3_3		: std_logic_vector(13 downto 0)		:= "000000" & "00000000";
	signal ADC_GR4_0		: std_logic_vector(13 downto 0)		:= "000000" & "00000000";
	signal ADC_GR4_1		: std_logic_vector(13 downto 0)		:= "000000" & "00000000";
	signal ADC_GR4_2		: std_logic_vector(13 downto 0)		:= "000000" & "00000000";
	signal ADC_GR4_3		: std_logic_vector(13 downto 0)		:= "000000" & "00000000";
	signal ADC_GR5_0		: std_logic_vector(13 downto 0)		:= "000000" & "00000000";
	signal ADC_GR5_1		: std_logic_vector(13 downto 0)		:= "000000" & "00000000";
	signal ADC_GR5_2		: std_logic_vector(13 downto 0)		:= "000000" & "00000000";
	signal ADC_GR5_3		: std_logic_vector(13 downto 0)		:= "000000" & "00000000";
	signal ADC_GR6_0		: std_logic_vector(13 downto 0)		:= "000000" & "00000000";
	signal ADC_GR6_1		: std_logic_vector(13 downto 0)		:= "000000" & "00000000";
	signal ADC_GR6_2		: std_logic_vector(13 downto 0)		:= "000000" & "00000000";
	signal ADC_GR6_3		: std_logic_vector(13 downto 0)		:= "000000" & "00000000";
	signal ADC_GR7_0		: std_logic_vector(13 downto 0)		:= "000000" & "00000000";
	signal ADC_GR7_1		: std_logic_vector(13 downto 0)		:= "000000" & "00000000";
	signal ADC_GR7_2		: std_logic_vector(13 downto 0)		:= "000000" & "00000000";
	signal ADC_GR7_3		: std_logic_vector(13 downto 0)		:= "000000" & "00000000";
	signal ADC_GR8_0		: std_logic_vector(13 downto 0)		:= "000000" & "00000000";
	signal ADC_GR8_1		: std_logic_vector(13 downto 0)		:= "000000" & "00000000";
	signal ADC_GR8_2		: std_logic_vector(13 downto 0)		:= "000000" & "00000000";
	signal ADC_GR8_3		: std_logic_vector(13 downto 0)		:= "000000" & "00000000";
	signal ADC_GR9_0		: std_logic_vector(13 downto 0)		:= "000000" & "00000000";
	signal ADC_GR9_1		: std_logic_vector(13 downto 0)		:= "000000" & "00000000";
	signal ADC_GR9_2		: std_logic_vector(13 downto 0)		:= "000000" & "00000000";
	signal ADC_GR9_3		: std_logic_vector(13 downto 0)		:= "000000" & "00000000";

	signal UPDATE_EN		: std_logic							:= '1';
	signal CONV_OK			: std_logic							:= '0';
	signal CONV_OK_B		: std_logic							:= '0';
	signal CONV_OK_BB		: std_logic							:= '0';
	signal CONV_OK_BBB		: std_logic							:= '0';


	signal AD0_ERR			: std_logic							:= '0';
	signal AD1_ERR			: std_logic							:= '0';
	signal AD2_ERR			: std_logic							:= '0';
	signal AD3_ERR			: std_logic							:= '0';
	signal AD4_ERR			: std_logic							:= '0';
	signal AD5_ERR			: std_logic							:= '0';
	signal AD6_ERR			: std_logic							:= '0';
	signal AD7_ERR			: std_logic							:= '0';
	signal AD8_ERR			: std_logic							:= '0';
	signal AD9_ERR			: std_logic							:= '0';
	signal AD10_ERR			: std_logic							:= '0';
	signal AD11_ERR			: std_logic							:= '0';
	signal AD12_ERR			: std_logic							:= '0';
	signal AD13_ERR			: std_logic							:= '0';
	signal AD14_ERR			: std_logic							:= '0';
	signal AD15_ERR			: std_logic							:= '0';
	signal AD16_ERR			: std_logic							:= '0';
	signal AD17_ERR			: std_logic							:= '0';
	signal AD18_ERR			: std_logic							:= '0';
	signal AD19_ERR			: std_logic							:= '0';
	signal AD20_ERR			: std_logic							:= '0';
	signal AD21_ERR			: std_logic							:= '0';
	signal AD22_ERR			: std_logic							:= '0';
	signal AD23_ERR			: std_logic							:= '0';
	signal AD24_ERR			: std_logic							:= '0';
	signal AD25_ERR			: std_logic							:= '0';
	signal AD26_ERR			: std_logic							:= '0';
	signal AD27_ERR			: std_logic							:= '0';
	signal AD28_ERR			: std_logic							:= '0';
	signal AD29_ERR			: std_logic							:= '0';
	signal AD30_ERR			: std_logic							:= '0';
	signal AD31_ERR			: std_logic							:= '0';
	signal AD32_ERR			: std_logic							:= '0';
	signal AD33_ERR			: std_logic							:= '0';
	signal AD34_ERR			: std_logic							:= '0';
	signal AD35_ERR			: std_logic							:= '0';
	signal AD36_ERR			: std_logic							:= '0';
	signal AD37_ERR			: std_logic							:= '0';
	signal AD38_ERR			: std_logic							:= '0';
	signal AD39_ERR			: std_logic							:= '0';
	
	signal AD_ERR			: std_logic							:= '0';
	
	signal ERR_CLEAR		: std_logic							:= '0';

	signal AD0_DATA			: integer range 0 to 16383			:= 0;
	signal AD1_DATA			: integer range 0 to 16383			:= 0;
	signal AD2_DATA			: integer range 0 to 16383			:= 0;
	signal AD3_DATA			: integer range 0 to 16383			:= 0;
	signal AD4_DATA			: integer range 0 to 16383			:= 0;
	signal AD5_DATA			: integer range 0 to 16383			:= 0;
	signal AD6_DATA			: integer range 0 to 16383			:= 0;
	signal AD7_DATA			: integer range 0 to 16383			:= 0;
	signal AD8_DATA			: integer range 0 to 16383			:= 0;
	signal AD9_DATA			: integer range 0 to 16383			:= 0;
	signal AD10_DATA		: integer range 0 to 16383			:= 0;
	signal AD11_DATA		: integer range 0 to 16383			:= 0;
	signal AD12_DATA		: integer range 0 to 16383			:= 0;
	signal AD13_DATA		: integer range 0 to 16383			:= 0;
	signal AD14_DATA		: integer range 0 to 16383			:= 0;
	signal AD15_DATA		: integer range 0 to 16383			:= 0;
	signal AD16_DATA		: integer range 0 to 16383			:= 0;
	signal AD17_DATA		: integer range 0 to 16383			:= 0;
	signal AD18_DATA		: integer range 0 to 16383			:= 0;
	signal AD19_DATA		: integer range 0 to 16383			:= 0;
	signal AD20_DATA		: integer range 0 to 16383			:= 0;
	signal AD21_DATA		: integer range 0 to 16383			:= 0;
	signal AD22_DATA		: integer range 0 to 16383			:= 0;
	signal AD23_DATA		: integer range 0 to 16383			:= 0;
	signal AD24_DATA		: integer range 0 to 16383			:= 0;
	signal AD25_DATA		: integer range 0 to 16383			:= 0;
	signal AD26_DATA		: integer range 0 to 16383			:= 0;
	signal AD27_DATA		: integer range 0 to 16383			:= 0;
	signal AD28_DATA		: integer range 0 to 16383			:= 0;
	signal AD29_DATA		: integer range 0 to 16383			:= 0;
	signal AD30_DATA		: integer range 0 to 16383			:= 0;
	signal AD31_DATA		: integer range 0 to 16383			:= 0;
	signal AD32_DATA		: integer range 0 to 16383			:= 0;
	signal AD33_DATA		: integer range 0 to 16383			:= 0;
	signal AD34_DATA		: integer range 0 to 16383			:= 0;
	signal AD35_DATA		: integer range 0 to 16383			:= 0;
	signal AD36_DATA		: integer range 0 to 16383			:= 0;
	signal AD37_DATA		: integer range 0 to 16383			:= 0;
	signal AD38_DATA		: integer range 0 to 16383			:= 0;
	signal AD39_DATA		: integer range 0 to 16383			:= 0;
	signal AD0_DATA_B		: integer range 0 to 16383			:= 0;
	signal AD1_DATA_B		: integer range 0 to 16383			:= 0;
	signal AD2_DATA_B		: integer range 0 to 16383			:= 0;
	signal AD3_DATA_B		: integer range 0 to 16383			:= 0;
	signal AD4_DATA_B		: integer range 0 to 16383			:= 0;
	signal AD5_DATA_B		: integer range 0 to 16383			:= 0;
	signal AD6_DATA_B		: integer range 0 to 16383			:= 0;
	signal AD7_DATA_B		: integer range 0 to 16383			:= 0;
	signal AD8_DATA_B		: integer range 0 to 16383			:= 0;
	signal AD9_DATA_B		: integer range 0 to 16383			:= 0;
	signal AD10_DATA_B		: integer range 0 to 16383			:= 0;
	signal AD11_DATA_B		: integer range 0 to 16383			:= 0;
	signal AD12_DATA_B		: integer range 0 to 16383			:= 0;
	signal AD13_DATA_B		: integer range 0 to 16383			:= 0;
	signal AD14_DATA_B		: integer range 0 to 16383			:= 0;
	signal AD15_DATA_B		: integer range 0 to 16383			:= 0;
	signal AD16_DATA_B		: integer range 0 to 16383			:= 0;
	signal AD17_DATA_B		: integer range 0 to 16383			:= 0;
	signal AD18_DATA_B		: integer range 0 to 16383			:= 0;
	signal AD19_DATA_B		: integer range 0 to 16383			:= 0;
	signal AD20_DATA_B		: integer range 0 to 16383			:= 0;
	signal AD21_DATA_B		: integer range 0 to 16383			:= 0;
	signal AD22_DATA_B		: integer range 0 to 16383			:= 0;
	signal AD23_DATA_B		: integer range 0 to 16383			:= 0;
	signal AD24_DATA_B		: integer range 0 to 16383			:= 0;
	signal AD25_DATA_B		: integer range 0 to 16383			:= 0;
	signal AD26_DATA_B		: integer range 0 to 16383			:= 0;
	signal AD27_DATA_B		: integer range 0 to 16383			:= 0;
	signal AD28_DATA_B		: integer range 0 to 16383			:= 0;
	signal AD29_DATA_B		: integer range 0 to 16383			:= 0;
	signal AD30_DATA_B		: integer range 0 to 16383			:= 0;
	signal AD31_DATA_B		: integer range 0 to 16383			:= 0;
	signal AD32_DATA_B		: integer range 0 to 16383			:= 0;
	signal AD33_DATA_B		: integer range 0 to 16383			:= 0;
	signal AD34_DATA_B		: integer range 0 to 16383			:= 0;
	signal AD35_DATA_B		: integer range 0 to 16383			:= 0;
	signal AD36_DATA_B		: integer range 0 to 16383			:= 0;
	signal AD37_DATA_B		: integer range 0 to 16383			:= 0;
	signal AD38_DATA_B		: integer range 0 to 16383			:= 0;
	signal AD39_DATA_B		: integer range 0 to 16383			:= 0;
	signal AD0_OV_LEVEL		: integer range 0 to 16383			:= 16383;
	signal AD24_OV_LEVEL	: integer range 0 to 16383			:= 16383;
	signal AD25_OV_LEVEL	: integer range 0 to 16383			:= 16383;
	signal AD28_OV_LEVEL	: integer range 0 to 16383			:= 16383;
	signal AD31_OV_LEVEL	: integer range 0 to 16383			:= 16383;
	signal AD34_OV_LEVEL	: integer range 0 to 16383			:= 16383;
	signal AD37_OV_LEVEL	: integer range 0 to 16383			:= 16383;
	signal AD0_UV_LEVEL		: integer range 0 to 16383			:= 0;
	signal AD24_UV_LEVEL	: integer range 0 to 16383			:= 0;
	signal AD25_UV_LEVEL	: integer range 0 to 16383			:= 0;
	signal AD28_UV_LEVEL	: integer range 0 to 16383			:= 0;
	signal AD31_UV_LEVEL	: integer range 0 to 16383			:= 0;
	signal AD34_UV_LEVEL	: integer range 0 to 16383			:= 0;
	signal AD37_UV_LEVEL	: integer range 0 to 16383			:= 0;

	signal RESET_B		: std_logic							:= '0';
	signal RESET_200	: std_logic							:= '0';



begin

	process(CLK100)
	begin
		if CLK100'event and CLK100 = '1' then
			if RD = '1' then
				case A is
					when x"00"	=>	DR <= CONV_std_logic_vector(AD0_DATA_B,16);
					when x"01"	=>	DR <= CONV_std_logic_vector(AD1_DATA_B,16);
					when x"02"	=>	DR <= CONV_std_logic_vector(AD2_DATA_B,16);
					when x"03"	=>	DR <= CONV_std_logic_vector(AD3_DATA_B,16);
					when x"04"	=>	DR <= CONV_std_logic_vector(AD4_DATA_B,16);
					when x"05"	=>	DR <= CONV_std_logic_vector(AD5_DATA_B,16);
					when x"06"	=>	DR <= CONV_std_logic_vector(AD6_DATA_B,16);
					when x"07"	=>	DR <= CONV_std_logic_vector(AD7_DATA_B,16);
					when x"08"	=>	DR <= CONV_std_logic_vector(AD8_DATA_B,16);
					when x"09"	=>	DR <= CONV_std_logic_vector(AD9_DATA_B,16);
					when x"0A"	=>	DR <= CONV_std_logic_vector(AD10_DATA_B,16);
					when x"0B"	=>	DR <= CONV_std_logic_vector(AD11_DATA_B,16);
					when x"0C"	=>	DR <= CONV_std_logic_vector(AD12_DATA_B,16);
					when x"0D"	=>	DR <= CONV_std_logic_vector(AD13_DATA_B,16);
					when x"0E"	=>	DR <= CONV_std_logic_vector(AD14_DATA_B,16);
					when x"0F"	=>	DR <= CONV_std_logic_vector(AD15_DATA_B,16);
					when x"10"	=>	DR <= CONV_std_logic_vector(AD16_DATA_B,16);
					when x"11"	=>	DR <= CONV_std_logic_vector(AD17_DATA_B,16);
					when x"12"	=>	DR <= CONV_std_logic_vector(AD18_DATA_B,16);
					when x"13"	=>	DR <= CONV_std_logic_vector(AD19_DATA_B,16);
					when x"14"	=>	DR <= CONV_std_logic_vector(AD20_DATA_B,16);
					when x"15"	=>	DR <= CONV_std_logic_vector(AD21_DATA_B,16);
					when x"16"	=>	DR <= CONV_std_logic_vector(AD22_DATA_B,16);
					when x"17"	=>	DR <= CONV_std_logic_vector(AD23_DATA_B,16);
					when x"18"	=>	DR <= CONV_std_logic_vector(AD24_DATA_B,16);
					when x"19"	=>	DR <= CONV_std_logic_vector(AD25_DATA_B,16);
					when x"1A"	=>	DR <= CONV_std_logic_vector(AD26_DATA_B,16);
					when x"1B"	=>	DR <= CONV_std_logic_vector(AD27_DATA_B,16);
					when x"1C"	=>	DR <= CONV_std_logic_vector(AD28_DATA_B,16);
					when x"1D"	=>	DR <= CONV_std_logic_vector(AD29_DATA_B,16);
					when x"1E"	=>	DR <= CONV_std_logic_vector(AD30_DATA_B,16);
					when x"1F"	=>	DR <= CONV_std_logic_vector(AD31_DATA_B,16);
					when x"20"	=>	DR <= CONV_std_logic_vector(AD32_DATA_B,16);
					when x"21"	=>	DR <= CONV_std_logic_vector(AD33_DATA_B,16);
					when x"22"	=>	DR <= CONV_std_logic_vector(AD34_DATA_B,16);
					when x"23"	=>	DR <= CONV_std_logic_vector(AD35_DATA_B,16);
					when x"24"	=>	DR <= CONV_std_logic_vector(AD36_DATA_B,16);
					when x"25"	=>	DR <= CONV_std_logic_vector(AD37_DATA_B,16);
					when x"26"	=>	DR <= CONV_std_logic_vector(AD38_DATA_B,16);
					when x"27"	=>	DR <= CONV_std_logic_vector(AD39_DATA_B,16);
				
					when x"30"	=>	DR <= "00000000" & "0000000" & UPDATE_EN;
					when x"3F"	=>	DR <= SPR;
					
					when x"40"	=>	DR2 <= CONV_std_logic_vector(AD0_OV_LEVEL,16);
					when x"41"	=>	DR2 <= CONV_std_logic_vector(AD0_UV_LEVEL,16);
					when x"42"	=>	DR2 <= CONV_std_logic_vector(AD24_OV_LEVEL,16);
					when x"43"	=>	DR2 <= CONV_std_logic_vector(AD24_UV_LEVEL,16);
					when x"44"	=>	DR2 <= CONV_std_logic_vector(AD25_OV_LEVEL,16);
					when x"45"	=>	DR2 <= CONV_std_logic_vector(AD25_UV_LEVEL,16);
					when x"46"	=>	DR2 <= CONV_std_logic_vector(AD28_OV_LEVEL,16);
					when x"47"	=>	DR2 <= CONV_std_logic_vector(AD28_UV_LEVEL,16);
					when x"48"	=>	DR2 <= CONV_std_logic_vector(AD31_OV_LEVEL,16);
					when x"49"	=>	DR2 <= CONV_std_logic_vector(AD31_UV_LEVEL,16);
					when x"4A"	=>	DR2 <= CONV_std_logic_vector(AD34_OV_LEVEL,16);
					when x"4B"	=>	DR2 <= CONV_std_logic_vector(AD34_UV_LEVEL,16);
					when x"4C"	=>	DR2 <= CONV_std_logic_vector(AD37_OV_LEVEL,16);
					when x"4D"	=>	DR2 <= CONV_std_logic_vector(AD37_UV_LEVEL,16);
				
					when x"50"	=>	DR2 <= "00000000" & "0000000" & AD_ERR;
					when x"51"	=>	DR2 <= AD15_ERR & AD14_ERR & AD13_ERR & AD12_ERR & AD11_ERR & AD10_ERR & AD9_ERR & AD8_ERR 
											& AD7_ERR & AD6_ERR & AD5_ERR & AD4_ERR & AD3_ERR & AD2_ERR & AD1_ERR & AD0_ERR;
					when x"52"	=>	DR2 <= AD31_ERR & AD30_ERR & AD29_ERR & AD28_ERR & AD27_ERR & AD26_ERR & AD25_ERR & AD24_ERR 
											& AD23_ERR & AD22_ERR & AD21_ERR & AD20_ERR & AD19_ERR & AD18_ERR & AD17_ERR & AD16_ERR;
					when x"53"	=>	DR2 <= "00000000"
											& AD39_ERR & AD38_ERR & AD37_ERR & AD36_ERR & AD35_ERR & AD34_ERR & AD33_ERR & AD32_ERR;

					when x"FF"	=>	DR2 <= SPR;
					when others	=>	DR2 <= "00000000" & "00000000";
				end case;
			end if;

			if RESET = '1' then
				UPDATE_EN		<= '1';
				AD_CH_EN		<= "1111111111";
				SPR				<= "00000000" & "00000000";
				
				AD0_OV_LEVEL	<= 16383;
				AD24_OV_LEVEL	<= 16383;
				AD25_OV_LEVEL	<= 16383;
				AD28_OV_LEVEL	<= 16383;
				AD31_OV_LEVEL	<= 16383;
				AD34_OV_LEVEL	<= 16383;
				AD37_OV_LEVEL	<= 16383;
				AD0_UV_LEVEL	<= 0;
				AD24_UV_LEVEL	<= 0;
				AD25_UV_LEVEL	<= 0;
				AD28_UV_LEVEL	<= 0;
				AD31_UV_LEVEL	<= 0;
				AD34_UV_LEVEL	<= 0;
				AD37_UV_LEVEL	<= 0;
				
				ERR_CLEAR		<= '0';
			elsif WR = '1' then
				case A is
					when x"30"	=>	UPDATE_EN	<= DW(0);
					when x"31"	=>	AD_CH_EN	<= DW(9 downto 0);						-- only for debug use
					when x"40"	=>	AD0_OV_LEVEL	<= CONV_INTEGER(DW(13 downto 0));
					when x"41"	=>	AD0_UV_LEVEL	<= CONV_INTEGER(DW(13 downto 0));
					when x"42"	=>	AD24_OV_LEVEL	<= CONV_INTEGER(DW(13 downto 0));
					when x"43"	=>	AD24_UV_LEVEL	<= CONV_INTEGER(DW(13 downto 0));
					when x"44"	=>	AD25_OV_LEVEL	<= CONV_INTEGER(DW(13 downto 0));
					when x"45"	=>	AD25_UV_LEVEL	<= CONV_INTEGER(DW(13 downto 0));
					when x"46"	=>	AD28_OV_LEVEL	<= CONV_INTEGER(DW(13 downto 0));
					when x"47"	=>	AD28_UV_LEVEL	<= CONV_INTEGER(DW(13 downto 0));
					when x"48"	=>	AD31_OV_LEVEL	<= CONV_INTEGER(DW(13 downto 0));
					when x"49"	=>	AD31_UV_LEVEL	<= CONV_INTEGER(DW(13 downto 0));
					when x"4A"	=>	AD34_OV_LEVEL	<= CONV_INTEGER(DW(13 downto 0));
					when x"4B"	=>	AD34_UV_LEVEL	<= CONV_INTEGER(DW(13 downto 0));
					when x"4C"	=>	AD37_OV_LEVEL	<= CONV_INTEGER(DW(13 downto 0));
					when x"4D"	=>	AD37_UV_LEVEL	<= CONV_INTEGER(DW(13 downto 0));
					when x"50"	=>	ERR_CLEAR	<= DW(0);
					when x"FF"	=>	SPR			<= DW;
					when others		=> null;
				end case;
			else
				ERR_CLEAR		<= '0';
			end if;
			
			
			CONV_OK_B	<= CONV_OK;
			CONV_OK_BB	<= CONV_OK_B;
			CONV_OK_BBB	<= CONV_OK_BB;
			
			
			if RESET = '1' then
				AD0_DATA	<= 0;
				AD1_DATA	<= 0;
				AD2_DATA	<= 0;
				AD3_DATA	<= 0;
				AD4_DATA	<= 0;
				AD5_DATA	<= 0;
				AD6_DATA	<= 0;
				AD7_DATA	<= 0;
				AD8_DATA	<= 0;
				AD9_DATA	<= 0;
				AD10_DATA	<= 0;
				AD11_DATA	<= 0;
				AD12_DATA	<= 0;
				AD13_DATA	<= 0;
				AD14_DATA	<= 0;
				AD15_DATA	<= 0;
				AD16_DATA	<= 0;
				AD17_DATA	<= 0;
				AD18_DATA	<= 0;
				AD19_DATA	<= 0;
				AD20_DATA	<= 0;
				AD21_DATA	<= 0;
				AD22_DATA	<= 0;
				AD23_DATA	<= 0;
				AD24_DATA	<= 0;
				AD25_DATA	<= 0;
				AD26_DATA	<= 0;
				AD27_DATA	<= 0;
				AD28_DATA	<= 0;
				AD29_DATA	<= 0;
				AD30_DATA	<= 0;
				AD31_DATA	<= 0;
				AD32_DATA	<= 0;
				AD33_DATA	<= 0;
				AD34_DATA	<= 0;
				AD35_DATA	<= 0;
				AD36_DATA	<= 0;
				AD37_DATA	<= 0;
				AD38_DATA	<= 0;
				AD39_DATA	<= 0;
				
			elsif (CONV_OK_BBB = '0') and (CONV_OK_BB = '1') then
				
				AD0_DATA	<= CONV_INTEGER(ADC_GR0_0);
				AD1_DATA	<= CONV_INTEGER(ADC_GR0_1);
				AD2_DATA	<= CONV_INTEGER(ADC_GR0_2);
				AD3_DATA	<= CONV_INTEGER(ADC_GR0_3);
				AD4_DATA	<= CONV_INTEGER(ADC_GR1_0);
				AD5_DATA	<= CONV_INTEGER(ADC_GR1_1);
				AD6_DATA	<= CONV_INTEGER(ADC_GR1_2);
				AD7_DATA	<= CONV_INTEGER(ADC_GR1_3);
				AD8_DATA	<= CONV_INTEGER(ADC_GR2_0);
				AD9_DATA	<= CONV_INTEGER(ADC_GR2_1);
				AD10_DATA	<= CONV_INTEGER(ADC_GR2_2);
				AD11_DATA	<= CONV_INTEGER(ADC_GR2_3);
				AD12_DATA	<= CONV_INTEGER(ADC_GR3_0);
				AD13_DATA	<= CONV_INTEGER(ADC_GR3_1);
				AD14_DATA	<= CONV_INTEGER(ADC_GR3_2);
				AD15_DATA	<= CONV_INTEGER(ADC_GR3_3);
				AD16_DATA	<= CONV_INTEGER(ADC_GR4_0);
				AD17_DATA	<= CONV_INTEGER(ADC_GR4_1);
				AD18_DATA	<= CONV_INTEGER(ADC_GR4_2);
				AD19_DATA	<= CONV_INTEGER(ADC_GR4_3);
				AD20_DATA	<= CONV_INTEGER(ADC_GR5_0);
				AD21_DATA	<= CONV_INTEGER(ADC_GR5_1);
				AD22_DATA	<= CONV_INTEGER(ADC_GR5_2);
				AD23_DATA	<= CONV_INTEGER(ADC_GR5_3);
				AD24_DATA	<= CONV_INTEGER(ADC_GR6_0);
				AD25_DATA	<= CONV_INTEGER(ADC_GR6_1);
				AD26_DATA	<= CONV_INTEGER(ADC_GR6_2);
				AD27_DATA	<= CONV_INTEGER(ADC_GR6_3);
				AD28_DATA	<= CONV_INTEGER(ADC_GR7_0);
				AD29_DATA	<= CONV_INTEGER(ADC_GR7_1);
				AD30_DATA	<= CONV_INTEGER(ADC_GR7_2);
				AD31_DATA	<= CONV_INTEGER(ADC_GR7_3);
				AD32_DATA	<= CONV_INTEGER(ADC_GR8_0);
				AD33_DATA	<= CONV_INTEGER(ADC_GR8_1);
				AD34_DATA	<= CONV_INTEGER(ADC_GR8_2);
				AD35_DATA	<= CONV_INTEGER(ADC_GR8_3);
				AD36_DATA	<= CONV_INTEGER(ADC_GR9_0);
				AD37_DATA	<= CONV_INTEGER(ADC_GR9_1);
				AD38_DATA	<= CONV_INTEGER(ADC_GR9_2);
				AD39_DATA	<= CONV_INTEGER(ADC_GR9_3);
			end if;

			if RESET = '1' then
				AD0_ERR		<= '0';
				AD1_ERR		<= '0';
				AD2_ERR		<= '0';
				AD3_ERR		<= '0';
				AD4_ERR		<= '0';
				AD5_ERR		<= '0';
				AD6_ERR		<= '0';
				AD7_ERR		<= '0';
				AD8_ERR		<= '0';
				AD9_ERR		<= '0';
				AD10_ERR	<= '0';
				AD11_ERR	<= '0';
				AD12_ERR	<= '0';
				AD13_ERR	<= '0';
				AD14_ERR	<= '0';
				AD15_ERR	<= '0';
				AD16_ERR	<= '0';
				AD17_ERR	<= '0';
				AD18_ERR	<= '0';
				AD19_ERR	<= '0';
				AD20_ERR	<= '0';
				AD21_ERR	<= '0';
				AD22_ERR	<= '0';
				AD23_ERR	<= '0';
				AD24_ERR	<= '0';
				AD25_ERR	<= '0';
				AD26_ERR	<= '0';
				AD27_ERR	<= '0';
				AD28_ERR	<= '0';
				AD29_ERR	<= '0';
				AD30_ERR	<= '0';
				AD31_ERR	<= '0';
				AD32_ERR	<= '0';
				AD33_ERR	<= '0';
				AD34_ERR	<= '0';
				AD35_ERR	<= '0';
				AD36_ERR	<= '0';
				AD37_ERR	<= '0';
				AD38_ERR	<= '0';
				AD39_ERR	<= '0';
			
			else
				if (AD0_DATA > AD0_OV_LEVEL) or (AD0_DATA < AD0_UV_LEVEL) then
					AD0_ERR	<= '1';
				elsif ERR_CLEAR = '1' then
					AD0_ERR	<= '0';
				end if;
				
				if (AD1_DATA > AD0_OV_LEVEL) or (AD1_DATA < AD0_UV_LEVEL) then
					AD1_ERR	<= '1';
				elsif ERR_CLEAR = '1' then
					AD1_ERR	<= '0';
				end if;
				
				if (AD2_DATA > AD0_OV_LEVEL) or (AD2_DATA < AD0_UV_LEVEL) then
					AD2_ERR	<= '1';
				elsif ERR_CLEAR = '1' then
					AD2_ERR	<= '0';
				end if;
				
				if (AD3_DATA > AD0_OV_LEVEL) or (AD3_DATA < AD0_UV_LEVEL) then
					AD3_ERR	<= '1';
				elsif ERR_CLEAR = '1' then
					AD3_ERR	<= '0';
				end if;
				
				if (AD4_DATA > AD0_OV_LEVEL) or (AD4_DATA < AD0_UV_LEVEL) then
					AD4_ERR	<= '1';
				elsif ERR_CLEAR = '1' then
					AD4_ERR	<= '0';
				end if;
				
				if (AD5_DATA > AD0_OV_LEVEL) or (AD5_DATA < AD0_UV_LEVEL) then
					AD5_ERR	<= '1';
				elsif ERR_CLEAR = '1' then
					AD5_ERR	<= '0';
				end if;
				
				if (AD6_DATA > AD0_OV_LEVEL) or (AD6_DATA < AD0_UV_LEVEL) then
					AD6_ERR	<= '1';
				elsif ERR_CLEAR = '1' then
					AD6_ERR	<= '0';
				end if;
				
				if (AD7_DATA > AD0_OV_LEVEL) or (AD7_DATA < AD0_UV_LEVEL) then
					AD7_ERR	<= '1';
				elsif ERR_CLEAR = '1' then
					AD7_ERR	<= '0';
				end if;
				
				if (AD8_DATA > AD0_OV_LEVEL) or (AD8_DATA < AD0_UV_LEVEL) then
					AD8_ERR	<= '1';
				elsif ERR_CLEAR = '1' then
					AD8_ERR	<= '0';
				end if;
				
				if (AD9_DATA > AD0_OV_LEVEL) or (AD9_DATA < AD0_UV_LEVEL) then
					AD9_ERR	<= '1';
				elsif ERR_CLEAR = '1' then
					AD9_ERR	<= '0';
				end if;
				
				
				
				if (AD10_DATA > AD0_OV_LEVEL) or (AD10_DATA < AD0_UV_LEVEL) then
					AD10_ERR	<= '1';
				elsif ERR_CLEAR = '1' then
					AD10_ERR	<= '0';
				end if;
				
				if (AD11_DATA > AD0_OV_LEVEL) or (AD11_DATA < AD0_UV_LEVEL) then
					AD11_ERR	<= '1';
				elsif ERR_CLEAR = '1' then
					AD11_ERR	<= '0';
				end if;
				
				if (AD12_DATA > AD0_OV_LEVEL) or (AD12_DATA < AD0_UV_LEVEL) then
					AD12_ERR	<= '1';
				elsif ERR_CLEAR = '1' then
					AD12_ERR	<= '0';
				end if;
				
				if (AD13_DATA > AD0_OV_LEVEL) or (AD13_DATA < AD0_UV_LEVEL) then
					AD13_ERR	<= '1';
				elsif ERR_CLEAR = '1' then
					AD13_ERR	<= '0';
				end if;
				
				if (AD14_DATA > AD0_OV_LEVEL) or (AD14_DATA < AD0_UV_LEVEL) then
					AD14_ERR	<= '1';
				elsif ERR_CLEAR = '1' then
					AD14_ERR	<= '0';
				end if;
				
				if (AD15_DATA > AD0_OV_LEVEL) or (AD15_DATA < AD0_UV_LEVEL) then
					AD15_ERR	<= '1';
				elsif ERR_CLEAR = '1' then
					AD15_ERR	<= '0';
				end if;
				
				if (AD16_DATA > AD0_OV_LEVEL) or (AD16_DATA < AD0_UV_LEVEL) then
					AD16_ERR	<= '1';
				elsif ERR_CLEAR = '1' then
					AD16_ERR	<= '0';
				end if;
				
				if (AD17_DATA > AD0_OV_LEVEL) or (AD17_DATA < AD0_UV_LEVEL) then
					AD17_ERR	<= '1';
				elsif ERR_CLEAR = '1' then
					AD17_ERR	<= '0';
				end if;
				
				if (AD18_DATA > AD0_OV_LEVEL) or (AD18_DATA < AD0_UV_LEVEL) then
					AD18_ERR	<= '1';
				elsif ERR_CLEAR = '1' then
					AD18_ERR	<= '0';
				end if;
				
				if (AD19_DATA > AD0_OV_LEVEL) or (AD19_DATA < AD0_UV_LEVEL) then
					AD19_ERR	<= '1';
				elsif ERR_CLEAR = '1' then
					AD19_ERR	<= '0';
				end if;
				
				
				
				if (AD20_DATA > AD0_OV_LEVEL) or (AD20_DATA < AD0_UV_LEVEL) then
					AD20_ERR	<= '1';
				elsif ERR_CLEAR = '1' then
					AD20_ERR	<= '0';
				end if;
				
				if (AD21_DATA > AD0_OV_LEVEL) or (AD21_DATA < AD0_UV_LEVEL) then
					AD21_ERR	<= '1';
				elsif ERR_CLEAR = '1' then
					AD21_ERR	<= '0';
				end if;
				
				if (AD22_DATA > AD0_OV_LEVEL) or (AD22_DATA < AD0_UV_LEVEL) then
					AD22_ERR	<= '1';
				elsif ERR_CLEAR = '1' then
					AD22_ERR	<= '0';
				end if;
				
				if (AD23_DATA > AD0_OV_LEVEL) or (AD23_DATA < AD0_UV_LEVEL) then
					AD23_ERR	<= '1';
				elsif ERR_CLEAR = '1' then
					AD23_ERR	<= '0';
				end if;
				
				if (AD24_DATA > AD24_OV_LEVEL) or (AD24_DATA < AD24_UV_LEVEL) then
					AD24_ERR	<= '1';
				elsif ERR_CLEAR = '1' then
					AD24_ERR	<= '0';
				end if;
				
				if (AD25_DATA > AD25_OV_LEVEL) or (AD25_DATA < AD25_UV_LEVEL) then
					AD25_ERR	<= '1';
				elsif ERR_CLEAR = '1' then
					AD25_ERR	<= '0';
				end if;
				
				if (AD26_DATA > AD25_OV_LEVEL) or (AD26_DATA < AD25_UV_LEVEL) then
					AD26_ERR	<= '1';
				elsif ERR_CLEAR = '1' then
					AD26_ERR	<= '0';
				end if;
				
				if (AD27_DATA > AD25_OV_LEVEL) or (AD27_DATA < AD25_UV_LEVEL) then
					AD27_ERR	<= '1';
				elsif ERR_CLEAR = '1' then
					AD27_ERR	<= '0';
				end if;
				
				if (AD28_DATA > AD28_OV_LEVEL) or (AD28_DATA < AD28_UV_LEVEL) then
					AD28_ERR	<= '1';
				elsif ERR_CLEAR = '1' then
					AD28_ERR	<= '0';
				end if;
				
				if (AD29_DATA > AD28_OV_LEVEL) or (AD29_DATA < AD28_UV_LEVEL) then
					AD29_ERR	<= '1';
				elsif ERR_CLEAR = '1' then
					AD29_ERR	<= '0';
				end if;
				
				
				
				if (AD30_DATA > AD28_OV_LEVEL) or (AD30_DATA < AD28_UV_LEVEL) then
					AD30_ERR	<= '1';
				elsif ERR_CLEAR = '1' then
					AD30_ERR	<= '0';
				end if;
				
				if (AD31_DATA > AD31_OV_LEVEL) or (AD31_DATA < AD31_UV_LEVEL) then
					AD31_ERR	<= '1';
				elsif ERR_CLEAR = '1' then
					AD31_ERR	<= '0';
				end if;
				
				if (AD32_DATA > AD31_OV_LEVEL) or (AD32_DATA < AD31_UV_LEVEL) then
					AD32_ERR	<= '1';
				elsif ERR_CLEAR = '1' then
					AD32_ERR	<= '0';
				end if;
				
				if (AD33_DATA > AD31_OV_LEVEL) or (AD33_DATA < AD31_UV_LEVEL) then
					AD33_ERR	<= '1';
				elsif ERR_CLEAR = '1' then
					AD33_ERR	<= '0';
				end if;
				
				if (AD34_DATA > AD34_OV_LEVEL) or (AD34_DATA < AD34_UV_LEVEL) then
					AD34_ERR	<= '1';
				elsif ERR_CLEAR = '1' then
					AD34_ERR	<= '0';
				end if;
				
				if (AD35_DATA > AD34_OV_LEVEL) or (AD35_DATA < AD34_UV_LEVEL) then
					AD35_ERR	<= '1';
				elsif ERR_CLEAR = '1' then
					AD35_ERR	<= '0';
				end if;
				
				if (AD36_DATA > AD34_OV_LEVEL) or (AD36_DATA < AD34_UV_LEVEL) then
					AD36_ERR	<= '1';
				elsif ERR_CLEAR = '1' then
					AD36_ERR	<= '0';
				end if;
				
				if (AD37_DATA > AD37_OV_LEVEL) or (AD37_DATA < AD37_UV_LEVEL) then
					AD37_ERR	<= '1';
				elsif ERR_CLEAR = '1' then
					AD37_ERR	<= '0';
				end if;
				
				if (AD38_DATA > AD37_OV_LEVEL) or (AD38_DATA < AD37_UV_LEVEL) then
					AD38_ERR	<= '1';
				elsif ERR_CLEAR = '1' then
					AD38_ERR	<= '0';
				end if;
				
				if (AD39_DATA > AD37_OV_LEVEL) or (AD39_DATA < AD37_UV_LEVEL) then
					AD39_ERR	<= '1';
				elsif ERR_CLEAR = '1' then
					AD39_ERR	<= '0';
				end if;
				
			end if;
			
			if RESET = '1' then
				AD0_DATA_B	<= 0;
				AD1_DATA_B	<= 0;
				AD2_DATA_B	<= 0;
				AD3_DATA_B	<= 0;
				AD4_DATA_B	<= 0;
				AD5_DATA_B	<= 0;
				AD6_DATA_B	<= 0;
				AD7_DATA_B	<= 0;
				AD8_DATA_B	<= 0;
				AD9_DATA_B	<= 0;
				AD10_DATA_B	<= 0;
				AD11_DATA_B	<= 0;
				AD12_DATA_B	<= 0;
				AD13_DATA_B	<= 0;
				AD14_DATA_B	<= 0;
				AD15_DATA_B	<= 0;
				AD16_DATA_B	<= 0;
				AD17_DATA_B	<= 0;
				AD18_DATA_B	<= 0;
				AD19_DATA_B	<= 0;
				AD20_DATA_B	<= 0;
				AD21_DATA_B	<= 0;
				AD22_DATA_B	<= 0;
				AD23_DATA_B	<= 0;
				AD24_DATA_B	<= 0;
				AD25_DATA_B	<= 0;
				AD26_DATA_B	<= 0;
				AD27_DATA_B	<= 0;
				AD28_DATA_B	<= 0;
				AD29_DATA_B	<= 0;
				AD30_DATA_B	<= 0;
				AD31_DATA_B	<= 0;
				AD32_DATA_B	<= 0;
				AD33_DATA_B	<= 0;
				AD34_DATA_B	<= 0;
				AD35_DATA_B	<= 0;
				AD36_DATA_B	<= 0;
				AD37_DATA_B	<= 0;
				AD38_DATA_B	<= 0;
				AD39_DATA_B	<= 0;
			elsif (UPDATE_EN = '1') then
				AD0_DATA_B	<= AD0_DATA;
				AD1_DATA_B	<= AD1_DATA;
				AD2_DATA_B	<= AD2_DATA;
				AD3_DATA_B	<= AD3_DATA;
				AD4_DATA_B	<= AD4_DATA;
				AD5_DATA_B	<= AD5_DATA;
				AD6_DATA_B	<= AD6_DATA;
				AD7_DATA_B	<= AD7_DATA;
				AD8_DATA_B	<= AD8_DATA;
				AD9_DATA_B	<= AD9_DATA;
				AD10_DATA_B	<= AD10_DATA;
				AD11_DATA_B	<= AD11_DATA;
				AD12_DATA_B	<= AD12_DATA;
				AD13_DATA_B	<= AD13_DATA;
				AD14_DATA_B	<= AD14_DATA;
				AD15_DATA_B	<= AD15_DATA;
				AD16_DATA_B	<= AD16_DATA;
				AD17_DATA_B	<= AD17_DATA;
				AD18_DATA_B	<= AD18_DATA;
				AD19_DATA_B	<= AD19_DATA;
				AD20_DATA_B	<= AD20_DATA;
				AD21_DATA_B	<= AD21_DATA;
				AD22_DATA_B	<= AD22_DATA;
				AD23_DATA_B	<= AD23_DATA;
				AD24_DATA_B	<= AD24_DATA;
				AD25_DATA_B	<= AD25_DATA;
				AD26_DATA_B	<= AD26_DATA;
				AD27_DATA_B	<= AD27_DATA;
				AD28_DATA_B	<= AD28_DATA;
				AD29_DATA_B	<= AD29_DATA;
				AD30_DATA_B	<= AD30_DATA;
				AD31_DATA_B	<= AD31_DATA;
				AD32_DATA_B	<= AD32_DATA;
				AD33_DATA_B	<= AD33_DATA;
				AD34_DATA_B	<= AD34_DATA;
				AD35_DATA_B	<= AD35_DATA;
				AD36_DATA_B	<= AD36_DATA;
				AD37_DATA_B	<= AD37_DATA;
				AD38_DATA_B	<= AD38_DATA;
				AD39_DATA_B	<= AD39_DATA;
			end if;
		end if;
	end process;

	AD_ERR	<= '0' when AD0_ERR = '0' and AD1_ERR = '0' and AD2_ERR = '0' and AD3_ERR = '0' and AD4_ERR = '0' and AD5_ERR = '0' and AD6_ERR = '0' and AD7_ERR = '0' and 
					AD8_ERR = '0' and AD9_ERR = '0' and AD10_ERR = '0' and AD11_ERR = '0' and AD12_ERR = '0' and AD13_ERR = '0' and AD14_ERR = '0' and AD15_ERR = '0' and 
					AD16_ERR = '0' and AD17_ERR = '0' and AD18_ERR = '0' and AD19_ERR = '0' and AD20_ERR = '0' and AD21_ERR = '0' and AD22_ERR = '0' and AD23_ERR = '0' and 
					AD24_ERR = '0' and AD25_ERR = '0' and AD26_ERR = '0' and AD27_ERR = '0' and AD28_ERR = '0' and AD29_ERR = '0' and AD30_ERR = '0' and AD31_ERR = '0' and 
					AD32_ERR = '0' and AD33_ERR = '0' and AD34_ERR = '0' and AD35_ERR = '0' and AD36_ERR = '0' and AD37_ERR = '0' and AD38_ERR = '0' and AD39_ERR = '0'
				else '1';

	ADC_ERR	<= AD_ERR;

	-- ====================================================
	-- AD Conversion Control
	-- ====================================================
	process(CLK150)
	begin
		if CLK150'event and CLK150 = '1' then
			
			RESET_B		<= RESET;
			RESET_200	<= RESET_B;
			
			if RESET_200 = '1' then
				CNT_ADPHASE		<= 0;
				AD_CS			<= '0';
				AD_CLK			<= '1';
			else
				if CNT_ADPHASE = 70 then
					CNT_ADPHASE	<= 0;
				else
					CNT_ADPHASE	<= CNT_ADPHASE + 1;
				end if;
				
				case CNT_ADPHASE is
					when 0 => 
										AD_CS		<= '0';
										AD_CLK		<= '1';
					when 6 => 			AD_CS		<= '1';				-- アクイジション時間を40ns確保
					when 8 => 			AD_CLK		<= '0';				-- 1
					when 10 => 			AD_CLK		<= '1';
					when 12 => 			AD_CLK		<= '0';				-- 2
					when 14 => 			AD_CLK		<= '1';
					when 16 => 			AD_CLK		<= '0';				-- 3
					when 18 => 			AD_CLK		<= '1';
					when 20 => 			AD_CLK		<= '0';				-- 4
					when 22 => 			AD_CLK		<= '1';
					when 24 => 			AD_CLK		<= '0';				-- 5
					when 26 => 			AD_CLK		<= '1';
					when 28 => 			AD_CLK		<= '0';				-- 6
					when 30 => 			AD_CLK		<= '1';
					when 32 => 			AD_CLK		<= '0';				-- 7
					when 34 => 			AD_CLK		<= '1';
					when 36 => 			AD_CLK		<= '0';				-- 8
					when 38 => 			AD_CLK		<= '1';
					when 40 => 			AD_CLK		<= '0';				-- 9
					when 42 => 			AD_CLK		<= '1';
					when 44 => 			AD_CLK		<= '0';				-- 10
					when 46 => 			AD_CLK		<= '1';
					when 48 => 			AD_CLK		<= '0';				-- 11
					when 50 => 			AD_CLK		<= '1';
					when 52 => 			AD_CLK		<= '0';				-- 12
					when 54 => 			AD_CLK		<= '1';
					when 56 => 			AD_CLK		<= '0';				-- 13
					when 58 => 			AD_CLK		<= '1';
					when 60 => 			AD_CLK		<= '0';				-- 14
					when 62 => 			AD_CLK		<= '1';
					when 64 => 			AD_CLK		<= '0';				-- 15
					when 66 => 			AD_CLK		<= '1';
					when 68 => 			AD_CLK		<= '0';				-- 16
					when 70 => 			AD_CLK		<= '1';
					when others  => null;
				end case;
				
			end if;
			
			AD0_CS_B		<= ADC_GR0_CS;
			AD0_CS_BB		<= AD0_CS_B;
			AD0_CS_BBB		<= AD0_CS_BB;
			
			AD0_SCLK_B		<= ADC_GR0_SCLK;
			AD0_SCLK_BB		<= AD0_SCLK_B;
			AD0_SCLK_BBB	<= AD0_SCLK_BB;
			
			ADC0_DATA_B		<= ADC_GR0_DATA;
			
			if RESET_200 = '1' then
				CONV_CNT0	<= 0;
				
				AD0_SR0			<= "000000" & "00000000";
				AD0_SR1			<= "000000" & "00000000";
				AD0_SR2			<= "000000" & "00000000";
				AD0_SR3			<= "000000" & "00000000";
				
				ADC_GR0_0		<= "000000" & "00000000";
				ADC_GR0_1		<= "000000" & "00000000";
				ADC_GR0_2		<= "000000" & "00000000";
				ADC_GR0_3		<= "000000" & "00000000";

				--追加ver2--
				Vc_UP1			<= "000000" & "00000000";		--U相上アームコンデンサ電圧Vc
				Vc_UP2			<= "000000" & "00000000";
				Vc_UP3			<= "000000" & "00000000";
				Vc_UP4			<= "000000" & "00000000";
				-----------
			else
				
				if AD0_CS_BBB = '0' and AD0_CS_BB = '1' then
					CONV_CNT0	<= 0;
					CONV_OK		<= '0';
				elsif AD0_SCLK_BBB = '0' and AD0_SCLK_BB = '1' then
					AD0_SR0		<= AD0_SR0(12 downto 0) & ADC0_DATA_B(0);
					AD0_SR1		<= AD0_SR1(12 downto 0) & ADC0_DATA_B(1);
					AD0_SR2		<= AD0_SR2(12 downto 0) & ADC0_DATA_B(2);
					AD0_SR3		<= AD0_SR3(12 downto 0) & ADC0_DATA_B(3);
					CONV_CNT0	<= CONV_CNT0 + 1;
				elsif CONV_CNT0 = 15 then
					ADC_GR0_0	<= AD0_SR0;
					ADC_GR0_1	<= AD0_SR1;
					ADC_GR0_2	<= AD0_SR2;
					ADC_GR0_3	<= AD0_SR3;
					--追加ver2--
					Vc_UP1		<= AD0_SR0;		--U相上アームコンデンサ電圧Vc
					Vc_UP2		<= AD0_SR1;
					Vc_UP3		<= AD0_SR2;
					Vc_UP4		<= AD0_SR3;
					-----------
				elsif CONV_CNT0 = 16 then
					CONV_OK		<= '1';
					CONV_CNT0	<= 0;
				end if;
			end if;
			
			AD1_CS_B		<= ADC_GR1_CS;
			AD1_CS_BB		<= AD1_CS_B;
			AD1_CS_BBB		<= AD1_CS_BB;
			
			AD1_SCLK_B		<= ADC_GR1_SCLK;
			AD1_SCLK_BB		<= AD1_SCLK_B;
			AD1_SCLK_BBB	<= AD1_SCLK_BB;
			
			ADC1_DATA_B		<= ADC_GR1_DATA;
			
			if RESET_200 = '1' then
				CONV_CNT1	<= 0;
				
				AD1_SR0			<= "000000" & "00000000";
				AD1_SR1			<= "000000" & "00000000";
				AD1_SR2			<= "000000" & "00000000";
				AD1_SR3			<= "000000" & "00000000";
				
				ADC_GR1_0		<= "000000" & "00000000";
				ADC_GR1_1		<= "000000" & "00000000";
				ADC_GR1_2		<= "000000" & "00000000";
				ADC_GR1_3		<= "000000" & "00000000";
				
				--追加ver2--
				Vc_UN1			<= "000000" & "00000000";		--U相下アームコンデンサ電圧Vc
				Vc_UN2			<= "000000" & "00000000";
				Vc_UN3			<= "000000" & "00000000";
				Vc_UN4			<= "000000" & "00000000";
				-----------
			else
				
				if AD1_CS_BBB = '0' and AD1_CS_BB = '1' then
					CONV_CNT1	<= 0;
				elsif AD1_SCLK_BBB = '0' and AD1_SCLK_BB = '1' then
					AD1_SR0		<= AD1_SR0(12 downto 0) & ADC1_DATA_B(0);
					AD1_SR1		<= AD1_SR1(12 downto 0) & ADC1_DATA_B(1);
					AD1_SR2		<= AD1_SR2(12 downto 0) & ADC1_DATA_B(2);
					AD1_SR3		<= AD1_SR3(12 downto 0) & ADC1_DATA_B(3);
					CONV_CNT1	<= CONV_CNT1 + 1;
				elsif CONV_CNT1 = 15 then
					ADC_GR1_0	<= AD1_SR0;
					ADC_GR1_1	<= AD1_SR1;
					ADC_GR1_2	<= AD1_SR2;
					ADC_GR1_3	<= AD1_SR3;
					--追加ver2--
					Vc_UN1		<= AD1_SR0;		--U相下アームコンデンサ電圧Vc
					Vc_UN2		<= AD1_SR1;
					Vc_UN3		<= AD1_SR2;
					Vc_UN4		<= AD1_SR3;
					-----------
				elsif CONV_CNT1 = 16 then
					CONV_CNT1	<= 0;
				end if;
			end if;
			
			AD2_CS_B		<= ADC_GR2_CS;
			AD2_CS_BB		<= AD2_CS_B;
			AD2_CS_BBB		<= AD2_CS_BB;
			
			AD2_SCLK_B		<= ADC_GR2_SCLK;
			AD2_SCLK_BB		<= AD2_SCLK_B;
			AD2_SCLK_BBB	<= AD2_SCLK_BB;
			
			ADC2_DATA_B		<= ADC_GR2_DATA;
			
			if RESET_200 = '1' then
				CONV_CNT2	<= 0;
				
				AD2_SR0			<= "000000" & "00000000";
				AD2_SR1			<= "000000" & "00000000";
				AD2_SR2			<= "000000" & "00000000";
				AD2_SR3			<= "000000" & "00000000";
				
				ADC_GR2_0		<= "000000" & "00000000";
				ADC_GR2_1		<= "000000" & "00000000";
				ADC_GR2_2		<= "000000" & "00000000";
				ADC_GR2_3		<= "000000" & "00000000";
				
				--追加ver2--
				Vc_VP1			<= "000000" & "00000000";		--V相上アームコンデンサ電圧Vc
				Vc_VP2			<= "000000" & "00000000";
				Vc_VP3			<= "000000" & "00000000";
				Vc_VP4			<= "000000" & "00000000";
				-----------
			else
				
				if AD2_CS_BBB = '0' and AD2_CS_BB = '1' then
					CONV_CNT2	<= 0;
				elsif AD2_SCLK_BBB = '0' and AD2_SCLK_BB = '1' then
					AD2_SR0		<= AD2_SR0(12 downto 0) & ADC2_DATA_B(0);
					AD2_SR1		<= AD2_SR1(12 downto 0) & ADC2_DATA_B(1);
					AD2_SR2		<= AD2_SR2(12 downto 0) & ADC2_DATA_B(2);
					AD2_SR3		<= AD2_SR3(12 downto 0) & ADC2_DATA_B(3);
					CONV_CNT2	<= CONV_CNT2 + 1;
				elsif CONV_CNT2 = 15 then
					ADC_GR2_0	<= AD2_SR0;
					ADC_GR2_1	<= AD2_SR1;
					ADC_GR2_2	<= AD2_SR2;
					ADC_GR2_3	<= AD2_SR3;
					
					--追加ver2--
					Vc_VP1		<= AD2_SR0;		--V相上アームコンデンサ電圧Vc
					Vc_VP2		<= AD2_SR1;
					Vc_VP3		<= AD2_SR2;
					Vc_VP4		<= AD2_SR3;
					-----------
				elsif CONV_CNT2 = 16 then
					CONV_CNT2	<= 0;
				end if;
			end if;
			
			AD3_CS_B		<= ADC_GR3_CS;
			AD3_CS_BB		<= AD3_CS_B;
			AD3_CS_BBB		<= AD3_CS_BB;
			
			AD3_SCLK_B		<= ADC_GR3_SCLK;
			AD3_SCLK_BB		<= AD3_SCLK_B;
			AD3_SCLK_BBB	<= AD3_SCLK_BB;
			
			ADC3_DATA_B		<= ADC_GR3_DATA;
			
			if RESET_200 = '1' then
				CONV_CNT3	<= 0;
				
				AD3_SR0			<= "000000" & "00000000";
				AD3_SR1			<= "000000" & "00000000";
				AD3_SR2			<= "000000" & "00000000";
				AD3_SR3			<= "000000" & "00000000";
				
				ADC_GR3_0		<= "000000" & "00000000";
				ADC_GR3_1		<= "000000" & "00000000";
				ADC_GR3_2		<= "000000" & "00000000";
				ADC_GR3_3		<= "000000" & "00000000";
				
				--追加ver2--
				Vc_VN1			<= "000000" & "00000000";		--V相下アームコンデンサ電圧Vc
				Vc_VN2			<= "000000" & "00000000";
				Vc_VN3			<= "000000" & "00000000";
				Vc_VN4			<= "000000" & "00000000";
				-----------
			else
				
				if AD3_CS_BBB = '0' and AD3_CS_BB = '1' then
					CONV_CNT3	<= 0;
				elsif AD3_SCLK_BBB = '0' and AD3_SCLK_BB = '1' then
					AD3_SR0		<= AD3_SR0(12 downto 0) & ADC3_DATA_B(0);
					AD3_SR1		<= AD3_SR1(12 downto 0) & ADC3_DATA_B(1);
					AD3_SR2		<= AD3_SR2(12 downto 0) & ADC3_DATA_B(2);
					AD3_SR3		<= AD3_SR3(12 downto 0) & ADC3_DATA_B(3);
					CONV_CNT3	<= CONV_CNT3 + 1;
				elsif CONV_CNT3 = 15 then
					ADC_GR3_0	<= AD3_SR0;
					ADC_GR3_1	<= AD3_SR1;
					ADC_GR3_2	<= AD3_SR2;
					ADC_GR3_3	<= AD3_SR3;
					
					--追加ver2--
					Vc_VN1		<= AD3_SR0;		--V相下アームコンデンサ電圧Vc
					Vc_VN2		<= AD3_SR1;
					Vc_VN3		<= AD3_SR2;
					Vc_VN4		<= AD3_SR3;
					-----------
				elsif CONV_CNT3 = 16 then
					CONV_CNT3	<= 0;
				end if;
			end if;
			
			AD4_CS_B		<= ADC_GR4_CS;
			AD4_CS_BB		<= AD4_CS_B;
			AD4_CS_BBB		<= AD4_CS_BB;
			
			AD4_SCLK_B		<= ADC_GR4_SCLK;
			AD4_SCLK_BB		<= AD4_SCLK_B;
			AD4_SCLK_BBB	<= AD4_SCLK_BB;
			
			ADC4_DATA_B		<= ADC_GR4_DATA;
			
			if RESET_200 = '1' then
				CONV_CNT4	<= 0;
				
				AD4_SR0			<= "000000" & "00000000";
				AD4_SR1			<= "000000" & "00000000";
				AD4_SR2			<= "000000" & "00000000";
				AD4_SR3			<= "000000" & "00000000";
				
				ADC_GR4_0		<= "000000" & "00000000";
				ADC_GR4_1		<= "000000" & "00000000";
				ADC_GR4_2		<= "000000" & "00000000";
				ADC_GR4_3		<= "000000" & "00000000";
				
				--追加ver2--
				Vc_WP1			<= "000000" & "00000000";		--W相上アームコンデンサ電圧Vc
				Vc_WP2			<= "000000" & "00000000";
				Vc_WP3			<= "000000" & "00000000";
				Vc_WP4			<= "000000" & "00000000";
				-----------
			else
				
				if AD4_CS_BBB = '0' and AD4_CS_BB = '1' then
					CONV_CNT4	<= 0;
				elsif AD4_SCLK_BBB = '0' and AD4_SCLK_BB = '1' then
					AD4_SR0		<= AD4_SR0(12 downto 0) & ADC4_DATA_B(0);
					AD4_SR1		<= AD4_SR1(12 downto 0) & ADC4_DATA_B(1);
					AD4_SR2		<= AD4_SR2(12 downto 0) & ADC4_DATA_B(2);
					AD4_SR3		<= AD4_SR3(12 downto 0) & ADC4_DATA_B(3);
					CONV_CNT4	<= CONV_CNT4 + 1;
				elsif CONV_CNT4 = 15 then
					ADC_GR4_0	<= AD4_SR0;
					ADC_GR4_1	<= AD4_SR1;
					ADC_GR4_2	<= AD4_SR2;
					ADC_GR4_3	<= AD4_SR3;
					
					--追加ver2--
					Vc_WP1		<= AD4_SR0;		--W相上アームコンデンサ電圧Vc
					Vc_WP2		<= AD4_SR1;
					Vc_WP3		<= AD4_SR2;
					Vc_WP4		<= AD4_SR3;
					-----------
				elsif CONV_CNT4 = 16 then
					CONV_CNT4	<= 0;
				end if;
			end if;
			
			AD5_CS_B		<= ADC_GR5_CS;
			AD5_CS_BB		<= AD5_CS_B;
			AD5_CS_BBB		<= AD5_CS_BB;
			
			AD5_SCLK_B		<= ADC_GR5_SCLK;
			AD5_SCLK_BB		<= AD5_SCLK_B;
			AD5_SCLK_BBB	<= AD5_SCLK_BB;
			
			ADC5_DATA_B		<= ADC_GR5_DATA;
			
			if RESET_200 = '1' then
				CONV_CNT5	<= 0;
				
				AD5_SR0			<= "000000" & "00000000";
				AD5_SR1			<= "000000" & "00000000";
				AD5_SR2			<= "000000" & "00000000";
				AD5_SR3			<= "000000" & "00000000";
				
				ADC_GR5_0		<= "000000" & "00000000";
				ADC_GR5_1		<= "000000" & "00000000";
				ADC_GR5_2		<= "000000" & "00000000";
				ADC_GR5_3		<= "000000" & "00000000";
				
				--追加ver2--
				Vc_WN1			<= "000000" & "00000000";		--W相下アームコンデンサ電圧Vc
				Vc_WN2			<= "000000" & "00000000";
				Vc_WN3			<= "000000" & "00000000";
				Vc_WN4			<= "000000" & "00000000";
				-----------
			else
				
				if AD5_CS_BBB = '0' and AD5_CS_BB = '1' then
					CONV_CNT5	<= 0;
				elsif AD5_SCLK_BBB = '0' and AD5_SCLK_BB = '1' then
					AD5_SR0		<= AD5_SR0(12 downto 0) & ADC5_DATA_B(0);
					AD5_SR1		<= AD5_SR1(12 downto 0) & ADC5_DATA_B(1);
					AD5_SR2		<= AD5_SR2(12 downto 0) & ADC5_DATA_B(2);
					AD5_SR3		<= AD5_SR3(12 downto 0) & ADC5_DATA_B(3);
					CONV_CNT5	<= CONV_CNT5 + 1;
				elsif CONV_CNT5 = 15 then
					ADC_GR5_0	<= AD5_SR0;
					ADC_GR5_1	<= AD5_SR1;
					ADC_GR5_2	<= AD5_SR2;
					ADC_GR5_3	<= AD5_SR3;
					
					--追加ver2--
					Vc_WN1		<= AD5_SR0;		--W相下アームコンデンサ電圧Vc
					Vc_WN2		<= AD5_SR1;
					Vc_WN3		<= AD5_SR2;
					Vc_WN4		<= AD5_SR3;
					-----------
				elsif CONV_CNT5 = 16 then
					CONV_CNT5	<= 0;
				end if;
			end if;
			
			AD6_CS_B		<= ADC_GR6_CS;
			AD6_CS_BB		<= AD6_CS_B;
			AD6_CS_BBB		<= AD6_CS_BB;
			
			AD6_SCLK_B		<= ADC_GR6_SCLK;
			AD6_SCLK_BB		<= AD6_SCLK_B;
			AD6_SCLK_BBB	<= AD6_SCLK_BB;
			
			ADC6_DATA_B		<= ADC_GR6_DATA;
			
			if RESET_200 = '1' then
				CONV_CNT6	<= 0;
				
				AD6_SR0			<= "000000" & "00000000";
				AD6_SR1			<= "000000" & "00000000";
				AD6_SR2			<= "000000" & "00000000";
				AD6_SR3			<= "000000" & "00000000";
				
				ADC_GR6_0		<= "000000" & "00000000";
				ADC_GR6_1		<= "000000" & "00000000";
				ADC_GR6_2		<= "000000" & "00000000";
				ADC_GR6_3		<= "000000" & "00000000";
				
				--追加ver2--
				Vdc				<= "000000" & "00000000";		--直流電源電圧Vdc
				Vuv_mmc			<= "000000" & "00000000";		--MMC出力線間電圧Vo_mmc
				Vwu_mmc			<= "000000" & "00000000";
				Vvw_mmc			<= "000000" & "00000000";
				-----------
			else
				
				if AD6_CS_BBB = '0' and AD6_CS_BB = '1' then
					CONV_CNT6	<= 0;
				elsif AD6_SCLK_BBB = '0' and AD6_SCLK_BB = '1' then
					AD6_SR0		<= AD6_SR0(12 downto 0) & ADC6_DATA_B(0);
					AD6_SR1		<= AD6_SR1(12 downto 0) & ADC6_DATA_B(1);
					AD6_SR2		<= AD6_SR2(12 downto 0) & ADC6_DATA_B(2);
					AD6_SR3		<= AD6_SR3(12 downto 0) & ADC6_DATA_B(3);
					CONV_CNT6	<= CONV_CNT6 + 1;
				elsif CONV_CNT6 = 15 then
					ADC_GR6_0	<= AD6_SR0;
					ADC_GR6_1	<= AD6_SR1;
					ADC_GR6_2	<= AD6_SR2;
					ADC_GR6_3	<= AD6_SR3;
					
					--追加ver2--
					Vdc			<= AD6_SR0;		--直流電源電圧Vdc
					Vuv_mmc		<= AD6_SR1;		--MMC出力線間電圧Vo_mmc
					Vwu_mmc		<= AD6_SR2;
					Vvw_mmc		<= AD6_SR3;
					-----------
				elsif CONV_CNT6 = 16 then
					CONV_CNT6	<= 0;
				end if;
			end if;
			
			AD7_CS_B		<= ADC_GR7_CS;
			AD7_CS_BB		<= AD7_CS_B;
			AD7_CS_BBB		<= AD7_CS_BB;
			
			AD7_SCLK_B		<= ADC_GR7_SCLK;
			AD7_SCLK_BB		<= AD7_SCLK_B;
			AD7_SCLK_BBB	<= AD7_SCLK_BB;
			
			ADC7_DATA_B		<= ADC_GR7_DATA;
			
			if RESET_200 = '1' then
				CONV_CNT7	<= 0;
				
				AD7_SR0			<= "000000" & "00000000";
				AD7_SR1			<= "000000" & "00000000";
				AD7_SR2			<= "000000" & "00000000";
				AD7_SR3			<= "000000" & "00000000";
				
				ADC_GR7_0		<= "000000" & "00000000";
				ADC_GR7_1		<= "000000" & "00000000";
				ADC_GR7_2		<= "000000" & "00000000";
				ADC_GR7_3		<= "000000" & "00000000";
				
				--追加ver2--
				Vwu_inv			<= "000000" & "00000000";		--インバータ出力線間電圧Vo_inv
				Vvw_inv			<= "000000" & "00000000";
				Vuv_inv			<= "000000" & "00000000";
				Iz_UP				<= "000000" & "00000000";		--U相上アーム電流
				-----------
			else
				
				if AD7_CS_BBB = '0' and AD7_CS_BB = '1' then
					CONV_CNT7	<= 0;
				elsif AD7_SCLK_BBB = '0' and AD7_SCLK_BB = '1' then
					AD7_SR0		<= AD7_SR0(12 downto 0) & ADC7_DATA_B(0);
					AD7_SR1		<= AD7_SR1(12 downto 0) & ADC7_DATA_B(1);
					AD7_SR2		<= AD7_SR2(12 downto 0) & ADC7_DATA_B(2);
					AD7_SR3		<= AD7_SR3(12 downto 0) & ADC7_DATA_B(3);
					CONV_CNT7	<= CONV_CNT7 + 1;
				elsif CONV_CNT7 = 15 then
					ADC_GR7_0	<= AD7_SR0;
					ADC_GR7_1	<= AD7_SR1;
					ADC_GR7_2	<= AD7_SR2;
					ADC_GR7_3	<= AD7_SR3;
					
					--追加ver2--
					Vwu_inv		<= AD7_SR0;		--インバータ出力線間電圧Vo_inv
					Vvw_inv		<= AD7_SR1;		
					Vuv_inv		<= AD7_SR2;
					Iz_UP			<= AD7_SR3;		--U相上アーム電流
					-----------
				elsif CONV_CNT7 = 16 then
					CONV_CNT7	<= 0;
				end if;
			end if;
			
			AD8_CS_B		<= ADC_GR8_CS;
			AD8_CS_BB		<= AD8_CS_B;
			AD8_CS_BBB		<= AD8_CS_BB;
			
			AD8_SCLK_B		<= ADC_GR8_SCLK;
			AD8_SCLK_BB		<= AD8_SCLK_B;
			AD8_SCLK_BBB	<= AD8_SCLK_BB;
			
			ADC8_DATA_B		<= ADC_GR8_DATA;
			
			if RESET_200 = '1' then
				CONV_CNT8	<= 0;
				
				AD8_SR0			<= "000000" & "00000000";
				AD8_SR1			<= "000000" & "00000000";
				AD8_SR2			<= "000000" & "00000000";
				AD8_SR3			<= "000000" & "00000000";
				
				ADC_GR8_0		<= "000000" & "00000000";
				ADC_GR8_1		<= "000000" & "00000000";
				ADC_GR8_2		<= "000000" & "00000000";
				ADC_GR8_3		<= "000000" & "00000000";
				
				--追加ver2--
				Iz_VP				<= "000000" & "00000000";		--V相上アーム電流
				Iz_WP				<= "000000" & "00000000";		--W相上アーム電流
				Iz_UN				<= "000000" & "00000000";		--U相下アーム電流
				Iz_VN				<= "000000" & "00000000";		--V相下アーム電流
				-----------
			else
				
				if AD8_CS_BBB = '0' and AD8_CS_BB = '1' then
					CONV_CNT8	<= 0;
				elsif AD8_SCLK_BBB = '0' and AD8_SCLK_BB = '1' then
					AD8_SR0		<= AD8_SR0(12 downto 0) & ADC8_DATA_B(0);
					AD8_SR1		<= AD8_SR1(12 downto 0) & ADC8_DATA_B(1);
					AD8_SR2		<= AD8_SR2(12 downto 0) & ADC8_DATA_B(2);
					AD8_SR3		<= AD8_SR3(12 downto 0) & ADC8_DATA_B(3);
					CONV_CNT8	<= CONV_CNT8 + 1;
				elsif CONV_CNT8 = 15 then
					ADC_GR8_0	<= AD8_SR0;
					ADC_GR8_1	<= AD8_SR1;
					ADC_GR8_2	<= AD8_SR2;
					ADC_GR8_3	<= AD8_SR3;
					
					--追加ver2--
					Iz_VP			<= AD8_SR0;		--V相上アーム電流
					Iz_WP			<= AD8_SR1;		--W相上アーム電流
					Iz_UN			<= AD8_SR2;		--U相下アーム電流
					Iz_VN			<= AD8_SR3;		--V相下アーム電流
					-----------
				elsif CONV_CNT8 = 16 then
					CONV_CNT8	<= 0;
				end if;
			end if;
			
			AD9_CS_B		<= ADC_GR9_CS;
			AD9_CS_BB		<= AD9_CS_B;
			AD9_CS_BBB		<= AD9_CS_BB;
			
			AD9_SCLK_B		<= ADC_GR9_SCLK;
			AD9_SCLK_BB		<= AD9_SCLK_B;
			AD9_SCLK_BBB	<= AD9_SCLK_BB;
			
			ADC9_DATA_B		<= ADC_GR9_DATA;
			
			if RESET_200 = '1' then
				CONV_CNT9	<= 0;
				
				AD9_SR0			<= "000000" & "00000000";
				AD9_SR1			<= "000000" & "00000000";
				AD9_SR2			<= "000000" & "00000000";
				AD9_SR3			<= "000000" & "00000000";
				
				ADC_GR9_0		<= "000000" & "00000000";
				ADC_GR9_1		<= "000000" & "00000000";
				ADC_GR9_2		<= "000000" & "00000000";
				ADC_GR9_3		<= "000000" & "00000000";
				
				--追加ver2--
				Iz_WN				<= "000000" & "00000000";		--W相下アーム電流
				IU_inv			<= "000000" & "00000000";		--インバータ出力電流Io_inv
				IV_inv			<= "000000" & "00000000";
				IW_inv			<= "000000" & "00000000";
				-----------
			else
				
				if AD9_CS_BBB = '0' and AD9_CS_BB = '1' then
					CONV_CNT9	<= 0;
				elsif AD9_SCLK_BBB = '0' and AD9_SCLK_BB = '1' then
					AD9_SR0		<= AD9_SR0(12 downto 0) & ADC9_DATA_B(0);
					AD9_SR1		<= AD9_SR1(12 downto 0) & ADC9_DATA_B(1);
					AD9_SR2		<= AD9_SR2(12 downto 0) & ADC9_DATA_B(2);
					AD9_SR3		<= AD9_SR3(12 downto 0) & ADC9_DATA_B(3);
					CONV_CNT9	<= CONV_CNT9 + 1;
				elsif CONV_CNT9 = 15 then
					ADC_GR9_0	<= AD9_SR0;
					ADC_GR9_1	<= AD9_SR1;
					ADC_GR9_2	<= AD9_SR2;
					ADC_GR9_3	<= AD9_SR3;
					
					--追加ver2--
					Iz_WN			<= AD9_SR0;		--W相下アーム電流
					IU_inv		<= AD9_SR1;		--インバータ出力電流Io_inv
					IV_inv		<= AD9_SR2;		
					IW_inv		<= AD9_SR3;		
					-----------
				elsif CONV_CNT9 = 16 then
					CONV_CNT9	<= 0;
				end if;
			end if;

		end if;
	end process;


end Behavioral;



\end{mylisting}
%%Root main.tex
%---------------------------------------------------------------------------
%付録A　付録
%---------------------------------------------------------------------------
%
%\lstnewenvironment{mylisting}[1][]
%    {\lstset{
%        frame=single,
%        numbersep=10pt,
%        tabsize=2,
%        #1
%    }
%}{}
\begin{mylisting}[language=VHDL,caption=WDTモジュール,label=appF_WDT_module_VHDL]

----------------------------------------------------------------------------------
-- FPGA Logic for MMC Controller Board
-- MMC WDT Logic
--
--		Logic No.:		HSP0056-DA0-0001
--		Target Board:	HSP0056-B00-0001
--		Target FPGA:	XC6SLX75-2FGG676C
--
--		Version:		1.0 Build1
--		Date:			2016/09/05
--		Copyright:		Headspring Inc.
--
----------------------------------------------------------------------------------
library IEEE;
Library UNISIM;
use IEEE.STD_LOGIC_1164.ALL;
use IEEE.STD_LOGIC_ARITH.ALL;
use IEEE.STD_LOGIC_SIGNED.ALL;
use UNISIM.vcomponents.all;

entity MMC_WDT is
	Port (
		CLK100		: in	std_logic;
		RESET		: in	std_logic;
		A			: in	std_logic_vector(7 downto 0);
		DW			: in	std_logic_vector(15 downto 0);
		DR			: out	std_logic_vector(15 downto 0)	:= "00000000" & "00000000";
		RD			: in	std_logic;
		WR			: in	std_logic;
		
		WDT_ON		: out	std_logic						:= '0';
		ERROR		: out	std_logic						:= '0'
		);
end MMC_WDT;


architecture Behavioral of MMC_WDT is

	signal SPR					: std_logic_vector(15 downto 0)		:= "00000000" & "00000000";
	signal WDT_OV_CYCLE			: std_logic_vector(15 downto 0)		:= "00000000" & "00000000";
	signal WDT_EN				: std_logic							:= '0';
	signal WDT_ERR				: std_logic							:= '0';
	signal WDT_CLEAR			: std_logic							:= '0';
	signal PRE_CNT				: integer range 0 to 99				:= 99;
	signal WDT_CNT				: std_logic_vector(15 downto 0)		:= "11111111" & "11111111";



begin
	process(CLK100)
	begin
		if CLK100'event and CLK100 = '1' then
			if RD = '1' then
				case A is
					when x"00"	=>	DR <= WDT_OV_CYCLE;
					when x"01"	=>	DR <= "00000000" & "0000000" & WDT_EN;
					when x"03"	=>	DR <= "00000000" & "0000000" & WDT_ERR;
				
					when x"FF"	=>	DR <= SPR;
					when others	=>	DR <= "00000000" & "00000000";
				end case;
			end if;

			if RESET = '1' then
				SPR				<= "00000000" & "00000000";
				WDT_CLEAR		<= '0';
				WDT_OV_CYCLE	<= "00000000" & "00000000";
				WDT_EN			<= '0';
			elsif WR = '1' then
				case A is
					when x"00"	=>	if WDT_EN = '0' then						-- 一度 WDT_ENを'1'にすると書き換え不可能
										WDT_OV_CYCLE		<= DW;
									end if;
					when x"01"	=>	if WDT_EN = '0' then						-- 一度 WDT_ENを'1'にすると書き換え不可能
										WDT_EN				<= DW(0);
									end if;
					when x"02"	=>	if DW = x"AA55" then
										WDT_CLEAR <= '1';
									end if;
				
					when x"FF"	=>	SPR			<= DW;
					when others		=> null;
				end case;
			else
				WDT_CLEAR <= '0';
			end if;
			
			
			if RESET = '1' then
				PRE_CNT		<= 99;
				WDT_CNT		<= "11111111" & "11111111";
				WDT_ERR		<= '0';
			elsif WDT_EN = '0' then
				PRE_CNT		<= 99;
				WDT_CNT		<= WDT_OV_CYCLE;
			else
				if WDT_CNT = "00000000" & "00000000" then						-- エラー状態になると解除不能になる
					WDT_ERR		<= '1';
				elsif WDT_CLEAR = '1' then
					PRE_CNT		<= 99;
					WDT_CNT		<= WDT_OV_CYCLE;
				elsif PRE_CNT = 0 then
					PRE_CNT		<= 99;
					WDT_CNT		<= WDT_CNT - '1';
				else
					PRE_CNT	<= PRE_CNT - 1;
				end if;
			end if;
			
			
			
		end if;
	end process;
	
	ERROR	<= WDT_ERR;
	WDT_ON	<= WDT_EN;


end Behavioral;



\end{mylisting}
\newpage
%Root main.tex
%---------------------------------------------------------------------------
%付録A　付録
%---------------------------------------------------------------------------

%\lstnewenvironment{mylisting}[1][]
%    {\lstset{
%        frame=single,
%        numbersep=10pt,
%        tabsize=2,
%        #1
%    }
%}{}
\begin{mylisting}[language=VHDL,caption=MSSRDCDB制御\&SVPWMモジュール,label=appF_Control_top2_module_VHDL]

-- -------------------------------------------------------------
-- 
-- File Name: hdl_prj\hdlsrc\MMC_9level_half_cel\Control_top2.vhd
-- Created: 2024-01-24 16:04:30
-- 
-- Generated by MATLAB 9.12 and HDL Coder 3.20
-- 
-- 
-- -------------------------------------------------------------
-- Rate and Clocking Details
-- -------------------------------------------------------------
-- Model base rate: 1e-08
-- Target subsystem base rate: 1e-08
-- Explicit user oversample request: 100x
-- 
-- 
-- Clock Enable  Sample Time
-- -------------------------------------------------------------
-- ce_out        1e-06
-- -------------------------------------------------------------
-- 
-- 
-- Output Signal                 Clock Enable  Sample Time
-- -------------------------------------------------------------
-- MMC_PWM_OUT1                  ce_out        1e-06
-- MMC_PWM_OUT2                  ce_out        1e-06
-- Vuv                           ce_out        1e-06
-- Vvw                           ce_out        1e-06
-- Ioutu                         ce_out        1e-06
-- Ioutv                         ce_out        1e-06
-- ILu                           ce_out        1e-06
-- ILv                           ce_out        1e-06
-- Vrefa                         ce_out        1e-06
-- Vrefb                         ce_out        1e-06
-- deltaTa                       ce_out        1e-06
-- deltaTb                       ce_out        1e-06
-- deltaTa_hold                  ce_out        1e-06
-- deltaTb_hold                  ce_out        1e-06
-- switch_control1               ce_out        1e-06
-- -------------------------------------------------------------
-- 
-- -------------------------------------------------------------


-- -------------------------------------------------------------
-- 
-- Module: Control_top2
-- Source Path: MMC_9level_half_cel/Control/Control_top2
-- Hierarchy Level: 0
-- 
-- -------------------------------------------------------------
LIBRARY IEEE;
USE IEEE.std_logic_1164.ALL;
USE IEEE.numeric_std.ALL;
USE work.Control_top2_pkg.ALL;

ENTITY Control_top2 IS
  PORT( clk                               :   IN    std_logic;
        reset                             :   IN    std_logic;
        clk_enable                        :   IN    std_logic;
        V_uv                              :   IN    std_logic_vector(13 DOWNTO 0);  -- ufix14
        V_uw                              :   IN    std_logic_vector(13 DOWNTO 0);  -- ufix14
        iout_u                            :   IN    std_logic_vector(13 DOWNTO 0);  -- ufix14
        iout_v                            :   IN    std_logic_vector(13 DOWNTO 0);  -- ufix14
        iL_u                              :   IN    std_logic_vector(13 DOWNTO 0);  -- ufix14
        iL_v                              :   IN    std_logic_vector(13 DOWNTO 0);  -- ufix14
        switch_control                    :   IN    std_logic_vector(13 DOWNTO 0);  -- ufix14
        Up_current                        :   IN    std_logic_vector(13 DOWNTO 0);  -- ufix14
        Un_current                        :   IN    std_logic_vector(13 DOWNTO 0);  -- ufix14
        Vp_current                        :   IN    std_logic_vector(13 DOWNTO 0);  -- ufix14
        Vn_current                        :   IN    std_logic_vector(13 DOWNTO 0);  -- ufix14
        Wp_current                        :   IN    std_logic_vector(13 DOWNTO 0);  -- ufix14
        UW_current                        :   IN    std_logic_vector(13 DOWNTO 0);  -- ufix14
        CV_U1                             :   IN    std_logic_vector(13 DOWNTO 0);  -- ufix14
        CV_U2                             :   IN    std_logic_vector(13 DOWNTO 0);  -- ufix14
        CV_U3                             :   IN    std_logic_vector(13 DOWNTO 0);  -- ufix14
        CV_U4                             :   IN    std_logic_vector(13 DOWNTO 0);  -- ufix14
        CV_U5                             :   IN    std_logic_vector(13 DOWNTO 0);  -- ufix14
        CV_U6                             :   IN    std_logic_vector(13 DOWNTO 0);  -- ufix14
        CV_U7                             :   IN    std_logic_vector(13 DOWNTO 0);  -- ufix14
        CV_U8                             :   IN    std_logic_vector(13 DOWNTO 0);  -- ufix14
        CV_V1                             :   IN    std_logic_vector(13 DOWNTO 0);  -- ufix14
        CV_V2                             :   IN    std_logic_vector(13 DOWNTO 0);  -- ufix14
        CV_V3                             :   IN    std_logic_vector(13 DOWNTO 0);  -- ufix14
        CV_V4                             :   IN    std_logic_vector(13 DOWNTO 0);  -- ufix14
        CV_V5                             :   IN    std_logic_vector(13 DOWNTO 0);  -- ufix14
        CV_V6                             :   IN    std_logic_vector(13 DOWNTO 0);  -- ufix14
        CV_V7                             :   IN    std_logic_vector(13 DOWNTO 0);  -- ufix14
        CV_V8                             :   IN    std_logic_vector(13 DOWNTO 0);  -- ufix14
        CV_W1                             :   IN    std_logic_vector(13 DOWNTO 0);  -- ufix14
        CV_W2                             :   IN    std_logic_vector(13 DOWNTO 0);  -- ufix14
        CV_W3                             :   IN    std_logic_vector(13 DOWNTO 0);  -- ufix14
        CV_W4                             :   IN    std_logic_vector(13 DOWNTO 0);  -- ufix14
        CV_W5                             :   IN    std_logic_vector(13 DOWNTO 0);  -- ufix14
        CV_W6                             :   IN    std_logic_vector(13 DOWNTO 0);  -- ufix14
        CV_W7                             :   IN    std_logic_vector(13 DOWNTO 0);  -- ufix14
        CV_W8                             :   IN    std_logic_vector(13 DOWNTO 0);  -- ufix14
        ce_out                            :   OUT   std_logic;
        MMC_PWM_OUT1                      :   OUT   std_logic_vector(47 DOWNTO 0);  -- ufix48
        MMC_PWM_OUT2                      :   OUT   std_logic_vector(47 DOWNTO 0);  -- ufix48
        Vuv                               :   OUT   std_logic_vector(19 DOWNTO 0);  -- sfix20
        Vvw                               :   OUT   std_logic_vector(19 DOWNTO 0);  -- sfix20
        Ioutu                             :   OUT   std_logic_vector(19 DOWNTO 0);  -- sfix20
        Ioutv                             :   OUT   std_logic_vector(19 DOWNTO 0);  -- sfix20
        ILu                               :   OUT   std_logic_vector(19 DOWNTO 0);  -- sfix20
        ILv                               :   OUT   std_logic_vector(19 DOWNTO 0);  -- sfix20
        Vrefa                             :   OUT   std_logic_vector(19 DOWNTO 0);  -- sfix20
        Vrefb                             :   OUT   std_logic_vector(19 DOWNTO 0);  -- sfix20
        deltaTa                           :   OUT   std_logic_vector(15 DOWNTO 0);  -- int16
        deltaTb                           :   OUT   std_logic_vector(15 DOWNTO 0);  -- int16
        deltaTa_hold                      :   OUT   std_logic_vector(15 DOWNTO 0);  -- int16
        deltaTb_hold                      :   OUT   std_logic_vector(15 DOWNTO 0);  -- int16
        switch_control1                   :   OUT   std_logic_vector(13 DOWNTO 0)  -- ufix14
        );
END Control_top2;


ARCHITECTURE rtl OF Control_top2 IS

  -- Component Declarations
  COMPONENT Control_top2_tc
    PORT( clk                             :   IN    std_logic;
          reset                           :   IN    std_logic;
          clk_enable                      :   IN    std_logic;
          enb                             :   OUT   std_logic;
          enb_1_100_0                     :   OUT   std_logic;
          enb_1_100_1                     :   OUT   std_logic
          );
  END COMPONENT;

  COMPONENT Sine_Wave_UV
    PORT( clk                             :   IN    std_logic;
          reset                           :   IN    std_logic;
          enb_1_100_0                     :   IN    std_logic;
          Sine_Wave_out1                  :   OUT   std_logic_vector(13 DOWNTO 0)  -- sfix14
          );
  END COMPONENT;

  COMPONENT Sine_Wave_VW
    PORT( clk                             :   IN    std_logic;
          reset                           :   IN    std_logic;
          enb_1_100_0                     :   IN    std_logic;
          Sine_Wave_out1                  :   OUT   std_logic_vector(13 DOWNTO 0)  -- sfix14
          );
  END COMPONENT;

  COMPONENT SRDBDC2
    PORT( V_uv                            :   IN    std_logic_vector(13 DOWNTO 0);  -- ufix14
          V_vw                            :   IN    std_logic_vector(13 DOWNTO 0);  -- ufix14
          Iout_u                          :   IN    std_logic_vector(13 DOWNTO 0);  -- ufix14
          Iout_v                          :   IN    std_logic_vector(13 DOWNTO 0);  -- ufix14
          IL_u                            :   IN    std_logic_vector(13 DOWNTO 0);  -- ufix14
          IL_v                            :   IN    std_logic_vector(13 DOWNTO 0);  -- ufix14
          Vref_a                          :   IN    std_logic_vector(13 DOWNTO 0);  -- sfix14
          Vref_b                          :   IN    std_logic_vector(13 DOWNTO 0);  -- sfix14
          deltaTa                         :   OUT   std_logic_vector(15 DOWNTO 0);  -- int16
          deltaTb                         :   OUT   std_logic_vector(15 DOWNTO 0);  -- int16
          Vuv                             :   OUT   std_logic_vector(19 DOWNTO 0);  -- sfix20
          Vvw                             :   OUT   std_logic_vector(19 DOWNTO 0);  -- sfix20
          Ioutu                           :   OUT   std_logic_vector(19 DOWNTO 0);  -- sfix20
          IoutV                           :   OUT   std_logic_vector(19 DOWNTO 0);  -- sfix20
          ILu                             :   OUT   std_logic_vector(19 DOWNTO 0);  -- sfix20
          ILv                             :   OUT   std_logic_vector(19 DOWNTO 0);  -- sfix20
          Vrefa                           :   OUT   std_logic_vector(19 DOWNTO 0);  -- sfix20
          Vrefb                           :   OUT   std_logic_vector(19 DOWNTO 0)  -- sfix20
          );
  END COMPONENT;

  COMPONENT nfp_convert_sfix_16_En0_to_double
    PORT( nfp_in                          :   IN    std_logic_vector(15 DOWNTO 0);  -- int16
          nfp_out                         :   OUT   std_logic_vector(63 DOWNTO 0)  -- double
          );
  END COMPONENT;

  COMPONENT nfp_convert_fix_8_En0_to_double
    PORT( nfp_in                          :   IN    std_logic_vector(7 DOWNTO 0);  -- uint8
          nfp_out                         :   OUT   std_logic_vector(63 DOWNTO 0)  -- double
          );
  END COMPONENT;

  COMPONENT hold_function3
    PORT( clk                             :   IN    std_logic;
          reset                           :   IN    std_logic;
          enb_1_100_0                     :   IN    std_logic;
          u1                              :   IN    std_logic_vector(63 DOWNTO 0);  -- double
          cnt                             :   IN    std_logic_vector(63 DOWNTO 0);  -- double
          y1                              :   OUT   std_logic_vector(15 DOWNTO 0)  -- int16
          );
  END COMPONENT;

  COMPONENT SVPWM_calcurate
    PORT( clk                             :   IN    std_logic;
          reset                           :   IN    std_logic;
          enb_1_100_0                     :   IN    std_logic;
          enb_1_100_1                     :   IN    std_logic;
          enb                             :   IN    std_logic;
          ua                              :   IN    std_logic_vector(15 DOWNTO 0);  -- int16
          ub                              :   IN    std_logic_vector(15 DOWNTO 0);  -- int16
          counter                         :   IN    std_logic_vector(7 DOWNTO 0);  -- uint8
          SU                              :   OUT   std_logic_vector(3 DOWNTO 0);  -- ufix4
          SV                              :   OUT   std_logic_vector(3 DOWNTO 0);  -- ufix4
          SW                              :   OUT   std_logic_vector(3 DOWNTO 0)  -- ufix4
          );
  END COMPONENT;

  COMPONENT Cell_balance
    PORT( clk                             :   IN    std_logic;
          reset                           :   IN    std_logic;
          enb_1_100_0                     :   IN    std_logic;
          SU_Level                        :   IN    std_logic_vector(3 DOWNTO 0);  -- ufix4
          SV_Level                        :   IN    std_logic_vector(3 DOWNTO 0);  -- ufix4
          SW_Level                        :   IN    std_logic_vector(3 DOWNTO 0);  -- ufix4
          cnt                             :   IN    std_logic_vector(7 DOWNTO 0);  -- uint8
          Up_current                      :   IN    std_logic_vector(15 DOWNTO 0);  -- int16
          Un_current                      :   IN    std_logic_vector(15 DOWNTO 0);  -- int16
          Vp_current                      :   IN    std_logic_vector(15 DOWNTO 0);  -- int16
          Vn_current                      :   IN    std_logic_vector(15 DOWNTO 0);  -- int16
          Wp_current                      :   IN    std_logic_vector(15 DOWNTO 0);  -- int16
          UW_current                      :   IN    std_logic_vector(15 DOWNTO 0);  -- int16
          CV_U1                           :   IN    std_logic_vector(13 DOWNTO 0);  -- ufix14
          CV_U2                           :   IN    std_logic_vector(13 DOWNTO 0);  -- ufix14
          CV_U3                           :   IN    std_logic_vector(13 DOWNTO 0);  -- ufix14
          CV_U4                           :   IN    std_logic_vector(13 DOWNTO 0);  -- ufix14
          CV_U5                           :   IN    std_logic_vector(13 DOWNTO 0);  -- ufix14
          CV_U6                           :   IN    std_logic_vector(13 DOWNTO 0);  -- ufix14
          CV_U7                           :   IN    std_logic_vector(13 DOWNTO 0);  -- ufix14
          CV_U8                           :   IN    std_logic_vector(13 DOWNTO 0);  -- ufix14
          CV_V1                           :   IN    std_logic_vector(13 DOWNTO 0);  -- ufix14
          CV_V2                           :   IN    std_logic_vector(13 DOWNTO 0);  -- ufix14
          CV_V3                           :   IN    std_logic_vector(13 DOWNTO 0);  -- ufix14
          CV_V4                           :   IN    std_logic_vector(13 DOWNTO 0);  -- ufix14
          CV_V5                           :   IN    std_logic_vector(13 DOWNTO 0);  -- ufix14
          CV_V6                           :   IN    std_logic_vector(13 DOWNTO 0);  -- ufix14
          CV_V7                           :   IN    std_logic_vector(13 DOWNTO 0);  -- ufix14
          CV_V8                           :   IN    std_logic_vector(13 DOWNTO 0);  -- ufix14
          CV_W1                           :   IN    std_logic_vector(13 DOWNTO 0);  -- ufix14
          CV_W2                           :   IN    std_logic_vector(13 DOWNTO 0);  -- ufix14
          CV_W3                           :   IN    std_logic_vector(13 DOWNTO 0);  -- ufix14
          CV_W4                           :   IN    std_logic_vector(13 DOWNTO 0);  -- ufix14
          CV_W5                           :   IN    std_logic_vector(13 DOWNTO 0);  -- ufix14
          CV_W6                           :   IN    std_logic_vector(13 DOWNTO 0);  -- ufix14
          CV_W7                           :   IN    std_logic_vector(13 DOWNTO 0);  -- ufix14
          CV_W8                           :   IN    std_logic_vector(13 DOWNTO 0);  -- ufix14
          MMC_PWM_OUT1                    :   OUT   std_logic_vector(47 DOWNTO 0);  -- ufix48
          MMC_PWM_OUT2                    :   OUT   std_logic_vector(47 DOWNTO 0)  -- ufix48
          );
  END COMPONENT;

  -- Component Configuration Statements
  FOR ALL : Control_top2_tc
    USE ENTITY work.Control_top2_tc(rtl);

  FOR ALL : Sine_Wave_UV
    USE ENTITY work.Sine_Wave_UV(rtl);

  FOR ALL : Sine_Wave_VW
    USE ENTITY work.Sine_Wave_VW(rtl);

  FOR ALL : SRDBDC2
    USE ENTITY work.SRDBDC2(rtl);

  FOR ALL : nfp_convert_sfix_16_En0_to_double
    USE ENTITY work.nfp_convert_sfix_16_En0_to_double(rtl);

  FOR ALL : nfp_convert_fix_8_En0_to_double
    USE ENTITY work.nfp_convert_fix_8_En0_to_double(rtl);

  FOR ALL : hold_function3
    USE ENTITY work.hold_function3(rtl);

  FOR ALL : SVPWM_calcurate
    USE ENTITY work.SVPWM_calcurate(rtl);

  FOR ALL : Cell_balance
    USE ENTITY work.Cell_balance(rtl);

  -- Signals
  SIGNAL enb_1_100_0                      : std_logic;
  SIGNAL enb_1_100_1                      : std_logic;
  SIGNAL enb                              : std_logic;
  SIGNAL switch_control_unsigned          : unsigned(13 DOWNTO 0);  -- ufix14
  SIGNAL Sine_Wave_UV_out1                : std_logic_vector(13 DOWNTO 0);  -- ufix14
  SIGNAL switch_compare_1                 : std_logic;
  SIGNAL Sine_Wave_UV_out1_signed         : signed(13 DOWNTO 0);  -- sfix14
  SIGNAL Data_Type_Conversion_out1        : signed(15 DOWNTO 0);  -- int16
  SIGNAL alpha3929_0_0041_mul_temp        : signed(31 DOWNTO 0);  -- sfix32_En10
  SIGNAL alpha3929_0_0041_out1            : signed(15 DOWNTO 0);  -- int16
  SIGNAL Sine_Wave_VW_out1                : std_logic_vector(13 DOWNTO 0);  -- ufix14
  SIGNAL SRDBDC2_out1                     : std_logic_vector(15 DOWNTO 0);  -- ufix16
  SIGNAL SRDBDC2_out2                     : std_logic_vector(15 DOWNTO 0);  -- ufix16
  SIGNAL SRDBDC2_out3                     : std_logic_vector(19 DOWNTO 0);  -- ufix20
  SIGNAL SRDBDC2_out4                     : std_logic_vector(19 DOWNTO 0);  -- ufix20
  SIGNAL SRDBDC2_out5                     : std_logic_vector(19 DOWNTO 0);  -- ufix20
  SIGNAL SRDBDC2_out6                     : std_logic_vector(19 DOWNTO 0);  -- ufix20
  SIGNAL SRDBDC2_out7                     : std_logic_vector(19 DOWNTO 0);  -- ufix20
  SIGNAL SRDBDC2_out8                     : std_logic_vector(19 DOWNTO 0);  -- ufix20
  SIGNAL SRDBDC2_out9                     : std_logic_vector(19 DOWNTO 0);  -- ufix20
  SIGNAL SRDBDC2_out10                    : std_logic_vector(19 DOWNTO 0);  -- ufix20
  SIGNAL count_step                       : unsigned(7 DOWNTO 0);  -- uint8
  SIGNAL count_from_1                     : unsigned(7 DOWNTO 0);  -- uint8
  SIGNAL Counter_Limited_out1             : unsigned(7 DOWNTO 0);  -- uint8
  SIGNAL count                            : unsigned(7 DOWNTO 0);  -- uint8
  SIGNAL needToWrap                       : std_logic;
  SIGNAL count_value                      : unsigned(7 DOWNTO 0);  -- uint8
  SIGNAL SRDBDC2_out1_signed              : signed(15 DOWNTO 0);  -- int16
  SIGNAL Switch_out1                      : signed(15 DOWNTO 0);  -- int16
  SIGNAL Data_Type_Conversion8_out1       : std_logic_vector(63 DOWNTO 0);  -- ufix64
  SIGNAL Data_Type_Conversion10_out1      : std_logic_vector(63 DOWNTO 0);  -- ufix64
  SIGNAL hold_function4_out1              : std_logic_vector(15 DOWNTO 0);  -- ufix16
  SIGNAL hold_function4_out1_signed       : signed(15 DOWNTO 0);  -- int16
  SIGNAL switch_compare_1_1               : std_logic;
  SIGNAL Sine_Wave_VW_out1_signed         : signed(13 DOWNTO 0);  -- sfix14
  SIGNAL Data_Type_Conversion1_out1       : signed(15 DOWNTO 0);  -- int16
  SIGNAL Gain1_mul_temp                   : signed(31 DOWNTO 0);  -- sfix32_En10
  SIGNAL Gain1_out1                       : signed(15 DOWNTO 0);  -- int16
  SIGNAL SRDBDC2_out2_signed              : signed(15 DOWNTO 0);  -- int16
  SIGNAL Switch1_out1                     : signed(15 DOWNTO 0);  -- int16
  SIGNAL Data_Type_Conversion12_out1      : std_logic_vector(63 DOWNTO 0);  -- ufix64
  SIGNAL Data_Type_Conversion13_out1      : std_logic_vector(63 DOWNTO 0);  -- ufix64
  SIGNAL hold_function3_out1              : std_logic_vector(15 DOWNTO 0);  -- ufix16
  SIGNAL hold_function3_out1_signed       : signed(15 DOWNTO 0);  -- int16
  SIGNAL SVPWM_calcurate_out1             : std_logic_vector(3 DOWNTO 0);  -- ufix4
  SIGNAL SVPWM_calcurate_out2             : std_logic_vector(3 DOWNTO 0);  -- ufix4
  SIGNAL SVPWM_calcurate_out3             : std_logic_vector(3 DOWNTO 0);  -- ufix4
  SIGNAL Up_current_unsigned              : unsigned(13 DOWNTO 0);  -- ufix14
  SIGNAL Saturation31_out1                : unsigned(13 DOWNTO 0);  -- ufix14
  SIGNAL Un_current_unsigned              : unsigned(13 DOWNTO 0);  -- ufix14
  SIGNAL Saturation33_out1                : unsigned(13 DOWNTO 0);  -- ufix14
  SIGNAL Vp_current_unsigned              : unsigned(13 DOWNTO 0);  -- ufix14
  SIGNAL Saturation3_out1                 : unsigned(13 DOWNTO 0);  -- ufix14
  SIGNAL Vn_current_unsigned              : unsigned(13 DOWNTO 0);  -- ufix14
  SIGNAL Saturation5_out1                 : unsigned(13 DOWNTO 0);  -- ufix14
  SIGNAL Wp_current_unsigned              : unsigned(13 DOWNTO 0);  -- ufix14
  SIGNAL Saturation2_out1                 : unsigned(13 DOWNTO 0);  -- ufix14
  SIGNAL UW_current_unsigned              : unsigned(13 DOWNTO 0);  -- ufix14
  SIGNAL Saturation4_out1                 : unsigned(13 DOWNTO 0);  -- ufix14
  SIGNAL Mux_out1                         : vector_of_unsigned14(0 TO 5);  -- ufix14 [6]
  SIGNAL Data_Type_Conversion9_out1       : vector_of_signed16(0 TO 5);  -- int16 [6]
  SIGNAL Constant_out1                    : signed(15 DOWNTO 0);  -- int16
  SIGNAL Sum14_out1                       : vector_of_signed16(0 TO 5);  -- int16 [6]
  SIGNAL CV_U1_unsigned                   : unsigned(13 DOWNTO 0);  -- ufix14
  SIGNAL Saturation7_out1                 : unsigned(13 DOWNTO 0);  -- ufix14
  SIGNAL CV_U2_unsigned                   : unsigned(13 DOWNTO 0);  -- ufix14
  SIGNAL Saturation9_out1                 : unsigned(13 DOWNTO 0);  -- ufix14
  SIGNAL CV_U3_unsigned                   : unsigned(13 DOWNTO 0);  -- ufix14
  SIGNAL Saturation6_out1                 : unsigned(13 DOWNTO 0);  -- ufix14
  SIGNAL CV_U4_unsigned                   : unsigned(13 DOWNTO 0);  -- ufix14
  SIGNAL Saturation8_out1                 : unsigned(13 DOWNTO 0);  -- ufix14
  SIGNAL CV_U5_unsigned                   : unsigned(13 DOWNTO 0);  -- ufix14
  SIGNAL Saturation11_out1                : unsigned(13 DOWNTO 0);  -- ufix14
  SIGNAL CV_U6_unsigned                   : unsigned(13 DOWNTO 0);  -- ufix14
  SIGNAL Saturation13_out1                : unsigned(13 DOWNTO 0);  -- ufix14
  SIGNAL CV_U7_unsigned                   : unsigned(13 DOWNTO 0);  -- ufix14
  SIGNAL Saturation10_out1                : unsigned(13 DOWNTO 0);  -- ufix14
  SIGNAL CV_U8_unsigned                   : unsigned(13 DOWNTO 0);  -- ufix14
  SIGNAL Saturation12_out1                : unsigned(13 DOWNTO 0);  -- ufix14
  SIGNAL CV_V1_unsigned                   : unsigned(13 DOWNTO 0);  -- ufix14
  SIGNAL Saturation27_out1                : unsigned(13 DOWNTO 0);  -- ufix14
  SIGNAL CV_V2_unsigned                   : unsigned(13 DOWNTO 0);  -- ufix14
  SIGNAL Saturation29_out1                : unsigned(13 DOWNTO 0);  -- ufix14
  SIGNAL CV_V3_unsigned                   : unsigned(13 DOWNTO 0);  -- ufix14
  SIGNAL Saturation18_out1                : unsigned(13 DOWNTO 0);  -- ufix14
  SIGNAL CV_V4_unsigned                   : unsigned(13 DOWNTO 0);  -- ufix14
  SIGNAL Saturation28_out1                : unsigned(13 DOWNTO 0);  -- ufix14
  SIGNAL CV_V5_unsigned                   : unsigned(13 DOWNTO 0);  -- ufix14
  SIGNAL Saturation15_out1                : unsigned(13 DOWNTO 0);  -- ufix14
  SIGNAL CV_V6_unsigned                   : unsigned(13 DOWNTO 0);  -- ufix14
  SIGNAL Saturation17_out1                : unsigned(13 DOWNTO 0);  -- ufix14
  SIGNAL CV_V7_unsigned                   : unsigned(13 DOWNTO 0);  -- ufix14
  SIGNAL Saturation14_out1                : unsigned(13 DOWNTO 0);  -- ufix14
  SIGNAL CV_V8_unsigned                   : unsigned(13 DOWNTO 0);  -- ufix14
  SIGNAL Saturation16_out1                : unsigned(13 DOWNTO 0);  -- ufix14
  SIGNAL CV_W1_unsigned                   : unsigned(13 DOWNTO 0);  -- ufix14
  SIGNAL Saturation24_out1                : unsigned(13 DOWNTO 0);  -- ufix14
  SIGNAL CV_W2_unsigned                   : unsigned(13 DOWNTO 0);  -- ufix14
  SIGNAL Saturation26_out1                : unsigned(13 DOWNTO 0);  -- ufix14
  SIGNAL CV_W3_unsigned                   : unsigned(13 DOWNTO 0);  -- ufix14
  SIGNAL Saturation23_out1                : unsigned(13 DOWNTO 0);  -- ufix14
  SIGNAL CV_W4_unsigned                   : unsigned(13 DOWNTO 0);  -- ufix14
  SIGNAL Saturation25_out1                : unsigned(13 DOWNTO 0);  -- ufix14
  SIGNAL CV_W5_unsigned                   : unsigned(13 DOWNTO 0);  -- ufix14
  SIGNAL Saturation20_out1                : unsigned(13 DOWNTO 0);  -- ufix14
  SIGNAL CV_W6_unsigned                   : unsigned(13 DOWNTO 0);  -- ufix14
  SIGNAL Saturation22_out1                : unsigned(13 DOWNTO 0);  -- ufix14
  SIGNAL CV_W7_unsigned                   : unsigned(13 DOWNTO 0);  -- ufix14
  SIGNAL Saturation19_out1                : unsigned(13 DOWNTO 0);  -- ufix14
  SIGNAL CV_W8_unsigned                   : unsigned(13 DOWNTO 0);  -- ufix14
  SIGNAL Saturation21_out1                : unsigned(13 DOWNTO 0);  -- ufix14
  SIGNAL Cell_balance_out1                : std_logic_vector(47 DOWNTO 0);  -- ufix48
  SIGNAL Cell_balance_out2                : std_logic_vector(47 DOWNTO 0);  -- ufix48
  SIGNAL SRDBDC2_out3_signed              : signed(19 DOWNTO 0);  -- sfix20
  SIGNAL delayMatch_reg                   : vector_of_signed20(0 TO 2);  -- sfix20 [3]
  SIGNAL delayMatch_reg_next              : vector_of_signed20(0 TO 2);  -- sfix20 [3]
  SIGNAL SRDBDC2_out3_1                   : signed(19 DOWNTO 0);  -- sfix20
  SIGNAL SRDBDC2_out4_signed              : signed(19 DOWNTO 0);  -- sfix20
  SIGNAL delayMatch1_reg                  : vector_of_signed20(0 TO 2);  -- sfix20 [3]
  SIGNAL delayMatch1_reg_next             : vector_of_signed20(0 TO 2);  -- sfix20 [3]
  SIGNAL SRDBDC2_out4_1                   : signed(19 DOWNTO 0);  -- sfix20
  SIGNAL SRDBDC2_out5_signed              : signed(19 DOWNTO 0);  -- sfix20
  SIGNAL delayMatch2_reg                  : vector_of_signed20(0 TO 2);  -- sfix20 [3]
  SIGNAL delayMatch2_reg_next             : vector_of_signed20(0 TO 2);  -- sfix20 [3]
  SIGNAL SRDBDC2_out5_1                   : signed(19 DOWNTO 0);  -- sfix20
  SIGNAL SRDBDC2_out6_signed              : signed(19 DOWNTO 0);  -- sfix20
  SIGNAL delayMatch3_reg                  : vector_of_signed20(0 TO 2);  -- sfix20 [3]
  SIGNAL delayMatch3_reg_next             : vector_of_signed20(0 TO 2);  -- sfix20 [3]
  SIGNAL SRDBDC2_out6_1                   : signed(19 DOWNTO 0);  -- sfix20
  SIGNAL SRDBDC2_out7_signed              : signed(19 DOWNTO 0);  -- sfix20
  SIGNAL delayMatch4_reg                  : vector_of_signed20(0 TO 2);  -- sfix20 [3]
  SIGNAL delayMatch4_reg_next             : vector_of_signed20(0 TO 2);  -- sfix20 [3]
  SIGNAL SRDBDC2_out7_1                   : signed(19 DOWNTO 0);  -- sfix20
  SIGNAL SRDBDC2_out8_signed              : signed(19 DOWNTO 0);  -- sfix20
  SIGNAL delayMatch5_reg                  : vector_of_signed20(0 TO 2);  -- sfix20 [3]
  SIGNAL delayMatch5_reg_next             : vector_of_signed20(0 TO 2);  -- sfix20 [3]
  SIGNAL SRDBDC2_out8_1                   : signed(19 DOWNTO 0);  -- sfix20
  SIGNAL SRDBDC2_out9_signed              : signed(19 DOWNTO 0);  -- sfix20
  SIGNAL delayMatch6_reg                  : vector_of_signed20(0 TO 2);  -- sfix20 [3]
  SIGNAL delayMatch6_reg_next             : vector_of_signed20(0 TO 2);  -- sfix20 [3]
  SIGNAL SRDBDC2_out9_1                   : signed(19 DOWNTO 0);  -- sfix20
  SIGNAL SRDBDC2_out10_signed             : signed(19 DOWNTO 0);  -- sfix20
  SIGNAL delayMatch7_reg                  : vector_of_signed20(0 TO 2);  -- sfix20 [3]
  SIGNAL delayMatch7_reg_next             : vector_of_signed20(0 TO 2);  -- sfix20 [3]
  SIGNAL SRDBDC2_out10_1                  : signed(19 DOWNTO 0);  -- sfix20
  SIGNAL Sum14_out1_0                     : signed(15 DOWNTO 0);  -- int16
  SIGNAL delayMatch8_reg                  : vector_of_signed16(0 TO 2);  -- sfix16 [3]
  SIGNAL delayMatch8_reg_next             : vector_of_signed16(0 TO 2);  -- sfix16 [3]
  SIGNAL Demux_out1                       : signed(15 DOWNTO 0);  -- int16
  SIGNAL Sum14_out1_1                     : signed(15 DOWNTO 0);  -- int16
  SIGNAL delayMatch9_reg                  : vector_of_signed16(0 TO 2);  -- sfix16 [3]
  SIGNAL delayMatch9_reg_next             : vector_of_signed16(0 TO 2);  -- sfix16 [3]
  SIGNAL Demux_out2                       : signed(15 DOWNTO 0);  -- int16
  SIGNAL delayMatch10_reg                 : vector_of_signed16(0 TO 2);  -- sfix16 [3]
  SIGNAL delayMatch10_reg_next            : vector_of_signed16(0 TO 2);  -- sfix16 [3]
  SIGNAL Data_Type_Conversion11_out1      : signed(15 DOWNTO 0);  -- int16
  SIGNAL delayMatch11_reg                 : vector_of_signed16(0 TO 2);  -- sfix16 [3]
  SIGNAL delayMatch11_reg_next            : vector_of_signed16(0 TO 2);  -- sfix16 [3]
  SIGNAL Data_Type_Conversion14_out1      : signed(15 DOWNTO 0);  -- int16
  SIGNAL delayMatch12_reg                 : vector_of_unsigned14(0 TO 2);  -- ufix14 [3]
  SIGNAL delayMatch12_reg_next            : vector_of_unsigned14(0 TO 2);  -- ufix14 [3]
  SIGNAL switch_control_1                 : unsigned(13 DOWNTO 0);  -- ufix14

BEGIN
  u_Control_top2_tc : Control_top2_tc
    PORT MAP( clk => clk,
              reset => reset,
              clk_enable => clk_enable,
              enb => enb,
              enb_1_100_0 => enb_1_100_0,
              enb_1_100_1 => enb_1_100_1
              );

  u_Sine_Wave_UV : Sine_Wave_UV
    PORT MAP( clk => clk,
              reset => reset,
              enb_1_100_0 => enb_1_100_0,
              Sine_Wave_out1 => Sine_Wave_UV_out1  -- sfix14
              );

  u_Sine_Wave_VW : Sine_Wave_VW
    PORT MAP( clk => clk,
              reset => reset,
              enb_1_100_0 => enb_1_100_0,
              Sine_Wave_out1 => Sine_Wave_VW_out1  -- sfix14
              );

  u_SRDBDC2 : SRDBDC2
    PORT MAP( V_uv => V_uv,  -- ufix14
              V_vw => V_uw,  -- ufix14
              Iout_u => iout_u,  -- ufix14
              Iout_v => iout_v,  -- ufix14
              IL_u => iL_u,  -- ufix14
              IL_v => iL_v,  -- ufix14
              Vref_a => Sine_Wave_UV_out1,  -- sfix14
              Vref_b => Sine_Wave_VW_out1,  -- sfix14
              deltaTa => SRDBDC2_out1,  -- int16
              deltaTb => SRDBDC2_out2,  -- int16
              Vuv => SRDBDC2_out3,  -- sfix20
              Vvw => SRDBDC2_out4,  -- sfix20
              Ioutu => SRDBDC2_out5,  -- sfix20
              IoutV => SRDBDC2_out6,  -- sfix20
              ILu => SRDBDC2_out7,  -- sfix20
              ILv => SRDBDC2_out8,  -- sfix20
              Vrefa => SRDBDC2_out9,  -- sfix20
              Vrefb => SRDBDC2_out10  -- sfix20
              );

  u_MMC_9level_half_cel_Control_Control_top2_nfp_convert_sfix_16_En0_to_double : nfp_convert_sfix_16_En0_to_double
    PORT MAP( nfp_in => std_logic_vector(Switch_out1),  -- int16
              nfp_out => Data_Type_Conversion8_out1  -- double
              );

  u_MMC_9level_half_cel_Control_Control_top2_nfp_convert_fix_8_En0_to_double : nfp_convert_fix_8_En0_to_double
    PORT MAP( nfp_in => std_logic_vector(Counter_Limited_out1),  -- uint8
              nfp_out => Data_Type_Conversion10_out1  -- double
              );

  u_hold_function4 : hold_function3
    PORT MAP( clk => clk,
              reset => reset,
              enb_1_100_0 => enb_1_100_0,
              u1 => Data_Type_Conversion8_out1,  -- double
              cnt => Data_Type_Conversion10_out1,  -- double
              y1 => hold_function4_out1  -- int16
              );

  u_MMC_9level_half_cel_Control_Control_top2_nfp_convert_sfix_16_En0_to_double_1 : nfp_convert_sfix_16_En0_to_double
    PORT MAP( nfp_in => std_logic_vector(Switch1_out1),  -- int16
              nfp_out => Data_Type_Conversion12_out1  -- double
              );

  u_MMC_9level_half_cel_Control_Control_top2_nfp_convert_fix_8_En0_to_double_1 : nfp_convert_fix_8_En0_to_double
    PORT MAP( nfp_in => std_logic_vector(Counter_Limited_out1),  -- uint8
              nfp_out => Data_Type_Conversion13_out1  -- double
              );

  u_hold_function3 : hold_function3
    PORT MAP( clk => clk,
              reset => reset,
              enb_1_100_0 => enb_1_100_0,
              u1 => Data_Type_Conversion12_out1,  -- double
              cnt => Data_Type_Conversion13_out1,  -- double
              y1 => hold_function3_out1  -- int16
              );

  u_SVPWM_calcurate : SVPWM_calcurate
    PORT MAP( clk => clk,
              reset => reset,
              enb_1_100_0 => enb_1_100_0,
              enb_1_100_1 => enb_1_100_1,
              enb => enb,
              ua => hold_function4_out1,  -- int16
              ub => hold_function3_out1,  -- int16
              counter => std_logic_vector(Counter_Limited_out1),  -- uint8
              SU => SVPWM_calcurate_out1,  -- ufix4
              SV => SVPWM_calcurate_out2,  -- ufix4
              SW => SVPWM_calcurate_out3  -- ufix4
              );

  u_Cell_balance : Cell_balance
    PORT MAP( clk => clk,
              reset => reset,
              enb_1_100_0 => enb_1_100_0,
              SU_Level => SVPWM_calcurate_out1,  -- ufix4
              SV_Level => SVPWM_calcurate_out2,  -- ufix4
              SW_Level => SVPWM_calcurate_out3,  -- ufix4
              cnt => std_logic_vector(Counter_Limited_out1),  -- uint8
              Up_current => std_logic_vector(Sum14_out1(0)),  -- int16
              Un_current => std_logic_vector(Sum14_out1(1)),  -- int16
              Vp_current => std_logic_vector(Sum14_out1(2)),  -- int16
              Vn_current => std_logic_vector(Sum14_out1(3)),  -- int16
              Wp_current => std_logic_vector(Sum14_out1(4)),  -- int16
              UW_current => std_logic_vector(Sum14_out1(5)),  -- int16
              CV_U1 => std_logic_vector(Saturation7_out1),  -- ufix14
              CV_U2 => std_logic_vector(Saturation9_out1),  -- ufix14
              CV_U3 => std_logic_vector(Saturation6_out1),  -- ufix14
              CV_U4 => std_logic_vector(Saturation8_out1),  -- ufix14
              CV_U5 => std_logic_vector(Saturation11_out1),  -- ufix14
              CV_U6 => std_logic_vector(Saturation13_out1),  -- ufix14
              CV_U7 => std_logic_vector(Saturation10_out1),  -- ufix14
              CV_U8 => std_logic_vector(Saturation12_out1),  -- ufix14
              CV_V1 => std_logic_vector(Saturation27_out1),  -- ufix14
              CV_V2 => std_logic_vector(Saturation29_out1),  -- ufix14
              CV_V3 => std_logic_vector(Saturation18_out1),  -- ufix14
              CV_V4 => std_logic_vector(Saturation28_out1),  -- ufix14
              CV_V5 => std_logic_vector(Saturation15_out1),  -- ufix14
              CV_V6 => std_logic_vector(Saturation17_out1),  -- ufix14
              CV_V7 => std_logic_vector(Saturation14_out1),  -- ufix14
              CV_V8 => std_logic_vector(Saturation16_out1),  -- ufix14
              CV_W1 => std_logic_vector(Saturation24_out1),  -- ufix14
              CV_W2 => std_logic_vector(Saturation26_out1),  -- ufix14
              CV_W3 => std_logic_vector(Saturation23_out1),  -- ufix14
              CV_W4 => std_logic_vector(Saturation25_out1),  -- ufix14
              CV_W5 => std_logic_vector(Saturation20_out1),  -- ufix14
              CV_W6 => std_logic_vector(Saturation22_out1),  -- ufix14
              CV_W7 => std_logic_vector(Saturation19_out1),  -- ufix14
              CV_W8 => std_logic_vector(Saturation21_out1),  -- ufix14
              MMC_PWM_OUT1 => Cell_balance_out1,  -- ufix48
              MMC_PWM_OUT2 => Cell_balance_out2  -- ufix48
              );

  switch_control_unsigned <= unsigned(switch_control);

  
  switch_compare_1 <= '1' WHEN switch_control_unsigned > to_unsigned(16#36B0#, 14) ELSE
      '0';

  Sine_Wave_UV_out1_signed <= signed(Sine_Wave_UV_out1);

  Data_Type_Conversion_out1 <= resize(Sine_Wave_UV_out1_signed, 16);

  alpha3929_0_0041_mul_temp <= to_signed(16#4400#, 16) * Data_Type_Conversion_out1;
  alpha3929_0_0041_out1 <= alpha3929_0_0041_mul_temp(25 DOWNTO 10);

  count_step <= to_unsigned(16#01#, 8);

  count_from_1 <= to_unsigned(16#00#, 8);

  count <= Counter_Limited_out1 + count_step;

  
  needToWrap <= '1' WHEN Counter_Limited_out1 >= to_unsigned(16#C7#, 8) ELSE
      '0';

  
  count_value <= count WHEN needToWrap = '0' ELSE
      count_from_1;

  -- Count limited, Unsigned Counter
  --  initial value   = 0
  --  step value      = 1
  --  count to value  = 199
  Counter_Limited_process : PROCESS (clk, reset)
  BEGIN
    IF reset = '1' THEN
      Counter_Limited_out1 <= to_unsigned(16#00#, 8);
    ELSIF clk'EVENT AND clk = '1' THEN
      IF enb_1_100_0 = '1' THEN
        Counter_Limited_out1 <= count_value;
      END IF;
    END IF;
  END PROCESS Counter_Limited_process;


  SRDBDC2_out1_signed <= signed(SRDBDC2_out1);

  
  Switch_out1 <= alpha3929_0_0041_out1 WHEN switch_compare_1 = '0' ELSE
      SRDBDC2_out1_signed;

  hold_function4_out1_signed <= signed(hold_function4_out1);

  
  switch_compare_1_1 <= '1' WHEN switch_control_unsigned > to_unsigned(16#36B0#, 14) ELSE
      '0';

  Sine_Wave_VW_out1_signed <= signed(Sine_Wave_VW_out1);

  Data_Type_Conversion1_out1 <= resize(Sine_Wave_VW_out1_signed, 16);

  Gain1_mul_temp <= to_signed(16#4400#, 16) * Data_Type_Conversion1_out1;
  Gain1_out1 <= Gain1_mul_temp(25 DOWNTO 10);

  SRDBDC2_out2_signed <= signed(SRDBDC2_out2);

  
  Switch1_out1 <= Gain1_out1 WHEN switch_compare_1_1 = '0' ELSE
      SRDBDC2_out2_signed;

  hold_function3_out1_signed <= signed(hold_function3_out1);

  Up_current_unsigned <= unsigned(Up_current);

  
  Saturation31_out1 <= to_unsigned(16#3FFE#, 14) WHEN Up_current_unsigned > to_unsigned(16#3FFE#, 14) ELSE
      Up_current_unsigned;

  Un_current_unsigned <= unsigned(Un_current);

  
  Saturation33_out1 <= to_unsigned(16#3FFE#, 14) WHEN Un_current_unsigned > to_unsigned(16#3FFE#, 14) ELSE
      Un_current_unsigned;

  Vp_current_unsigned <= unsigned(Vp_current);

  
  Saturation3_out1 <= to_unsigned(16#3FFE#, 14) WHEN Vp_current_unsigned > to_unsigned(16#3FFE#, 14) ELSE
      Vp_current_unsigned;

  Vn_current_unsigned <= unsigned(Vn_current);

  
  Saturation5_out1 <= to_unsigned(16#3FFE#, 14) WHEN Vn_current_unsigned > to_unsigned(16#3FFE#, 14) ELSE
      Vn_current_unsigned;

  Wp_current_unsigned <= unsigned(Wp_current);

  
  Saturation2_out1 <= to_unsigned(16#3FFE#, 14) WHEN Wp_current_unsigned > to_unsigned(16#3FFE#, 14) ELSE
      Wp_current_unsigned;

  UW_current_unsigned <= unsigned(UW_current);

  
  Saturation4_out1 <= to_unsigned(16#3FFE#, 14) WHEN UW_current_unsigned > to_unsigned(16#3FFE#, 14) ELSE
      UW_current_unsigned;

  Mux_out1(0) <= Saturation31_out1;
  Mux_out1(1) <= Saturation33_out1;
  Mux_out1(2) <= Saturation3_out1;
  Mux_out1(3) <= Saturation5_out1;
  Mux_out1(4) <= Saturation2_out1;
  Mux_out1(5) <= Saturation4_out1;

  Data_Type_Conversion9_out1(0) <= signed(resize(Mux_out1(0), 16));
  Data_Type_Conversion9_out1(1) <= signed(resize(Mux_out1(1), 16));
  Data_Type_Conversion9_out1(2) <= signed(resize(Mux_out1(2), 16));
  Data_Type_Conversion9_out1(3) <= signed(resize(Mux_out1(3), 16));
  Data_Type_Conversion9_out1(4) <= signed(resize(Mux_out1(4), 16));
  Data_Type_Conversion9_out1(5) <= signed(resize(Mux_out1(5), 16));

  Constant_out1 <= to_signed(16#22CE#, 16);

  Sum14_out1(0) <= Data_Type_Conversion9_out1(0) - Constant_out1;
  Sum14_out1(1) <= Data_Type_Conversion9_out1(1) - Constant_out1;
  Sum14_out1(2) <= Data_Type_Conversion9_out1(2) - Constant_out1;
  Sum14_out1(3) <= Data_Type_Conversion9_out1(3) - Constant_out1;
  Sum14_out1(4) <= Data_Type_Conversion9_out1(4) - Constant_out1;
  Sum14_out1(5) <= Data_Type_Conversion9_out1(5) - Constant_out1;

  CV_U1_unsigned <= unsigned(CV_U1);

  
  Saturation7_out1 <= to_unsigned(16#3FFE#, 14) WHEN CV_U1_unsigned > to_unsigned(16#3FFE#, 14) ELSE
      CV_U1_unsigned;

  CV_U2_unsigned <= unsigned(CV_U2);

  
  Saturation9_out1 <= to_unsigned(16#3FFE#, 14) WHEN CV_U2_unsigned > to_unsigned(16#3FFE#, 14) ELSE
      CV_U2_unsigned;

  CV_U3_unsigned <= unsigned(CV_U3);

  
  Saturation6_out1 <= to_unsigned(16#3FFE#, 14) WHEN CV_U3_unsigned > to_unsigned(16#3FFE#, 14) ELSE
      CV_U3_unsigned;

  CV_U4_unsigned <= unsigned(CV_U4);

  
  Saturation8_out1 <= to_unsigned(16#3FFE#, 14) WHEN CV_U4_unsigned > to_unsigned(16#3FFE#, 14) ELSE
      CV_U4_unsigned;

  CV_U5_unsigned <= unsigned(CV_U5);

  
  Saturation11_out1 <= to_unsigned(16#3FFE#, 14) WHEN CV_U5_unsigned > to_unsigned(16#3FFE#, 14) ELSE
      CV_U5_unsigned;

  CV_U6_unsigned <= unsigned(CV_U6);

  
  Saturation13_out1 <= to_unsigned(16#3FFE#, 14) WHEN CV_U6_unsigned > to_unsigned(16#3FFE#, 14) ELSE
      CV_U6_unsigned;

  CV_U7_unsigned <= unsigned(CV_U7);

  
  Saturation10_out1 <= to_unsigned(16#3FFE#, 14) WHEN CV_U7_unsigned > to_unsigned(16#3FFE#, 14) ELSE
      CV_U7_unsigned;

  CV_U8_unsigned <= unsigned(CV_U8);

  
  Saturation12_out1 <= to_unsigned(16#3FFE#, 14) WHEN CV_U8_unsigned > to_unsigned(16#3FFE#, 14) ELSE
      CV_U8_unsigned;

  CV_V1_unsigned <= unsigned(CV_V1);

  
  Saturation27_out1 <= to_unsigned(16#3FFE#, 14) WHEN CV_V1_unsigned > to_unsigned(16#3FFE#, 14) ELSE
      CV_V1_unsigned;

  CV_V2_unsigned <= unsigned(CV_V2);

  
  Saturation29_out1 <= to_unsigned(16#3FFE#, 14) WHEN CV_V2_unsigned > to_unsigned(16#3FFE#, 14) ELSE
      CV_V2_unsigned;

  CV_V3_unsigned <= unsigned(CV_V3);

  
  Saturation18_out1 <= to_unsigned(16#3FFE#, 14) WHEN CV_V3_unsigned > to_unsigned(16#3FFE#, 14) ELSE
      CV_V3_unsigned;

  CV_V4_unsigned <= unsigned(CV_V4);

  
  Saturation28_out1 <= to_unsigned(16#3FFE#, 14) WHEN CV_V4_unsigned > to_unsigned(16#3FFE#, 14) ELSE
      CV_V4_unsigned;

  CV_V5_unsigned <= unsigned(CV_V5);

  
  Saturation15_out1 <= to_unsigned(16#3FFE#, 14) WHEN CV_V5_unsigned > to_unsigned(16#3FFE#, 14) ELSE
      CV_V5_unsigned;

  CV_V6_unsigned <= unsigned(CV_V6);

  
  Saturation17_out1 <= to_unsigned(16#3FFE#, 14) WHEN CV_V6_unsigned > to_unsigned(16#3FFE#, 14) ELSE
      CV_V6_unsigned;

  CV_V7_unsigned <= unsigned(CV_V7);

  
  Saturation14_out1 <= to_unsigned(16#3FFE#, 14) WHEN CV_V7_unsigned > to_unsigned(16#3FFE#, 14) ELSE
      CV_V7_unsigned;

  CV_V8_unsigned <= unsigned(CV_V8);

  
  Saturation16_out1 <= to_unsigned(16#3FFE#, 14) WHEN CV_V8_unsigned > to_unsigned(16#3FFE#, 14) ELSE
      CV_V8_unsigned;

  CV_W1_unsigned <= unsigned(CV_W1);

  
  Saturation24_out1 <= to_unsigned(16#3FFE#, 14) WHEN CV_W1_unsigned > to_unsigned(16#3FFE#, 14) ELSE
      CV_W1_unsigned;

  CV_W2_unsigned <= unsigned(CV_W2);

  
  Saturation26_out1 <= to_unsigned(16#3FFE#, 14) WHEN CV_W2_unsigned > to_unsigned(16#3FFE#, 14) ELSE
      CV_W2_unsigned;

  CV_W3_unsigned <= unsigned(CV_W3);

  
  Saturation23_out1 <= to_unsigned(16#3FFE#, 14) WHEN CV_W3_unsigned > to_unsigned(16#3FFE#, 14) ELSE
      CV_W3_unsigned;

  CV_W4_unsigned <= unsigned(CV_W4);

  
  Saturation25_out1 <= to_unsigned(16#3FFE#, 14) WHEN CV_W4_unsigned > to_unsigned(16#3FFE#, 14) ELSE
      CV_W4_unsigned;

  CV_W5_unsigned <= unsigned(CV_W5);

  
  Saturation20_out1 <= to_unsigned(16#3FFE#, 14) WHEN CV_W5_unsigned > to_unsigned(16#3FFE#, 14) ELSE
      CV_W5_unsigned;

  CV_W6_unsigned <= unsigned(CV_W6);

  
  Saturation22_out1 <= to_unsigned(16#3FFE#, 14) WHEN CV_W6_unsigned > to_unsigned(16#3FFE#, 14) ELSE
      CV_W6_unsigned;

  CV_W7_unsigned <= unsigned(CV_W7);

  
  Saturation19_out1 <= to_unsigned(16#3FFE#, 14) WHEN CV_W7_unsigned > to_unsigned(16#3FFE#, 14) ELSE
      CV_W7_unsigned;

  CV_W8_unsigned <= unsigned(CV_W8);

  
  Saturation21_out1 <= to_unsigned(16#3FFE#, 14) WHEN CV_W8_unsigned > to_unsigned(16#3FFE#, 14) ELSE
      CV_W8_unsigned;

  SRDBDC2_out3_signed <= signed(SRDBDC2_out3);

  delayMatch_process : PROCESS (clk, reset)
  BEGIN
    IF reset = '1' THEN
      delayMatch_reg(0) <= to_signed(16#00000#, 20);
      delayMatch_reg(1) <= to_signed(16#00000#, 20);
      delayMatch_reg(2) <= to_signed(16#00000#, 20);
    ELSIF clk'EVENT AND clk = '1' THEN
      IF enb_1_100_0 = '1' THEN
        delayMatch_reg(0) <= delayMatch_reg_next(0);
        delayMatch_reg(1) <= delayMatch_reg_next(1);
        delayMatch_reg(2) <= delayMatch_reg_next(2);
      END IF;
    END IF;
  END PROCESS delayMatch_process;

  SRDBDC2_out3_1 <= delayMatch_reg(2);
  delayMatch_reg_next(0) <= SRDBDC2_out3_signed;
  delayMatch_reg_next(1) <= delayMatch_reg(0);
  delayMatch_reg_next(2) <= delayMatch_reg(1);

  Vuv <= std_logic_vector(SRDBDC2_out3_1);

  SRDBDC2_out4_signed <= signed(SRDBDC2_out4);

  delayMatch1_process : PROCESS (clk, reset)
  BEGIN
    IF reset = '1' THEN
      delayMatch1_reg(0) <= to_signed(16#00000#, 20);
      delayMatch1_reg(1) <= to_signed(16#00000#, 20);
      delayMatch1_reg(2) <= to_signed(16#00000#, 20);
    ELSIF clk'EVENT AND clk = '1' THEN
      IF enb_1_100_0 = '1' THEN
        delayMatch1_reg(0) <= delayMatch1_reg_next(0);
        delayMatch1_reg(1) <= delayMatch1_reg_next(1);
        delayMatch1_reg(2) <= delayMatch1_reg_next(2);
      END IF;
    END IF;
  END PROCESS delayMatch1_process;

  SRDBDC2_out4_1 <= delayMatch1_reg(2);
  delayMatch1_reg_next(0) <= SRDBDC2_out4_signed;
  delayMatch1_reg_next(1) <= delayMatch1_reg(0);
  delayMatch1_reg_next(2) <= delayMatch1_reg(1);

  Vvw <= std_logic_vector(SRDBDC2_out4_1);

  SRDBDC2_out5_signed <= signed(SRDBDC2_out5);

  delayMatch2_process : PROCESS (clk, reset)
  BEGIN
    IF reset = '1' THEN
      delayMatch2_reg(0) <= to_signed(16#00000#, 20);
      delayMatch2_reg(1) <= to_signed(16#00000#, 20);
      delayMatch2_reg(2) <= to_signed(16#00000#, 20);
    ELSIF clk'EVENT AND clk = '1' THEN
      IF enb_1_100_0 = '1' THEN
        delayMatch2_reg(0) <= delayMatch2_reg_next(0);
        delayMatch2_reg(1) <= delayMatch2_reg_next(1);
        delayMatch2_reg(2) <= delayMatch2_reg_next(2);
      END IF;
    END IF;
  END PROCESS delayMatch2_process;

  SRDBDC2_out5_1 <= delayMatch2_reg(2);
  delayMatch2_reg_next(0) <= SRDBDC2_out5_signed;
  delayMatch2_reg_next(1) <= delayMatch2_reg(0);
  delayMatch2_reg_next(2) <= delayMatch2_reg(1);

  Ioutu <= std_logic_vector(SRDBDC2_out5_1);

  SRDBDC2_out6_signed <= signed(SRDBDC2_out6);

  delayMatch3_process : PROCESS (clk, reset)
  BEGIN
    IF reset = '1' THEN
      delayMatch3_reg(0) <= to_signed(16#00000#, 20);
      delayMatch3_reg(1) <= to_signed(16#00000#, 20);
      delayMatch3_reg(2) <= to_signed(16#00000#, 20);
    ELSIF clk'EVENT AND clk = '1' THEN
      IF enb_1_100_0 = '1' THEN
        delayMatch3_reg(0) <= delayMatch3_reg_next(0);
        delayMatch3_reg(1) <= delayMatch3_reg_next(1);
        delayMatch3_reg(2) <= delayMatch3_reg_next(2);
      END IF;
    END IF;
  END PROCESS delayMatch3_process;

  SRDBDC2_out6_1 <= delayMatch3_reg(2);
  delayMatch3_reg_next(0) <= SRDBDC2_out6_signed;
  delayMatch3_reg_next(1) <= delayMatch3_reg(0);
  delayMatch3_reg_next(2) <= delayMatch3_reg(1);

  Ioutv <= std_logic_vector(SRDBDC2_out6_1);

  SRDBDC2_out7_signed <= signed(SRDBDC2_out7);

  delayMatch4_process : PROCESS (clk, reset)
  BEGIN
    IF reset = '1' THEN
      delayMatch4_reg(0) <= to_signed(16#00000#, 20);
      delayMatch4_reg(1) <= to_signed(16#00000#, 20);
      delayMatch4_reg(2) <= to_signed(16#00000#, 20);
    ELSIF clk'EVENT AND clk = '1' THEN
      IF enb_1_100_0 = '1' THEN
        delayMatch4_reg(0) <= delayMatch4_reg_next(0);
        delayMatch4_reg(1) <= delayMatch4_reg_next(1);
        delayMatch4_reg(2) <= delayMatch4_reg_next(2);
      END IF;
    END IF;
  END PROCESS delayMatch4_process;

  SRDBDC2_out7_1 <= delayMatch4_reg(2);
  delayMatch4_reg_next(0) <= SRDBDC2_out7_signed;
  delayMatch4_reg_next(1) <= delayMatch4_reg(0);
  delayMatch4_reg_next(2) <= delayMatch4_reg(1);

  ILu <= std_logic_vector(SRDBDC2_out7_1);

  SRDBDC2_out8_signed <= signed(SRDBDC2_out8);

  delayMatch5_process : PROCESS (clk, reset)
  BEGIN
    IF reset = '1' THEN
      delayMatch5_reg(0) <= to_signed(16#00000#, 20);
      delayMatch5_reg(1) <= to_signed(16#00000#, 20);
      delayMatch5_reg(2) <= to_signed(16#00000#, 20);
    ELSIF clk'EVENT AND clk = '1' THEN
      IF enb_1_100_0 = '1' THEN
        delayMatch5_reg(0) <= delayMatch5_reg_next(0);
        delayMatch5_reg(1) <= delayMatch5_reg_next(1);
        delayMatch5_reg(2) <= delayMatch5_reg_next(2);
      END IF;
    END IF;
  END PROCESS delayMatch5_process;

  SRDBDC2_out8_1 <= delayMatch5_reg(2);
  delayMatch5_reg_next(0) <= SRDBDC2_out8_signed;
  delayMatch5_reg_next(1) <= delayMatch5_reg(0);
  delayMatch5_reg_next(2) <= delayMatch5_reg(1);

  ILv <= std_logic_vector(SRDBDC2_out8_1);

  SRDBDC2_out9_signed <= signed(SRDBDC2_out9);

  delayMatch6_process : PROCESS (clk, reset)
  BEGIN
    IF reset = '1' THEN
      delayMatch6_reg(0) <= to_signed(16#00000#, 20);
      delayMatch6_reg(1) <= to_signed(16#00000#, 20);
      delayMatch6_reg(2) <= to_signed(16#00000#, 20);
    ELSIF clk'EVENT AND clk = '1' THEN
      IF enb_1_100_0 = '1' THEN
        delayMatch6_reg(0) <= delayMatch6_reg_next(0);
        delayMatch6_reg(1) <= delayMatch6_reg_next(1);
        delayMatch6_reg(2) <= delayMatch6_reg_next(2);
      END IF;
    END IF;
  END PROCESS delayMatch6_process;

  SRDBDC2_out9_1 <= delayMatch6_reg(2);
  delayMatch6_reg_next(0) <= SRDBDC2_out9_signed;
  delayMatch6_reg_next(1) <= delayMatch6_reg(0);
  delayMatch6_reg_next(2) <= delayMatch6_reg(1);

  Vrefa <= std_logic_vector(SRDBDC2_out9_1);

  SRDBDC2_out10_signed <= signed(SRDBDC2_out10);

  delayMatch7_process : PROCESS (clk, reset)
  BEGIN
    IF reset = '1' THEN
      delayMatch7_reg(0) <= to_signed(16#00000#, 20);
      delayMatch7_reg(1) <= to_signed(16#00000#, 20);
      delayMatch7_reg(2) <= to_signed(16#00000#, 20);
    ELSIF clk'EVENT AND clk = '1' THEN
      IF enb_1_100_0 = '1' THEN
        delayMatch7_reg(0) <= delayMatch7_reg_next(0);
        delayMatch7_reg(1) <= delayMatch7_reg_next(1);
        delayMatch7_reg(2) <= delayMatch7_reg_next(2);
      END IF;
    END IF;
  END PROCESS delayMatch7_process;

  SRDBDC2_out10_1 <= delayMatch7_reg(2);
  delayMatch7_reg_next(0) <= SRDBDC2_out10_signed;
  delayMatch7_reg_next(1) <= delayMatch7_reg(0);
  delayMatch7_reg_next(2) <= delayMatch7_reg(1);

  Vrefb <= std_logic_vector(SRDBDC2_out10_1);

  Sum14_out1_0 <= Sum14_out1(0);

  delayMatch8_process : PROCESS (clk, reset)
  BEGIN
    IF reset = '1' THEN
      delayMatch8_reg(0) <= to_signed(16#0000#, 16);
      delayMatch8_reg(1) <= to_signed(16#0000#, 16);
      delayMatch8_reg(2) <= to_signed(16#0000#, 16);
    ELSIF clk'EVENT AND clk = '1' THEN
      IF enb_1_100_0 = '1' THEN
        delayMatch8_reg(0) <= delayMatch8_reg_next(0);
        delayMatch8_reg(1) <= delayMatch8_reg_next(1);
        delayMatch8_reg(2) <= delayMatch8_reg_next(2);
      END IF;
    END IF;
  END PROCESS delayMatch8_process;

  Demux_out1 <= delayMatch8_reg(2);
  delayMatch8_reg_next(0) <= Sum14_out1_0;
  delayMatch8_reg_next(1) <= delayMatch8_reg(0);
  delayMatch8_reg_next(2) <= delayMatch8_reg(1);

  deltaTa <= std_logic_vector(Demux_out1);

  Sum14_out1_1 <= Sum14_out1(1);

  delayMatch9_process : PROCESS (clk, reset)
  BEGIN
    IF reset = '1' THEN
      delayMatch9_reg(0) <= to_signed(16#0000#, 16);
      delayMatch9_reg(1) <= to_signed(16#0000#, 16);
      delayMatch9_reg(2) <= to_signed(16#0000#, 16);
    ELSIF clk'EVENT AND clk = '1' THEN
      IF enb_1_100_0 = '1' THEN
        delayMatch9_reg(0) <= delayMatch9_reg_next(0);
        delayMatch9_reg(1) <= delayMatch9_reg_next(1);
        delayMatch9_reg(2) <= delayMatch9_reg_next(2);
      END IF;
    END IF;
  END PROCESS delayMatch9_process;

  Demux_out2 <= delayMatch9_reg(2);
  delayMatch9_reg_next(0) <= Sum14_out1_1;
  delayMatch9_reg_next(1) <= delayMatch9_reg(0);
  delayMatch9_reg_next(2) <= delayMatch9_reg(1);

  deltaTb <= std_logic_vector(Demux_out2);

  delayMatch10_process : PROCESS (clk, reset)
  BEGIN
    IF reset = '1' THEN
      delayMatch10_reg(0) <= to_signed(16#0000#, 16);
      delayMatch10_reg(1) <= to_signed(16#0000#, 16);
      delayMatch10_reg(2) <= to_signed(16#0000#, 16);
    ELSIF clk'EVENT AND clk = '1' THEN
      IF enb_1_100_0 = '1' THEN
        delayMatch10_reg(0) <= delayMatch10_reg_next(0);
        delayMatch10_reg(1) <= delayMatch10_reg_next(1);
        delayMatch10_reg(2) <= delayMatch10_reg_next(2);
      END IF;
    END IF;
  END PROCESS delayMatch10_process;

  Data_Type_Conversion11_out1 <= delayMatch10_reg(2);
  delayMatch10_reg_next(0) <= hold_function4_out1_signed;
  delayMatch10_reg_next(1) <= delayMatch10_reg(0);
  delayMatch10_reg_next(2) <= delayMatch10_reg(1);

  deltaTa_hold <= std_logic_vector(Data_Type_Conversion11_out1);

  delayMatch11_process : PROCESS (clk, reset)
  BEGIN
    IF reset = '1' THEN
      delayMatch11_reg(0) <= to_signed(16#0000#, 16);
      delayMatch11_reg(1) <= to_signed(16#0000#, 16);
      delayMatch11_reg(2) <= to_signed(16#0000#, 16);
    ELSIF clk'EVENT AND clk = '1' THEN
      IF enb_1_100_0 = '1' THEN
        delayMatch11_reg(0) <= delayMatch11_reg_next(0);
        delayMatch11_reg(1) <= delayMatch11_reg_next(1);
        delayMatch11_reg(2) <= delayMatch11_reg_next(2);
      END IF;
    END IF;
  END PROCESS delayMatch11_process;

  Data_Type_Conversion14_out1 <= delayMatch11_reg(2);
  delayMatch11_reg_next(0) <= hold_function3_out1_signed;
  delayMatch11_reg_next(1) <= delayMatch11_reg(0);
  delayMatch11_reg_next(2) <= delayMatch11_reg(1);

  deltaTb_hold <= std_logic_vector(Data_Type_Conversion14_out1);

  delayMatch12_process : PROCESS (clk, reset)
  BEGIN
    IF reset = '1' THEN
      delayMatch12_reg(0) <= to_unsigned(16#0000#, 14);
      delayMatch12_reg(1) <= to_unsigned(16#0000#, 14);
      delayMatch12_reg(2) <= to_unsigned(16#0000#, 14);
    ELSIF clk'EVENT AND clk = '1' THEN
      IF enb_1_100_0 = '1' THEN
        delayMatch12_reg(0) <= delayMatch12_reg_next(0);
        delayMatch12_reg(1) <= delayMatch12_reg_next(1);
        delayMatch12_reg(2) <= delayMatch12_reg_next(2);
      END IF;
    END IF;
  END PROCESS delayMatch12_process;

  switch_control_1 <= delayMatch12_reg(2);
  delayMatch12_reg_next(0) <= switch_control_unsigned;
  delayMatch12_reg_next(1) <= delayMatch12_reg(0);
  delayMatch12_reg_next(2) <= delayMatch12_reg(1);

  switch_control1 <= std_logic_vector(switch_control_1);

  ce_out <= enb_1_100_1;

  MMC_PWM_OUT1 <= Cell_balance_out1;

  MMC_PWM_OUT2 <= Cell_balance_out2;

END rtl;



\end{mylisting}
%\input{appendix/program/appF_SRDCDB.tex}
%\input{appendix/program/appF_SVPWM_VHDL.tex}
%\input{appendix/program/appF_Cell_balance_module_VHDL.tex}



\begin{comment}


\end{comment}

